% Lesson 3 — Vocabulary: Words and expressions related to interacting with mates, neighbours, adults and school authorities on community celebrations/festivals — Anglais 1ère A — Cours signé OnBuch+
% Programme officiel MINESEC. Compiler avec : tectonic ENA03-vocabulary-words-and-expressions-related-to-interacting.tex
\newcommand{\DOCMATIERE}{Anglais}
\newcommand{\DOCNIVEAU}{1ère A}
\newcommand{\DOCMODULE}{Chapter 1 : Family and Social Life}
\newcommand{\DOCLECON}{ENA03}
\newcommand{\DOCTITRE}{Lesson 3 — Vocabulary: Words and expressions related to interacting with mates, neighbours, adults and school authorities on community celebrations/festivals}
% OnBuch+ — préambule « cours signés » ENRICHI (compilé avec tectonic / XeLaTeX)
% Système visuel « Pop » d'OnBuch+ : fond crème (jamais blanc), bordures encre
% épaisses, ombres « dures » décalées, titres Archivo Black en capitales.
% Version MATHÉMATIQUES : graphes (pgfplots/tikz), boîtes pédagogiques
% (définition, propriété, méthode, attention, le savais-tu, exercice résolu,
% exercice, TP, bilan), page de garde et signature OnBuch+ ancrée en bas.
%
% Macros attendues dans main.tex AVANT \input{preamble} :
%   \DOCMATIERE  \DOCNIVEAU  \DOCMODULE  \DOCLECON (numéro)  \DOCTITRE
\documentclass[11pt]{article}

\usepackage{fontspec}
\usepackage[english,french]{babel} % français = langue principale ; l'anglais via \eng{..} / textebox
\usepackage[a4paper,margin=2cm,top=3.4cm,bottom=2.8cm,headheight=1.4cm,footskip=1.3cm]{geometry}
\usepackage{amsmath,amssymb,amsfonts}
\usepackage{mathtools} % \xmapsto, \coloneqq... souvent utilisés par les modèles
\usepackage{cancel} % simplifications barrées (bilans de forces)
% \abs{z} (module, valeur absolue) et \norm{u} (norme), utilisés par les modèles
\providecommand{\abs}{}\let\abs\relax\DeclarePairedDelimiter{\abs}{\lvert}{\rvert}
\providecommand{\norm}{}\let\norm\relax\DeclarePairedDelimiter{\norm}{\lVert}{\rVert}
\usepackage[table]{xcolor} % [table] : fournit aussi \rowcolors (alternance de lignes)
\usepackage{pagecolor}
\usepackage{tikz}
\usetikzlibrary{arrows.meta,positioning,calc,shapes.geometric,shapes.misc,shapes.arrows,decorations.pathmorphing,decorations.pathreplacing,decorations.markings,patterns,shadows,mindmap,trees,fit,backgrounds,matrix,intersections,angles,quotes}
\usepackage{pgfplots}
\usepackage[version=4]{mhchem} % équations nucléaires \ce{^{235}_{92}U}
\usepackage[european,siunitx]{circuitikz} % schémas électriques
\pgfplotsset{compat=1.18}
\usepgfplotslibrary{fillbetween,groupplots} % aires entre courbes (calcul intégral)
\usepackage{enumitem}
\usepackage{fancyhdr}
% [most] charge toutes les bibliothèques tcolorbox (listings, minted, raster,
% documentation...) et gonfle inutilement la mémoire TeX sur un document long
% (~300 boîtes) : on ne charge que ce qui est réellement utilisé.
\usepackage[skins,breakable]{tcolorbox}
\usepackage{graphicx}
\usepackage{colortbl}
\usepackage{booktabs}
\usepackage{tabularx}
\usepackage{diagbox} % case d'en-tête diagonale des tableaux à double entrée
% Colonnes extensibles centrée / gauche / droite (C, L, R), souvent utilisées par les modèles
\newcolumntype{C}{>{\centering\arraybackslash}X}
\newcolumntype{L}{>{\raggedright\arraybackslash}X}
\newcolumntype{R}{>{\raggedleft\arraybackslash}X}
\usepackage{array}
\usepackage{multicol}
\usepackage{multirow}
\usepackage{siunitx}
% parse-numbers=false : accepte aussi \SI{2,5 \times 10^{-3}}{...}
\sisetup{locale=FR,output-decimal-marker={,},group-minimum-digits=4,parse-numbers=false}
\DeclareSIUnit\ohm{\ensuremath{\symup{\Omega}}} % U+2126 absent de la police
\DeclareSIUnit\division{div} % réglages d'oscilloscope (ms/div, V/div)
\DeclareSIUnit\tr{tr} % tours (vitesse de rotation, ex. \SI{1500}{\tr\per\minute})
\DeclareSIUnit\molar{M} % concentration molaire (ex. \SI{3}{\molar}, \SI{1}{\micro\molar})
\DeclareSIUnit\an{an} % année (le modèle confond parfois avec \year, un compteur TeX) — ex. \SI{5}{\kilo\meter\per\an}
\DeclareSIUnit\FCFA{FCFA} % franc CFA (ex. \SI{500}{\FCFA\per\kilo\gram})
\DeclareSIUnit\degreeBrix{\textdegree Brix} % teneur en sucre d'un sirop/jus (réfractométrie)
\DeclareSIUnit\month{mois} % le modèle utilise parfois \month, un compteur TeX, comme unité de durée
\usepackage{qrcode}
% \degree : commande fréquemment utilisée par les modèles pour un angle ou une
% température en dehors de \SI{}{} (non fournie par les paquets chargés).
\providecommand{\degree}{^{\circ}}
% \qed (fin de démonstration), utilisé par les modèles sans amsthm
\providecommand{\qed}{\hfill$\square$}
% \barre : double barre des valeurs interdites dans un tableau de variations (tkz-tab)
\providecommand{\barre}{\ensuremath{\|}}
% environnement proof (amsthm non chargé) utilisé par les modèles
\ifdefined\proof\else\newenvironment{proof}[1][Démonstration]{\par\noindent\textit{#1.}\ }{\qed\par}\fi
\usepackage{lastpage}
\usepackage{hyperref}
\hypersetup{hidelinks}

% ============================================================
% Palette « Pop » OnBuch+ — mêmes hex que la classe P (plus_ui.dart)
% ============================================================
\definecolor{poporange}{HTML}{F59321}
\definecolor{popdark}{HTML}{E07A0C}
\definecolor{popcream}{HTML}{FFF6E8}
\definecolor{popink}{HTML}{1C1714}
\definecolor{poppurple}{HTML}{7A5AE0}
\definecolor{popgreen}{HTML}{2E9E6A}
\definecolor{poppink}{HTML}{E0457B}
\definecolor{popblue}{HTML}{2F7DE1}
\definecolor{popgold}{HTML}{C98A12}
\definecolor{popmuted}{HTML}{8A7F76}
\definecolor{popline}{HTML}{EADFD2}
\definecolor{popgray}{HTML}{8A7F76}
\colorlet{popgrey}{popgray}
% Alias tolérants (le modèle nomme parfois les couleurs différemment)
\colorlet{popwhite}{white}
\colorlet{popblack}{popink}
\colorlet{popred}{poppink}
\colorlet{popyellow}{popgold}
% Teintes claires (fonds de tableaux/encadrés) — le modèle nomme parfois
% les couleurs avec un suffixe L (ex. popblueL) sans les avoir définies.
\colorlet{poporangeL}{poporange!20}
\colorlet{popdarkL}{popdark!20}
\colorlet{poppurpleL}{poppurple!20}
\colorlet{popgreenL}{popgreen!20}
\colorlet{poppinkL}{poppink!20}
\colorlet{popblueL}{popblue!20}
\colorlet{popgoldL}{popgold!20}
\colorlet{popredL}{popred!20}
\colorlet{popgrayL}{popgray!20}
% Teintes claires pour fonds de tableaux / zones
\colorlet{popblueL}{popblue!10!white}
\colorlet{poporangeL}{poporange!14!white}
\colorlet{popgreenL}{popgreen!12!white}
\colorlet{poppurpleL}{poppurple!12!white}
\colorlet{poppinkL}{poppink!12!white}
\colorlet{popgoldL}{popgold!14!white}

\pagecolor{popcream}

% ============================================================
% Typographie
% ============================================================
\defaultfontfeatures{Ligatures=TeX}
\setmainfont{PlusJakartaSans-Medium.ttf}[Path=../../../fonts/,BoldFont=PlusJakartaSans-Bold.ttf]
% Symboles, API et emojis absents de la police principale : repli sur DejaVu Sans (caractères actifs)
\newfontface\symfont{DejaVuSans.ttf}[Path=../../../fonts/]
\makeatletter
\newcommand\symactive[1]{\catcode#1=13 \begingroup\lccode`\~=#1 \lowercase{\endgroup\def~}{{\symfont\char#1}}}
\newcommand\symblank[1]{\catcode#1=13 \begingroup\lccode`\~=#1 \lowercase{\endgroup\def~}{}}
\newcommand\symhyph[1]{\catcode#1=13 \begingroup\lccode`\~=#1 \lowercase{\endgroup\def~}{\mbox{-}}}
\newcount\symc
\newcommand\symrange[2]{\symc=#1 \@whilenum\symc<#2 \do{\expandafter\symactive\expandafter{\the\symc}\advance\symc by 1 }}
\newcommand\symblankrange[2]{\symc=#1 \@whilenum\symc<#2 \do{\expandafter\symblank\expandafter{\the\symc}\advance\symc by 1 }}
\makeatother
\symrange{"0250}{"02FF}\symactive{"2713}\symactive{"2714}\symactive{"2717}\symactive{"2718}\symactive{"2610}\symactive{"2611}\symactive{"2612}
\symhyph{"2011}\symhyph{"2E1A}\symblankrange{"1F300}{"1FAFF}\symblank{"FE0F}
% Maths en sans-serif (Fira Math) avec chiffres et lettres droites de Plus
% Jakarta, pour que les formules se fondent dans le texte.
\usepackage[math-style=ISO,bold-style=ISO]{unicode-math}
\setmathfont{FiraMath-Regular.otf}[Path=../../../fonts/]
\setmathfont{PlusJakartaSans-Medium.ttf}[Path=../../../fonts/,range={up/num,up/{Latin,latin},\mathup}]
\usepackage{icomma} % virgule décimale sans espace en mode math
% Caractères mathématiques Unicode absents des polices de texte (Plus Jakarta,
% Archivo Black) : rendus par leur équivalent mathématique, en texte comme en
% formule — sinon ils s'affichent en carré vide (ex. « Arithmétique dans ℤ »).
\usepackage{newunicodechar}
\newunicodechar{ℝ}{\ensuremath{\mathbb{R}}}
\newunicodechar{ℤ}{\ensuremath{\mathbb{Z}}}
\newunicodechar{ℂ}{\ensuremath{\mathbb{C}}}
\newunicodechar{ℕ}{\ensuremath{\mathbb{N}}}
\newunicodechar{ℚ}{\ensuremath{\mathbb{Q}}}
\newunicodechar{𝔻}{\ensuremath{\mathbb{D}}}
\newunicodechar{α}{\ensuremath{\alpha}}
\newunicodechar{β}{\ensuremath{\beta}}
\newunicodechar{γ}{\ensuremath{\gamma}}
\newunicodechar{δ}{\ensuremath{\delta}}
\newunicodechar{ε}{\ensuremath{\varepsilon}}
\newunicodechar{θ}{\ensuremath{\theta}}
\newunicodechar{λ}{\ensuremath{\lambda}}
\newunicodechar{μ}{\ensuremath{\mu}}
\newunicodechar{σ}{\ensuremath{\sigma}}
\newunicodechar{τ}{\ensuremath{\tau}}
\newunicodechar{φ}{\ensuremath{\varphi}}
\newunicodechar{ψ}{\ensuremath{\psi}}
\newunicodechar{ω}{\ensuremath{\omega}}
\newunicodechar{Σ}{\ensuremath{\Sigma}}
% majuscules produites par \MakeUppercase dans les titres (α-aminés -> Α)
\newunicodechar{Α}{\ensuremath{\alpha}}
\newunicodechar{Β}{\ensuremath{\beta}}
\newunicodechar{Γ}{\ensuremath{\Gamma}}
\newunicodechar{Δ}{\ensuremath{\Delta}}
\newunicodechar{Ε}{\ensuremath{\varepsilon}}
\newunicodechar{Θ}{\ensuremath{\Theta}}
\newunicodechar{Λ}{\ensuremath{\Lambda}}
\newunicodechar{Μ}{\ensuremath{\mu}}
\newunicodechar{Τ}{\ensuremath{\tau}}
\newunicodechar{Φ}{\ensuremath{\Phi}}
\newunicodechar{Ψ}{\ensuremath{\Psi}}
\newunicodechar{Ω}{\ensuremath{\Omega}}
\newunicodechar{↦}{\ensuremath{\mapsto}}
\newunicodechar{∈}{\ensuremath{\in}}
\newunicodechar{∉}{\ensuremath{\notin}}
\newunicodechar{∧}{\ensuremath{\wedge}}
\newunicodechar{⇔}{\ensuremath{\Leftrightarrow}}
\newunicodechar{⟺}{\ensuremath{\Longleftrightarrow}}
\newunicodechar{⇒}{\ensuremath{\Rightarrow}}
\newunicodechar{⟹}{\ensuremath{\Longrightarrow}}
\newunicodechar{∘}{\ensuremath{\circ}}
\newunicodechar{∪}{\ensuremath{\cup}}
\newunicodechar{∩}{\ensuremath{\cap}}
\newunicodechar{≡}{\ensuremath{\equiv}}
\newunicodechar{⊂}{\ensuremath{\subset}}
\newunicodechar{⊥}{\ensuremath{\perp}}
\newunicodechar{∃}{\ensuremath{\exists}}
\newunicodechar{∀}{\ensuremath{\forall}}
\newunicodechar{∛}{\ensuremath{\sqrt[3]{\,}}}
\newunicodechar{∜}{\ensuremath{\sqrt[4]{\,}}}
\newunicodechar{⃗}{\ensuremath{{}^{\to}}}
\newunicodechar{ⁿ}{\ifmmode^{n}\else\textsuperscript{\ensuremath{n}}\fi}
\newunicodechar{ˣ}{\ifmmode^{x}\else\textsuperscript{\ensuremath{x}}\fi}
\newunicodechar{ᵇ}{\ifmmode^{b}\else\textsuperscript{\ensuremath{b}}\fi}
\newunicodechar{ʳ}{\ifmmode^{r}\else\textsuperscript{\ensuremath{r}}\fi}
\newunicodechar{ᵘ}{\ifmmode^{u}\else\textsuperscript{\ensuremath{u}}\fi}
\newunicodechar{ʸ}{\ifmmode^{y}\else\textsuperscript{\ensuremath{y}}\fi}
\newunicodechar{ᵃ}{\ifmmode^{a}\else\textsuperscript{\ensuremath{a}}\fi}
\newunicodechar{ᵉ}{\ifmmode^{e}\else\textsuperscript{\ensuremath{e}}\fi}
\newunicodechar{ᵏ}{\ifmmode^{k}\else\textsuperscript{\ensuremath{k}}\fi}
\newunicodechar{ᵖ}{\ifmmode^{p}\else\textsuperscript{\ensuremath{p}}\fi}
\newunicodechar{ᴺ}{\ifmmode^{N}\else\textsuperscript{\ensuremath{N}}\fi}
\newunicodechar{ᶜ}{\ifmmode^{c}\else\textsuperscript{\ensuremath{c}}\fi}
\newunicodechar{ʰ}{\ifmmode^{h}\else\textsuperscript{\ensuremath{h}}\fi}
\newunicodechar{ᵠ}{\ifmmode^{q}\else\textsuperscript{\ensuremath{q}}\fi}
\newunicodechar{ₙ}{\ifmmode_{n}\else\textsubscript{\ensuremath{n}}\fi}
\newunicodechar{ₐ}{\ifmmode_{a}\else\textsubscript{\ensuremath{a}}\fi}
\newunicodechar{ᵢ}{\ifmmode_{i}\else\textsubscript{\ensuremath{i}}\fi}
\newunicodechar{ᵣ}{\ifmmode_{r}\else\textsubscript{\ensuremath{r}}\fi}
\newunicodechar{ₖ}{\ifmmode_{k}\else\textsubscript{\ensuremath{k}}\fi}
\newunicodechar{ᵦ}{\ifmmode_{\beta}\else\textsubscript{\ensuremath{\beta}}\fi}
\newunicodechar{ₓ}{\ifmmode_{x}\else\textsubscript{\ensuremath{x}}\fi}
\newfontfamily\popdisplay{ArchivoBlack-Regular.ttf}[Path=../../../fonts/]
\newfontfamily\poplogo{SpaceGrotesk-Bold.ttf}[Path=../../../fonts/]
\newfontfamily\popsemi{PlusJakartaSans-SemiBold.ttf}[Path=../../../fonts/]
\newfontfamily\popxbold{PlusJakartaSans-ExtraBold.ttf}[Path=../../../fonts/]
\newfontfamily\popxboldls{PlusJakartaSans-ExtraBold.ttf}[Path=../../../fonts/,LetterSpace=3.0]

\setlist[enumerate]{leftmargin=1.7em,itemsep=5pt,topsep=4pt}
\setlist[itemize]{leftmargin=1.7em,itemsep=4pt,topsep=4pt,label={\color{poporange}\textbullet}}
\setlength{\parindent}{0pt}
\setlength{\parskip}{5pt}
\renewcommand{\arraystretch}{1.25}

% pgfplots : style Pop par défaut
\pgfplotsset{
  every axis/.append style={axis line style={popink,line width=1pt},tick style={popink},
    grid style={popline,line width=0.5pt},label style={font=\small},tick label style={font=\footnotesize}},
  popcurve/.style={poporange,line width=1.8pt,no markers,smooth},
}
% Même style disponible dans un \draw TikZ simple (hors pgfplots)
\tikzset{popcurve/.style={poporange,line width=1.8pt}}
\tikzset{popgrid/.style={popline,very thin}}
\colorlet{popcurve}{poporange} % le nom du style est parfois utilisé comme couleur

% ============================================================
% Pill / squiggle
% ============================================================
\newcommand{\poppill}[3]{%
  \tikz[baseline=-0.55ex]\node[draw=popink,line width=1.2pt,rounded corners=2.6pt,%
    fill=#2,inner xsep=6.5pt,inner ysep=3pt]{%
    \popxboldls\fontsize{7}{7}\selectfont\color{#3}\MakeUppercase{#1}};%
}
\newcommand{\popsquiggle}[1]{%
  \tikz[baseline=-0.4ex]{
    \draw[#1,line width=2.6pt,line cap=round]
      plot [smooth] coordinates {(0,0.09) (0.34,0.01) (0.68,0.21) (1.02,0.01) (1.36,0.10)};
  }%
}
% Mot-clé surligné (marqueur orange sous le texte)
\newcommand{\cle}[1]{{\popxbold\color{popdark}#1}}

% ============================================================
% En-tête / pied
% ============================================================
\pagestyle{fancy}
\fancyhf{}
\renewcommand{\headrulewidth}{0pt}
\renewcommand{\footrulewidth}{0pt}
\lhead{\raisebox{-0.15ex}{\poplogo\fontsize{13}{13}\selectfont\color{popink}OnBuch\color{poporange}+}}
\rhead{\poppill{\DOCMATIERE\ \textbullet\ \DOCNIVEAU\ \textbullet\ Leçon \DOCLECON}{white}{popink}}
\lfoot{\popxbold\fontsize{7.6}{7.6}\selectfont\color{popmuted}\MakeUppercase{Cours signé OnBuch+}\ \textbullet\ {\popsemi\DOCTITRE}}
\rfoot{\tikz[baseline=-0.6ex]\node[rounded corners=8.5pt,fill=popink,minimum height=17pt,minimum width=17pt,inner xsep=5pt,inner sep=0]{\popxbold\fontsize{8}{8}\selectfont\color{popcream}\thepage\,/\,\pageref*{LastPage}};}

% ============================================================
% Titres
% ============================================================
\newcommand{\chaptitre}[2]{%
  \begin{tcolorbox}[enhanced,colback=poporange,colframe=popink,boxrule=2.6pt,arc=11pt,
      drop shadow=popink,left=15pt,right=15pt,top=12pt,bottom=11pt,before skip=3pt,after skip=18pt]
    \poppill{#2}{popink}{popcream}\par\vspace{9pt}
    {\raggedright\hyphenpenalty=10000\exhyphenpenalty=10000\popdisplay\fontsize{22}{24}\selectfont\color{popcream}\MakeUppercase{#1}\par}
  \end{tcolorbox}
}
% Section numérotée : pastille orange avec le numéro + titre Archivo
\newcounter{popsec}
\newcommand{\coursec}[1]{%
  \stepcounter{popsec}\setcounter{popsub}{0}%
  \par\vspace{16pt}\noindent
  \begin{minipage}{\linewidth}
  \tikz[baseline=-0.7ex]\node[draw=popink,line width=1.6pt,fill=poporange,rounded corners=4pt,minimum size=22pt,inner sep=0]{\popdisplay\fontsize{12}{12}\selectfont\color{popcream}\thepopsec};%
  \hspace{8pt}{\popdisplay\fontsize{15}{17}\selectfont\color{popink}\MakeUppercase{#1}}\par
  \vspace{4pt}\noindent{\color{poporange}\rule{36pt}{3.6pt}}
  \end{minipage}\par\nopagebreak\vspace{8pt}
}
\newcounter{popsub}
\newcommand{\courssub}[1]{%
  \stepcounter{popsub}%
  \par\vspace{9pt}\noindent{\popxbold\fontsize{11.5}{13}\selectfont\color{popink}{\color{poporange}\thepopsec.\thepopsub}\hspace{6pt}#1}\par\nopagebreak\vspace{4pt}
}

% ============================================================
% Boîtes pédagogiques — même recette : fond blanc, bordure épaisse colorée,
% ombre dure encre, badge plein. Titre optionnel [#1].
% ============================================================
\tcbset{popbase/.style={breakable,enhanced,colback=white,boxrule=2.3pt,arc=9pt,
  shadow={3pt}{-3pt}{0pt}{popink},left=14pt,right=14pt,top=11pt,bottom=11pt,before skip=11pt,after skip=15pt,
  fontupper=\popsemi\fontsize{10.6}{14.5}\selectfont\color{popink},
  fontlower=\popsemi\fontsize{10.6}{14.5}\selectfont\color{popink}}}
% Coche dessinée (le glyphe ✓ n'existe pas dans les polices du document)
\newcommand{\popcheck}[1]{\tikz[baseline=-0.5ex]\draw[#1,line width=1.6pt,line cap=round,line join=round](-0.13,0.0)--(-0.04,-0.09)--(0.13,0.1);}
\providecommand{\coche}{\popcheck{popgreen}}
\providecommand{\resultat}[1]{\par\smallskip\noindent\fcolorbox{popgreen}{popgreenL}{\parbox{\dimexpr\linewidth-2\fboxsep-2\fboxrule\relax}{\textbf{Résultat.} #1}}\par\smallskip}
\newcommand{\popboxhead}[3]{\poppill{#1}{#2}{white}\ifx\relax#3\relax\else\hspace{6pt}{\popxbold\fontsize{10.6}{12}\selectfont\color{popink}#3}\fi\par\vspace{7pt}}

\newtcolorbox{aretenir}[1][]{popbase,colframe=popgreen,before upper={\popboxhead{À retenir}{popgreen}{#1}}}
% Texte en anglais (document de lecture, dialogue, extrait) : cadre
% d'étude. Un titre optionnel (source, auteur) s'affiche dans l'en-tête.
\newtcolorbox{textebox}[1][]{popbase,colframe=popblue,before upper={\popboxhead{Text}{popblue}{#1}\begin{otherlanguage}{english}},after upper={\end{otherlanguage}}}
% Marques ✓ / ✗ (forme correcte / fautive) souvent utilisées par les modèles
\providecommand{\cross}{\textcolor{poppink}{\ensuremath{\times}}}
\providecommand{\xmark}{\cross}\providecommand{\crossmark}{\cross}\providecommand{\cmark}{\ensuremath{\checkmark}}
% Ligne de réponse / justification à compléter (inventée par les modèles)
\providecommand{\justification}[1]{\par\noindent\textit{Justification~:} \dotfill\par}
% \textipa (package tipa non chargé) : la police ne couvre pas l'API -> simple passage
\providecommand{\textipa}[1]{#1}
% Soulignement (ulem non chargé) : \uline, \ul
\providecommand{\uline}[1]{\underline{#1}}\providecommand{\ul}[1]{\underline{#1}}
% \glossaire{\item ...} inventée par les modèles : liste de glossaire
\providecommand{\glossaire}[1]{\begin{itemize}#1\end{itemize}}
% \alert (beamer) inventée par les modèles : texte mis en évidence
\providecommand{\alert}[1]{\textbf{\textcolor{poppink}{#1}}}
% variantes ASCII d'ordinaux écrites en macro
\providecommand{\iere}{\textsuperscript{re}}\providecommand{\ieme}{\textsuperscript{e}}\providecommand{\ere}{\textsuperscript{re}}\providecommand{\eme}{\textsuperscript{e}}\providecommand{\er}{\textsuperscript{er}}\providecommand{\ie}{\textsuperscript{e}}
% Ordinaux français écrits en macro par les modèles : 1\ière, 3\ième
\newcommand{\ière}{\textsuperscript{re}}\newcommand{\ième}{\textsuperscript{e}}
% \exemple (inventée par les modèles) : étiquette « Exemple : »
\providecommand{\exemple}{\textbf{Exemple~:}\ }
% \square utilisable aussi hors mode math (cases à cocher)
\AtBeginDocument{\let\squareMath\square\renewcommand{\square}{\ensuremath{\squareMath}}}
% Mot ou phrase en anglais à l'intérieur d'un paragraphe français
\newcommand{\eng}[1]{\ifmmode\text{\textit{#1}}\else\begingroup\selectlanguage{english}\itshape #1\endgroup\fi}
\let\esp\eng % alias toléré
% flèches d'intonation inventées par les modèles
\providecommand{\textsearrow}{}\renewcommand{\textsearrow}{\ensuremath{\searrow}}\providecommand{\textnearrow}{}\renewcommand{\textnearrow}{\ensuremath{\nearrow}}
% \fr{..} : mot français (inventé par les modèles)
\providecommand{\fr}[1]{\textit{#1}}
\newtcolorbox{exemplebox}[1][]{popbase,colframe=poporange,before upper={\popboxhead{Exemple}{poporange}{#1}}}
\newtcolorbox{definition}[1][]{popbase,colframe=popblue,before upper={\popboxhead{Définition}{popblue}{#1}}}
\newtcolorbox{propriete}[1][]{popbase,colframe=poppurple,before upper={\popboxhead{Propriété}{poppurple}{#1}}}
\newtcolorbox{methode}[1][]{popbase,colframe=popgold,before upper={\popboxhead{Méthode}{popgold}{#1}}}
\newtcolorbox{attention}[1][]{popbase,colframe=poppink,before upper={\popboxhead{Attention}{poppink}{#1}}}
\newtcolorbox{experience}[1][]{popbase,colframe=popblue,colback=popblueL,before upper={\popboxhead{Expérience}{popblue}{#1}}}
% « Le savais-tu ? » : carte violette pleine (contexte, culture, Cameroun)
\newtcolorbox{savaistu}[1][]{popbase,colback=poppurple,colframe=popink,
  fontupper=\popsemi\fontsize{10.4}{14.2}\selectfont\color{white},
  before upper={\poppill{Le savais-tu\,?}{white}{poppurple}\ifx\relax#1\relax\else\hspace{6pt}{\popxbold\color{white}#1}\fi\par\vspace{7pt}}}
% Exercice résolu : énoncé en haut, \tcblower puis la solution sur fond crème
\newcounter{popexo}\newcounter{popexr}
\newtcolorbox{exoresolu}[1][]{popbase,colframe=poporange,colbacklower=popcream,
  segmentation style={popink,line width=1.2pt,dashed},
  before upper={\stepcounter{popexr}\popboxhead{Exercice résolu \thepopexr}{poporange}{#1}},
  before lower={\poppill{Solution}{popink}{popcream}\par\vspace{6pt}}}
% Exercice (non résolu en place) — la correction est dans la partie Corrigés
\newtcolorbox{exercice}[1][]{popbase,colframe=popink,
  before upper={\stepcounter{popexo}\popboxhead{Exercice \thepopexo}{popink}{#1}}}
% Corrigé
\newtcolorbox{corrige}[1][]{popbase,colframe=popgreen,colback=popgreenL,
  before upper={\popboxhead{Corrigé}{popgreen}{#1}}}
% Objectifs / prérequis / bilan
\newtcolorbox{objectifs}{popbase,colframe=popink,colback=poporangeL,
  before upper={\popboxhead{Objectifs de la leçon}{popink}{}}}
\newtcolorbox{prerequis}{popbase,colframe=popink,colback=popblueL,
  before upper={\popboxhead{Prérequis}{popblue}{}}}
\newtcolorbox{bilan}{popbase,colframe=popink,colback=white,boxrule=2.8pt,
  before upper={\popboxhead{Fiche bilan}{poporange}{L'essentiel à connaître pour l'examen}}}
% Figure encadrée avec légende
\newtcolorbox{popfigure}[1][]{enhanced,breakable=false,colback=white,colframe=popline,boxrule=1.4pt,arc=8pt,
  left=10pt,right=10pt,top=10pt,bottom=8pt,before skip=10pt,after skip=12pt,halign=center,
  lower separated=false,halign lower=center,
  fontlower=\popsemi\fontsize{9}{11.5}\selectfont\color{popmuted},
  #1}
\newcommand{\legende}[1]{\par\vspace{6pt}{\popsemi\fontsize{9}{11.5}\selectfont\color{popmuted}#1}}

% ============================================================
% Signature OnBuch+
% ============================================================
\newcommand{\poplogoinline}[1]{{\poplogo\fontsize{#1}{#1}\selectfont\color{popink}OnBuch\color{poporange}+}}
\newcommand{\signatureonbuch}{%
  \begin{tcolorbox}[enhanced,colback=popink,colframe=popink,boxrule=2.2pt,arc=10pt,
      drop shadow=poporange,left=14pt,right=14pt,top=10pt,bottom=10pt,before skip=0pt,after skip=0pt]
    \tikz[baseline=-0.7ex]\node[circle,fill=popgreen,draw=popcream,line width=1.3pt,minimum size=22pt,inner sep=0]{\popcheck{white}};%
    \hspace{9pt}\begin{minipage}[c]{0.62\linewidth}
      {\popxbold\fontsize{12}{13}\selectfont\color{popcream}Signé\ \textbullet\ {\poplogo OnBuch\color{poporange}+}}\par\vspace{2pt}
      {\raggedright\popsemi\fontsize{8}{10.5}\selectfont\color{popline}\MakeUppercase{Équipe pédagogique OnBuch+}\\\MakeUppercase{Conforme au programme officiel MINESEC}\par}
    \end{minipage}\hfill
    {\poplogo\fontsize{22}{22}\selectfont\color{popcream}OnBuch\color{poporange}+}
  \end{tcolorbox}%
}

% ============================================================
% Page de garde
% #1 titre · #2 matière · #3 niveau · #4 module · #5 n° leçon · #6 durée ·
% #7 description / objectifs (texte ou itemize)
% ============================================================
\newcommand{\pagegarde}[7]{%
  \thispagestyle{empty}
  \newgeometry{a4paper,margin=1.7cm,top=1.6cm,bottom=1.4cm}%
  \begin{tikzpicture}[remember picture,overlay]
    \fill[poporange!18] ([xshift=-3.2cm,yshift=-3.6cm]current page.north east) circle (4.6cm);
    \fill[poppurple!14] ([xshift=2.4cm,yshift=7.6cm]current page.south west) circle (3.8cm);
  \end{tikzpicture}%
  {\poplogo\fontsize{30}{30}\selectfont\color{popink}OnBuch\color{poporange}+}\hfill
  \raisebox{0.6ex}{\poppill{Cours signé \textbullet\ Premium}{popink}{popcream}}\par
  \vspace{1pt}\popsquiggle{poppurple}\par
  \vspace{8pt}
  \begin{tcolorbox}[enhanced,colback=poporange,colframe=popink,boxrule=2.8pt,arc=13pt,
      drop shadow=popink,left=18pt,right=18pt,top=12pt,bottom=13pt,after skip=10pt]
    \poppill{#2 \textbullet\ #3}{popink}{popcream}\hspace{5pt}\poppill{#4}{white}{popink}\par\vspace{8pt}
    {\popdisplay\fontsize{14}{14}\selectfont\color{popink}LEÇON #5}\par\vspace{4pt}
    {\raggedright\hyphenpenalty=10000\exhyphenpenalty=10000\popdisplay\fontsize{25}{27}\selectfont\color{popcream}\MakeUppercase{#1}\par}
    \vspace{8pt}
    {\popxbold\fontsize{9.5}{9.5}\selectfont\color{popink}\MakeUppercase{Durée conseillée : #6}}
  \end{tcolorbox}
  \begin{tcolorbox}[enhanced,colback=white,colframe=popink,boxrule=2.2pt,arc=10pt,
      drop shadow=popink,left=16pt,right=16pt,top=10pt,bottom=10pt,after skip=8pt]
    \poppill{Dans cette leçon}{poppurple}{white}\par\vspace{5pt}
    \popsemi\fontsize{9.3}{12.6}\selectfont\color{popink}#7
  \end{tcolorbox}
  \noindent\begin{minipage}[t]{0.487\linewidth}
    \begin{tcolorbox}[enhanced,colback=popgreen,colframe=popink,boxrule=2.2pt,arc=10pt,
        drop shadow=popink,left=11pt,right=11pt,top=10pt,bottom=10pt,height=3.1cm,valign=center]
      \tcbox[colback=white,colframe=popink,boxrule=1.3pt,arc=5pt,boxsep=0pt,nobeforeafter,
        left=3pt,right=3pt,top=3pt,bottom=3pt,valign=center]{\qrcode[height=1.9cm]{https://play.google.com/store/apps/details?id=com.luvvix.onbuch&hl=fr}}%
      \hspace{8pt}\begin{minipage}[c]{0.5\linewidth}
        {\popdisplay\fontsize{11}{12.5}\selectfont\color{white}\MakeUppercase{Télécharge OnBuch}}\par\vspace{4pt}
        {\raggedright\popsemi\fontsize{8.4}{11}\selectfont\color{white}Gratuit \textbullet\ Google Play\\Tuteur IA, annales, concours, résultats\par}
      \end{minipage}
    \end{tcolorbox}
  \end{minipage}\hfill
  \begin{minipage}[t]{0.487\linewidth}
    \begin{tcolorbox}[enhanced,colback=poppurple,colframe=popink,boxrule=2.2pt,arc=10pt,
        drop shadow=popink,left=11pt,right=11pt,top=10pt,bottom=10pt,height=3.1cm,valign=center]
      \tikz[baseline=-0.7ex]\node[circle,fill=white,draw=popink,line width=1.3pt,minimum size=1.6cm,inner sep=0]{\popdisplay\fontsize{24}{24}\selectfont\color{poppurple}+};%
      \hspace{8pt}\begin{minipage}[c]{0.6\linewidth}
        {\popdisplay\fontsize{11}{12.5}\selectfont\color{white}\MakeUppercase{Passe à OnBuch+}}\par\vspace{4pt}
        {\raggedright\popsemi\fontsize{8.4}{11}\selectfont\color{white}Tous les cours signés, Tuteur IA illimité, fiches PDF — onglet OnBuch+ de l'app\par}
      \end{minipage}
    \end{tcolorbox}
  \end{minipage}
  \vspace{8pt}
  \popcontenu
  \vfill
  \signatureonbuch
  \restoregeometry
  \newpage
}

% Rangée « Ce cours contient » (page de garde)
\newcommand{\poptile}[3]{% #1 couleur #2 gros texte #3 légende
  \begin{tcolorbox}[enhanced,colback=#1,colframe=popink,boxrule=2pt,arc=9pt,shadow={2.5pt}{-2.5pt}{0pt}{popink},
      left=6pt,right=6pt,top=9pt,bottom=9pt,halign=center,valign=center,height=1.75cm,width=\linewidth,nobeforeafter]
    {\popdisplay\fontsize{12.5}{13}\selectfont\color{white}\MakeUppercase{#2}}\par\vspace{3pt}
    {\centering\popsemi\fontsize{7.8}{9.5}\selectfont\color{white}#3\par}
  \end{tcolorbox}}
\newcommand{\popcontenu}{%
  \par\noindent{\popxboldls\fontsize{7.5}{7.5}\selectfont\color{popmuted}CE COURS SIGNÉ CONTIENT}\par\vspace{7pt}
  \noindent\begin{minipage}[t]{0.235\linewidth}\poptile{popblue}{Cours}{complet, illustré, pas à pas}\end{minipage}\hfill
  \begin{minipage}[t]{0.235\linewidth}\poptile{poporange}{Méthodes}{et exercices résolus}\end{minipage}\hfill
  \begin{minipage}[t]{0.235\linewidth}\poptile{poppink}{Exercices}{type examen + corrigés}\end{minipage}\hfill
  \begin{minipage}[t]{0.235\linewidth}\poptile{popgreen}{Fiche bilan}{l'essentiel à retenir}\end{minipage}\par}

% Sommaire
\newtcolorbox{sommaire}{enhanced,colback=white,colframe=popink,boxrule=2.2pt,arc=10pt,
  drop shadow=popink,left=18pt,right=18pt,top=13pt,bottom=13pt,before skip=6pt,after skip=20pt,
  before upper={\poppill{Sommaire}{poppurple}{white}\par\vspace{9pt}\popsemi\fontsize{10.6}{14.5}\selectfont\color{popink}}}

% Signature de fin de cours — poussée tout en bas de la dernière page
\newcommand{\findecours}{%
  \par\vfill\nopagebreak
  \begin{center}\popsquiggle{poporange}\end{center}\vspace{4pt}
  \signatureonbuch
}
\AtBeginDocument{\def\llbracket{\mathopen{[\mkern-3.5mu[}}\def\rrbracket{\mathclose{]\mkern-3.5mu]}}} % intervalles d'entiers (glyphe ⟦ absent de la police)
% Glyphes absents de Fira Math : redéfinis à partir de glyphes disponibles
\AtBeginDocument{%
  \def\setminus{\mathbin{\backslash}}\let\smallsetminus\setminus
  \def\vdots{\mathord{\vbox{\baselineskip3.2pt\lineskiplimit0pt\kern4pt\hbox{.}\hbox{.}\hbox{.}}}}%
  \def\ddots{\mathinner{\mkern1mu\raise6.5pt\hbox{.}\mkern2mu\raise3.6pt\hbox{.}\mkern2mu\raise0.7pt\hbox{.}\mkern1mu}}%
  \def\triangle{\mathord{\scalebox{0.9}{$\symup{\Delta}$}}}%
  \def\nabla{\mathord{\raisebox{\depth}{\scalebox{1}[-1]{$\symup{\Delta}$}}}}%
  \def\checkmark{\popcheck{popdark}}%
  \def\star{\mathbin{\ast}}\def\bigstar{\mathord{\ast}}%
  \def\diamond{\mathbin{\scalebox{0.6}{$\Diamond$}}}%
  \def\bigcup{\mathop{\mathchoice{\raisebox{-0.3em}{\scalebox{1.7}{$\cup$}}}{\scalebox{1.25}{$\cup$}}{\cup}{\cup}}\displaylimits}%
  \def\bigcap{\mathop{\mathchoice{\raisebox{-0.3em}{\scalebox{1.7}{$\cap$}}}{\scalebox{1.25}{$\cap$}}{\cap}{\cap}}\displaylimits}%
  \def\lesseqqgtr{\mathrel{\lessgtr}}\def\bigoplus{\mathop{\oplus}\displaylimits}%
}
\providecommand{\ssi}{si et seulement si} % abréviation fréquente
\AtBeginDocument{\providecommand{\wideparen}{\overparen}} % arc de cercle

%% FIGFIX-BEGIN v4 — correctifs globaux des figures (docs/onbuch-generation/tools/figfix)
\usepackage{adjustbox}\usepackage{etoolbox}
% toute figure TikZ/pgfplots plus large que la ligne est réduite à la largeur disponible
\BeforeBeginEnvironment{tikzpicture}{\begin{adjustbox}{max width=\linewidth}}
\AfterEndEnvironment{tikzpicture}{\end{adjustbox}}
% légendes pgfplots lisibles par-dessus les courbes
\pgfplotsset{every axis/.append style={legend style={fill=white,fill opacity=0.92,text opacity=1,draw=black!15,inner sep=3pt}}}
% étiquettes d'arêtes (syntaxe edge["..."]) avec fond blanc
\tikzset{every edge quotes/.append style={fill=white,inner sep=1.2pt}}
%% FIGFIX-END
\begin{document}

\pagegarde{Lesson 3 — Vocabulary: Words and expressions related to interacting with mates, neighbours, adults and school authorities on community celebrations/festivals}{Anglais}{1ère A}{Chapter 1 : Family and Social Life}{ENA03}{2 h}{Cette leçon de vocabulaire permet aux élèves de 1ère A d'acquérir et d'employer le lexique des interactions sociales (camarades, voisins, adultes, autorités scolaires) autour des fêtes et festivals communautaires. Elle couvre les champs lexicaux, les collocations, les expressions idiomatiques, les faux amis et la prononciation, avec réemploi en contexte oral et écrit.
\vspace{4pt}
\begin{multicols}{2}\small
\begin{itemize}
\item À la fin de cette leçon, je sais nommer les principales fêtes et célébrations communautaires en anglais
\item je sais utiliser le vocabulaire approprié selon l'interlocuteur (mate, neighbour, adult, school authority)
\item je sais employer les collocations courantes liées aux fêtes (to throw a party, to attend a ceremony...)
\item je sais utiliser des expressions idiomatiques pour inviter, féliciter et remercier
\item je sais éviter les faux amis et les calques du français (assist, ceremony, eventually...)
\item je sais former et utiliser des mots de la même famille (celebrate, celebration, celebrant...)
\end{itemize}\end{multicols}}

\begin{sommaire}
\begin{multicols}{2}
\begin{enumerate}
\item Community celebrations and festivals: core vocabulary
\item Interacting with different people: register and politeness
\item Collocations and idiomatic expressions
\item False friends and common traps for francophones
\item Word families, pronunciation and spelling
\item Using the vocabulary in context
\item Activité d'intégration
\item Méthodes et exercices résolus
\item Exercices
\item Corrigés des exercices
\item Fiche bilan
\end{enumerate}
\end{multicols}
\end{sommaire}

\begin{objectifs}
\begin{itemize}
\item À la fin de cette leçon, je sais nommer les principales fêtes et célébrations communautaires en anglais
\item je sais utiliser le vocabulaire approprié selon l'interlocuteur (mate, neighbour, adult, school authority)
\item je sais employer les collocations courantes liées aux fêtes (to throw a party, to attend a ceremony...)
\item je sais utiliser des expressions idiomatiques pour inviter, féliciter et remercier
\item je sais éviter les faux amis et les calques du français (assist, ceremony, eventually...)
\item je sais former et utiliser des mots de la même famille (celebrate, celebration, celebrant...)
\item je sais prononcer correctement le lexique clé et respecter son orthographe
\item je sais réemployer ce lexique dans des phrases, dialogues et courts textes
\end{itemize}
\end{objectifs}

\begin{prerequis}
\begin{itemize}
\item Vocabulaire de base de la famille et de la vie sociale (Seconde, Chapter: Family and Social Life)
\item Les salutations et formules de politesse simples (Hello, How are you?, Thank you)
\item Le présent simple et le présent continu pour décrire des habitudes et des événements
\item La notion de registre de langue (familier / soutenu) vue en français et en Seconde
\end{itemize}
\end{prerequis}

\newpage

\coursec{Community celebrations and festivals: core vocabulary}

\textbf{Situation d'accroche.} It's December in Bamenda. The whole neighbourhood is preparing the end-of-year community festival: your mates are organising the football match, your neighbour Ma Grace is cooking eru for fifty guests, and the school principal has been invited as guest of honour. You want to invite your English penfriend, greet the elders politely, and ask the principal for permission to use the school field. But how do you speak to a mate, to an adult, and to a school authority in English? The words change with the person! This lesson gives you all the vocabulary you need to interact about community celebrations and festivals.

\textbf{Problématique.} Quels mots anglais utiliser pour nommer les fêtes, les organisateurs, les invités et les activités d'une célébration communautaire ? Comment varier son vocabulaire selon la personne à qui l'on parle ? C'est tout l'objet de cette première section : bâtir un socle lexical solide, correctement prononcé et correctement orthographié.

\courssub{Discovering the vocabulary through a text}

Avant toute liste de mots, lisons un court texte authentique. Soulignez mentalement tous les mots qui se rapportent aux fêtes.

\begin{textebox}[A harvest festival in our community]
Our community organises a harvest festival every December. The chief and the elders are the guests of honour. Neighbours cook traditional dishes, mates organise games, and the school choir performs songs.

This year, the festival will take place on the school field. The organising committee met last week to plan the event. They scheduled the opening ceremony for nine o'clock in the morning. The principal of Government High School kindly agreed to host the first part of the celebration. He will give a short speech to welcome everybody.

On the day before the festival, the young people decorate the square with palm branches and coloured cloth. The women prepare achu, ndolé and roasted plantains, while the men set up the tents and the chairs. A local businessman has agreed to sponsor the event: he is paying for the drinks and the sound system.

During the celebration, dance groups perform traditional dances and a famous artist sings. Children parade through the streets wearing beautiful costumes. At noon, the whole community shares a meal together. In the evening, families offer gifts to the chief and thank the organisers. Everybody agrees: it is the happiest gathering of the year, and nobody wants to miss the next anniversary of the festival.
\end{textebox}

\textbf{Glossaire.}

\begin{center}
\begin{tabularx}{\linewidth}{|l|l|X|}
\hline
\rowcolor{popblueL} \textbf{Mot anglais} & \textbf{Nature} & \textbf{Traduction} \\
\hline
harvest & noun & récolte \\
\hline
chief & noun & chef (traditionnel) \\
\hline
elder & noun & ancien, aîné \\
\hline
guest of honour & noun phrase & invité d'honneur \\
\hline
to schedule & verb & programmer, fixer au calendrier \\
\hline
to host & verb & accueillir (un événement) \\
\hline
to sponsor & verb & sponsoriser, financer \\
\hline
to parade & verb & défiler \\
\hline
gathering & noun & rassemblement, réunion \\
\hline
to miss & verb & manquer, rater \\
\hline
\end{tabularx}
\end{center}

\textbf{Comprehension questions.}

\begin{enumerate}
\item When does the community organise its harvest festival?
\item Who are the guests of honour?
\item Where will the festival take place this year?
\item What is the local businessman paying for?
\item Name two activities that happen during the celebration.
\end{enumerate}

\courssub{The lexical field of celebrations}

Observons d'abord les mots du texte qui \emph{nomment} les événements : \eng{festival}, \eng{celebration}, \eng{ceremony}, \eng{event}, \eng{gathering}, \eng{anniversary}. Tous désignent des moments festifs, mais ils ne sont pas interchangeables.

\begin{definition}[Festival]
A \cle{festival} is a series of public events (music, dance, food) celebrating a special occasion, often over several days. \eng{The Ngondo is a famous festival of the Sawa people.} « Le Ngondo est un festival célèbre du peuple sawa. »
\end{definition}

\begin{definition}[Ceremony]
A \cle{ceremony} is a formal act or a series of formal acts performed on important occasions (weddings, graduations, openings). \eng{The opening ceremony starts at nine o'clock.} « La cérémonie d'ouverture commence à neuf heures. »
\end{definition}

\begin{definition}[Celebration, feast, party, gathering, event, anniversary]
\begin{itemize}
\item \cle{celebration} : une fête, une célébration (le fait de fêter quelque chose). \eng{The wedding celebration lasted two days.}
\item \cle{feast} : un festin, un grand repas festif ; aussi une fête religieuse. \eng{They prepared a huge feast for the village.}
\item \cle{party} : une fête privée, entre amis ou en famille. \eng{We are having a birthday party on Saturday.}
\item \cle{gathering} : un rassemblement de personnes. \eng{The festival is the biggest gathering of the year.}
\item \cle{event} : un événement (terme général). \eng{The football match is the main event of the day.}
\item \cle{anniversary} : un anniversaire (de mariage, d'indépendance, d'un événement). \eng{The school celebrated its twentieth anniversary.}
\end{itemize}
\end{definition}

\begin{propriete}[Types de fêtes : Cameroun et monde anglophone]
\begin{itemize}
\item \cle{traditional festival} : festival traditionnel. \eng{The Nguon is a traditional festival of the Bamoun people.}
\item \cle{cultural week} : semaine culturelle. \eng{Our school organises a cultural week in March.}
\item \cle{harvest festival} : fête de la récolte. \eng{The harvest festival takes place in December.}
\item \cle{wedding ceremony} : cérémonie de mariage. \eng{The wedding ceremony was held in the church.}
\item \cle{funeral celebration} : célébration funéraire (au Cameroun, les funérailles sont souvent festives). \eng{Many relatives came for the funeral celebration.}
\item \cle{birthday party} : fête d'anniversaire. \eng{She invited thirty friends to her birthday party.}
\item \cle{Christmas} : Noël. \eng{At Christmas, families share a meal and offer gifts.}
\item \cle{New Yam Festival} : fête de l'igname nouvelle (très répandue en Afrique de l'Ouest anglophone, notamment au Nigeria). \eng{The New Yam Festival marks the beginning of the harvest season.}
\item \cle{Ngondo} : festival traditionnel du peuple Sawa (Douala), au bord du fleuve Wouri.
\item \cle{Nguon} : grande fête traditionnelle du peuple Bamoun (Foumban), autour du sultan.
\end{itemize}
\end{propriete}

\begin{attention}[Faux amis et calques]
\begin{itemize}
\item \eng{anniversary} ne signifie pas « anniversaire » au sens de date de naissance : on dit \eng{birthday}. \eng{Happy birthday!} et non « Happy anniversary! » pour un anniversaire de naissance.
\item \eng{feast} est plus fort que « fête » : c'est un \emph{festin}. Pour « faire la fête », dites \eng{to celebrate} ou \eng{to have a party}.
\item Ne dites pas « assist to the ceremony » (calque du français « assister à ») : dites \eng{to attend the ceremony}.
\end{itemize}
\end{attention}

\courssub{The lexical field of organisation}

\begin{definition}[Organising an event]
\begin{itemize}
\item \cle{to organise} : organiser. \eng{My mates are organising the football match.}
\item \cle{to plan} : planifier, prévoir. \eng{The committee planned the event carefully.}
\item \cle{to prepare} : préparer. \eng{Ma Grace is preparing eru for fifty guests.}
\item \cle{to decorate} : décorer. \eng{The pupils decorate the square with flowers.}
\item \cle{to invite} : inviter. \eng{We invited the principal and the teachers.}
\item \cle{to host} : accueillir (un événement, des invités). \eng{Our school will host the cultural week.}
\item \cle{to sponsor} : sponsoriser, financer. \eng{A local company sponsors the festival every year.}
\item \cle{to schedule} : programmer, fixer la date. \eng{They scheduled the ceremony for nine o'clock.}
\end{itemize}
\end{definition}

\begin{attention}[« Organiser » : prononciation et orthographe]
En anglais, \eng{to organise} (britannique) ou \eng{to organize} (américain) s'écrit avec \textbf{-ise/-ize}, jamais avec un « s » comme en français « organiser ». Le nom est \eng{organisation} et la personne est \eng{an organiser}. Prononciation : « o-ga-naï-zeur » pour \eng{organiser}.
\end{attention}

\courssub{The lexical field of participants}

\begin{definition}[People at a celebration]
\begin{itemize}
\item \cle{guest} : invité(e). \eng{Fifty guests came to the party.}
\item \cle{host} : hôte, celui qui reçoit. \eng{The host welcomed everybody warmly.}
\item \cle{organiser} : organisateur(-trice). \eng{The organisers worked for three months.}
\item \cle{participant} : participant(e). \eng{Each participant received a certificate.}
\item \cle{guest of honour} : invité(e) d'honneur. \eng{The principal is our guest of honour.}
\item \cle{chief} : chef traditionnel. \eng{The chief blessed the new yams.}
\item \cle{elder} : ancien, aîné. \eng{We greeted the elders respectfully.}
\item \cle{invitee} : personne invitée (terme formel, moins courant que \eng{guest}). \eng{All invitees received a written invitation.}
\end{itemize}
\end{definition}

\begin{attention}[Host : attention au genre et au sens]
\eng{host} désigne celui qui \emph{reçoit}, pas celui qui est invité. La forme féminine traditionnelle est \eng{hostess}, mais \eng{host} est aujourd'hui utilisé pour les deux. Ne confondez pas avec le français « hôte », qui peut désigner l'invité : en anglais, l'invité est toujours \eng{guest}.
\end{attention}

\courssub{The lexical field of activities}

\begin{definition}[What people do at a festival]
\begin{itemize}
\item \cle{to dance} : danser. \eng{The dance groups danced all night.}
\item \cle{to sing} : chanter. \eng{The school choir sang three songs.}
\item \cle{to perform} : se produire, donner un spectacle. \eng{A famous artist performed at the festival.}
\item \cle{to parade} : défiler. \eng{Children paraded through the streets.}
\item \cle{to share a meal} : partager un repas. \eng{The whole community shared a meal at noon.}
\item \cle{to give a speech} : prononcer un discours. \eng{The principal gave a short speech.}
\item \cle{to offer gifts} : offrir des cadeaux. \eng{Families offered gifts to the chief.}
\end{itemize}
\end{definition}

\begin{propriete}[Tableau récapitulatif : événement et personnes impliquées]
\begin{center}
\begin{tabularx}{\linewidth}{|X|X|}
\hline
\rowcolor{popblueL} \textbf{Event (événement)} & \textbf{People involved (personnes impliquées)} \\
\hline
traditional festival (Ngondo, Nguon) & chief, elders, dancers, the whole community \\
\hline
harvest festival & farmers, organisers, guests of honour \\
\hline
wedding ceremony & bride and groom, families, guests, pastor \\
\hline
birthday party & host, mates, neighbours, guests \\
\hline
cultural week & principal, teachers, pupils, artists \\
\hline
funeral celebration & relatives, elders, mourners, the community \\
\hline
\end{tabularx}
\end{center}
\end{propriete}

\begin{popfigure}
\begin{tikzpicture}[mindmap, grow cyclic, every node/.style={concept}, root concept/.append style={concept color=popblue, font=\small\bfseries, text=white}, level 1/.append style={level distance=3.4cm, sibling angle=90}, level 2/.append style={level distance=2.2cm, sibling angle=45}, concept color=popblue, text=white]
\node[root concept]{COMMUNITY CELEBRATIONS}
  child[concept color=poporange] { node {Types of events}
    child { node[concept color=poporangeL, text=popdark, font=\scriptsize] {festival, ceremony}
    }
    child { node[concept color=poporangeL, text=popdark, font=\scriptsize] {party, feast, gathering}
    }
  }
  child[concept color=popgreen] { node {People}
    child { node[concept color=popgreenL, text=popdark, font=\scriptsize] {host, guest, invitee}
    }
    child { node[concept color=popgreenL, text=popdark, font=\scriptsize] {chief, elder, organiser}
    }
  }
  child[concept color=poppurple] { node {Organisation}
    child { node[concept color=poppurpleL, text=popdark, font=\scriptsize] {plan, organise, schedule}
    }
    child { node[concept color=poppurpleL, text=popdark, font=\scriptsize] {invite, host, sponsor}
    }
  }
  child[concept color=poppink] { node {Activities}
    child { node[concept color=poppinkL, text=popdark, font=\scriptsize] {dance, sing, perform}
    }
    child { node[concept color=poppinkL, text=popdark, font=\scriptsize] {parade, give a speech}
    }
  };
\end{tikzpicture}
\legende{Carte mentale du vocabulaire des célébrations communautaires : quatre champs lexicaux à mémoriser ensemble.}
\end{popfigure}

\courssub{Guided pronunciation of key words}

\begin{propriete}[Prononciation et orthographe des mots clés]
\begin{center}
\begin{tabularx}{\linewidth}{|l|X|X|}
\hline
\rowcolor{popblueL} \textbf{Mot} & \textbf{Prononciation (lettres françaises)} & \textbf{Traduction} \\
\hline
celebration & « sé-li-bré-cheune » (4 syllabes, accent sur « bré ») & célébration \\
\hline
ceremony & « sé-ri-mo-ni » (accent sur la 1\up{re} syllabe) & cérémonie \\
\hline
neighbour & « né-beur » (le \emph{gh} est muet) & voisin \\
\hline
authorities & « o-zo-ri-tiz » & autorités \\
\hline
festival & « fès-ti-vel » & festival \\
\hline
guest & « ghest » (un seul son, le \emph{u} est muet) & invité \\
\hline
chief & « tchif » & chef \\
\hline
elder & « èl-deur » & aîné, ancien \\
\hline
speech & « spitch » & discours \\
\hline
choir & « kouaïeur » (orthographe piège !) & chœur, chorale \\
\hline
\end{tabularx}
\end{center}
\end{propriete}

\begin{attention}[Ceremony : trois ou quatre syllabes ?]
En anglais britannique, \eng{ceremony} se prononce « sé-ri-mo-ni » : la dernière syllabe est réduite, contrairement au français « cé-ré-mo-nie » où toutes les syllabes sont nettes. L'accent tonique tombe sur la \textbf{première} syllabe, pas sur l'avant-dernière. De même, \eng{celebration} porte l'accent sur « bré » : « sé-li-BRÉ-cheune ». Un mauvais accent tonique rend le mot difficile à comprendre pour un anglophone.
\end{attention}

\begin{methode}[Apprendre le vocabulaire par champs lexicaux]
\begin{enumerate}
\item \textbf{Regroupez} les mots par thème : \emph{people} (host, guest, elder), \emph{food} (feast, dishes, meal), \emph{activities} (dance, parade, perform), \emph{organisation} (plan, invite, sponsor).
\item \textbf{Illustrez} chaque mot par une phrase personnelle, liée à votre vie : \eng{My neighbour Ma Grace hosts a party every Christmas.}
\item \textbf{Dessinez} une carte mentale comme celle de cette leçon pour visualiser les familles de mots.
\item \textbf{Réemployez} le lexique à l'oral : décrivez à un camarade la dernière fête de votre quartier en utilisant au moins huit mots de la leçon.
\item \textbf{Vérifiez} la prononciation à voix haute, en marquant l'accent tonique.
\end{enumerate}
\end{methode}

\begin{exoresolu}[Vocabulary in context]
Énoncé. Complétez les phrases avec le mot qui convient, choisi dans la liste : \eng{host}, \eng{guests}, \eng{ceremony}, \eng{sponsor}, \eng{parade}, \eng{speech}.

\begin{enumerate}
\item The opening \dots\ of the festival starts at nine o'clock.
\item Our school will \dots\ the cultural week this year.
\item A local bank agreed to \dots\ the event and pay for the tents.
\item Fifty \dots\ attended Ma Grace's birthday party.
\item The principal gave a short \dots\ to welcome the parents.
\item Children \dots\ through the streets in beautiful costumes.
\end{enumerate}

\tcblower
Solution détaillée.

\begin{enumerate}
\item \eng{ceremony} — « cérémonie » : on parle de la cérémonie d'ouverture (\eng{opening ceremony}).
\item \eng{host} — « accueillir » : l'école accueille l'événement ; attention, \eng{host} est ici un verbe.
\item \eng{sponsor} — « sponsoriser » : la banque finance l'événement.
\item \eng{guests} — « invités » : pluriel avec \eng{fifty}.
\item \eng{speech} — « discours » : collocation \eng{to give a speech}.
\item \eng{parade} — « défilent » : verbe conjugué au présent avec \eng{children}.
\end{enumerate}

Piège à éviter : à la question 4, ne mettez pas \eng{hosts} : ce sont les invités qui viennent à une fête, pas les hôtes. Le \eng{host} est la personne qui reçoit (ici, Ma Grace).
\end{exoresolu}

\begin{savaistu}[Traditional festivals in Cameroon]
Au Cameroun anglophone et dans tout le pays, les \eng{traditional festivals} jouent un rôle social immense. Le \cle{Ngondo}, festival du peuple Sawa, se déroule au bord du fleuve Wouri à Douala : courses de pirogues, danses et rituels rassemblent toute la communauté autour des chefs traditionnels. Le \cle{Nguon}, grande fête du peuple Bamoun à Foumban, célèbre la monarchie et la culture bamoun depuis des siècles. Ces festivals montrent qu'en anglais comme en français, une \eng{celebration} est d'abord un moment de \emph{communauté} : \eng{the whole community gathers around its chiefs and elders}.
\end{savaistu}

\begin{aretenir}[L'essentiel de la section]
\begin{itemize}
\item Un \eng{festival} est une série d'événements publics ; une \eng{ceremony} est un acte formel ; une \eng{party} est privée ; un \eng{gathering} est un rassemblement.
\item On \eng{organises}, \eng{plans}, \eng{schedules}, \eng{hosts} et \eng{sponsors} un événement ; on \eng{invites} des \eng{guests}.
\item Le \eng{host} reçoit, le \eng{guest} est reçu ; le \eng{chief} et les \eng{elders} sont les \eng{guests of honour}.
\item Pendant la fête, on \eng{dances}, \eng{sings}, \eng{performs}, \eng{parades}, on \eng{shares a meal}, on \eng{gives a speech} et on \eng{offers gifts}.
\item Prononciation : « sé-li-bré-cheune », « sé-ri-mo-ni », « né-beur », « o-zo-ri-tiz ».
\end{itemize}
\end{aretenir}

\begin{exercice}[Your turn]
Décrivez en cinq phrases une fête de votre quartier ou de votre village, en utilisant au moins six mots de cette section (par exemple : \eng{festival}, \eng{chief}, \eng{guests}, \eng{perform}, \eng{share a meal}, \eng{offer gifts}). Soulignez chaque mot de vocabulaire employé.
\end{exercice}

\coursec{Interacting with different people: register and politeness}

\textbf{Transition.} Nous savons maintenant \emph{nommer} les fêtes, les personnes et les activités. Mais parler d'une fête, ce n'est pas seulement employer les bons noms : c'est aussi \emph{s'adresser correctement} à son interlocuteur. On ne parle pas à son meilleur ami comme au principal du lycée ! C'est la question du \cle{register} (le registre de langue).

\courssub{Who is who: mates, neighbours, adults, authorities}

\begin{definition}[Les interlocuteurs]
\begin{itemize}
\item \cle{mate} : camarade, copain (registre familier, surtout britannique). \eng{My mates and I are organising the games.}
\item \cle{neighbour} : voisin(e). \eng{Our neighbour Ma Grace cooks for the whole street.}
\item \cle{adult} : adulte. \eng{Children should greet adults respectfully.}
\item \cle{elder} : aîné, ancien (personne âgée respectée dans la communauté). \eng{The elders sit under the big tent.}
\item \cle{school authorities} : les autorités scolaires (le principal, les censeurs, les enseignants). \eng{We asked the school authorities for permission.}
\item \cle{principal} (GB) / \cle{headmaster/headmistress} : principal(e) du lycée. \eng{The principal opened the ceremony.}
\end{itemize}
\end{definition}

\begin{attention}[« Principal » : faux ami partiel]
\eng{principal} (nom) désigne bien le \emph{principal} d'un établissement scolaire. Mais l'adjectif \eng{principal} signifie « principal » au sens de « le plus important » : \eng{the principal reason} = « la raison principale ». Ne confondez pas avec \eng{principle} (« principe », règle morale), qui se prononce exactement pareil !
\end{attention}

\courssub{Formal and informal register}

\begin{propriete}[Adapter son langage à l'interlocuteur]
En anglais, le choix des mots dépend de la \emph{relation} avec l'interlocuteur :
\begin{itemize}
\item \textbf{Registre familier} (\emph{informal}) : avec les \eng{mates} et les amis. Phrases courtes, contractions (\eng{Hi!}, \eng{What's up?}, \eng{Come and join us!}).
\item \textbf{Registre neutre} : avec les \eng{neighbours} et les connaissances. \eng{Hello, Ma Grace. How are you today?}
\item \textbf{Registre soutenu} (\emph{formal}) : avec les \eng{adults}, les \eng{elders} et les \eng{school authorities}. Phrases complètes, formules de politesse, pas de contractions. \eng{Good morning, Sir. May I ask you a question?}
\end{itemize}
\end{propriete}

\begin{center}
\begin{tabularx}{\linewidth}{|X|X|X|}
\hline
\rowcolor{popblueL} \textbf{Situation} & \textbf{Informal (mates)} & \textbf{Formal (adults, authorities)} \\
\hline
Saluer & Hi! / Hello, guys! & Good morning, Sir. / Good afternoon, Madam. \\
\hline
Inviter & Come to the party on Saturday! & We would be honoured if you could attend our celebration. \\
\hline
Demander la permission & Can I use the field? & May I ask for permission to use the school field, Sir? \\
\hline
Remercier & Thanks a lot! & Thank you very much for your kindness. \\
\hline
Prendre congé & See you! / Bye! & Goodbye, Sir. Have a nice day. \\
\hline
\end{tabularx}
\end{center}

\begin{definition}[Formules de politesse utiles]
\begin{itemize}
\item \cle{May I ...?} : « Puis-je ... ? » (demande de permission polie). \eng{May I come in, Sir?}
\item \cle{Could you ...?} : « Pourriez-vous ... ? ». \eng{Could you help us decorate the hall, please?}
\item \cle{I would like to ...} : « Je voudrais ... ». \eng{I would like to invite you to our festival.}
\item \cle{Excuse me} : « Excusez-moi » (pour attirer l'attention). \eng{Excuse me, Madam, where is the ceremony taking place?}
\item \cle{You are welcome} : « Je vous en prie » (réponse à \eng{thank you}). \eng{— Thank you for the gift. — You are welcome.}
\end{itemize}
\end{definition}

\begin{attention}[Les contractions et le registre soutenu]
Dans un échange formel (avec le principal, un chef, un ancien), évitez les contractions : dites \eng{I would like} et non \eng{I'd like}, \eng{we cannot} et non \eng{we can't}. À l'écrit formel (lettre d'invitation au principal), les contractions sont à proscrire. À l'oral entre \eng{mates}, elles sont naturelles et recommandées.
\end{attention}

\begin{exemplebox}[Trois invitations, trois registres]
\begin{itemize}
\item À un camarade : \eng{Hey, Peter! We're having a party at my place on Saturday. Come and have fun!} « Hé, Peter ! On fait une fête chez moi samedi. Viens t'amuser ! »
\item À une voisine : \eng{Hello, Ma Grace. We are organising a small celebration on Saturday. Would you like to join us?} « Bonjour, Ma Grace. Nous organisons une petite fête samedi. Voudriez-vous vous joindre à nous ? »
\item Au principal : \eng{Good morning, Sir. The students of Form Five would like to invite you to our end-of-year celebration. We would be honoured by your presence.} « Bonjour, Monsieur. Les élèves de Première voudraient vous inviter à leur fête de fin d'année. Nous serions honorés de votre présence. »
\end{itemize}
\end{exemplebox}

\begin{experience}[Role play: inviting different people]
Par groupes de trois, jouez la scène suivante : l'élève A invite l'élève B (son \eng{mate}) au festival du quartier ; puis A et B invitent ensemble l'élève C, qui joue le rôle du \eng{principal}, à être le \eng{guest of honour}. Utilisez au moins deux formules de politesse (\eng{May I ...?}, \eng{We would be honoured if ...}) et adaptez le registre à chaque interlocuteur. La classe écoute et relève les formules employées.
\end{experience}

\coursec{Collocations and idiomatic expressions}

\textbf{Transition.} Un mot ne vit jamais seul : il s'associe à d'autres mots selon des habitudes fixes de la langue. Ces associations s'appellent des \cle{collocations}. Les connaître, c'est parler un anglais naturel — et éviter les calques du français.

\courssub{Essential collocations about celebrations}

\begin{propriete}[Collocations à mémoriser]
\begin{itemize}
\item \cle{to hold a ceremony / a festival} : tenir, organiser une cérémonie. \eng{The festival is held every year in December.} (et non « make a ceremony »)
\item \cle{to take place} : avoir lieu. \eng{The celebration takes place on the school field.} (et non « take place in » sans complément de lieu précis : \eng{take place} + complément de lieu ou de temps)
\item \cle{to attend a ceremony / a party} : assister à. \eng{Fifty guests attended the wedding.}
\item \cle{to give / deliver a speech} : prononcer un discours. \eng{The chief delivered a long speech.}
\item \cle{to throw / have a party} : organiser une fête (familier). \eng{We are throwing a party for her birthday.}
\item \cle{to celebrate an anniversary / a victory} : fêter. \eng{The school celebrated its twentieth anniversary.}
\item \cle{to wish someone a happy birthday} : souhaiter un joyeux anniversaire. \eng{We wished her a happy birthday.}
\item \cle{to have fun / to enjoy oneself} : s'amuser. \eng{Everybody enjoyed themselves at the festival.}
\item \cle{to pay for} : payer (quelque chose). \eng{The sponsor paid for the drinks.} (jamais « pay the drinks »)
\item \cle{to look forward to + V-ing} : attendre avec impatience. \eng{We are looking forward to seeing you at the festival.}
\end{itemize}
\end{propriete}

\begin{attention}[Collocations : les erreurs typiques des francophones]
\begin{itemize}
\item « assister à la fête » : \eng{to attend the party}, et non « assist to the party ».
\item « avoir lieu » : \eng{to take place}, et non « to have place ».
\item « prononcer un discours » : \eng{to give a speech}, et non « to pronounce a speech » (\eng{pronounce} = prononcer un mot).
\item « fêter son anniversaire » : \eng{to celebrate one's birthday}, et non « to feast one's birthday ».
\item « s'amuser » : \eng{to have fun} ou \eng{to enjoy oneself}, et non « to amuse oneself » (qui signifie « se distraire en trompant l'ennui »).
\end{itemize}
\end{attention}

\courssub{Idiomatic expressions}

\begin{definition}[Expressions idiomatiques de la fête]
\begin{itemize}
\item \cle{the more, the merrier} : « plus on est de fous, plus on rit ». \eng{Can I bring my cousin? — Sure, the more, the merrier!}
\item \cle{to have a great time} : passer un excellent moment. \eng{We had a great time at the harvest festival.}
\item \cle{to dance the night away} : danser toute la nuit. \eng{The young people danced the night away.}
\item \cle{to be the life of the party} : être le boute-en-train, l'âme de la fête. \eng{My uncle is always the life of the party.}
\item \cle{in high spirits} : de très bonne humeur, joyeux. \eng{The whole village was in high spirits during the festival.}
\item \cle{a memorable occasion} : une occasion mémorable. \eng{The chief's visit made it a memorable occasion.}
\end{itemize}
\end{definition}

\begin{exoresolu}[Choose the right collocation]
Énoncé. Choisissez le verbe correct entre parenthèses.

\begin{enumerate}
\item The community (makes / holds) a festival every December.
\item The ceremony (takes / has) place in the school hall.
\item The principal (gave / pronounced) a speech to open the celebration.
\item We are (throwing / launching) a party for our teacher's birthday.
\item I look forward to (see / seeing) you at the festival.
\end{enumerate}

\tcblower
Solution détaillée.

\begin{enumerate}
\item \eng{holds} — collocation : \eng{to hold a festival}. « Make a festival » est un calque du français « faire une fête ».
\item \eng{takes} — \eng{to take place} = « avoir lieu ». « Has place » n'existe pas.
\item \eng{gave} — \eng{to give a speech}. \eng{Pronounce} signifie « prononcer » un mot, pas un discours.
\item \eng{throwing} — \eng{to throw a party} (familier) = « organiser une fête ». \eng{Launch} s'emploie pour un produit ou un projet.
\item \eng{seeing} — après \eng{to look forward to}, le verbe est en \emph{-ing}, car \eng{to} est ici une préposition, pas la marque de l'infinitif.
\end{enumerate}
\end{exoresolu}

\coursec{False friends and common traps for francophones}

\textbf{Transition.} Certains mots anglais ressemblent à des mots français mais ont un sens différent : ce sont les \cle{false friends} (faux amis). En voici les plus dangereux dans le champ lexical des fêtes et des relations sociales.

\begin{propriete}[Tableau des faux amis]
\begin{center}
\begin{tabularx}{\linewidth}{|X|X|X|}
\hline
\rowcolor{popblueL} \textbf{Mot anglais} & \textbf{Sens réel} & \textbf{Piège français} \\
\hline
to assist (someone) & aider quelqu'un & « assister à » = \eng{to attend} \\
\hline
a celebration & une fête, une célébration & « célébrité » = \eng{a celebrity} \\
\hline
actually & en fait, réellement & « actuellement » = \eng{currently, at present} \\
\hline
an event & un événement & « éventuel » = \eng{possible} ; « éventuellement » = \eng{possibly} \\
\hline
to demand & exiger & « demander » = \eng{to ask (for)} \\
\hline
sensible & raisonnable, sensé & « sensible » = \eng{sensitive} \\
\hline
a lecture & un cours magistral, une conférence & « lecture » (action de lire) = \eng{reading} \\
\hline
to attend & assister à, être présent à & « attendre » = \eng{to wait (for)} \\
\hline
a host & celui qui reçoit & « hôte » au sens d'invité = \eng{a guest} \\
\hline
to achieve & accomplir, réaliser & « achever » au sens de finir = \eng{to finish, to complete} \\
\hline
\end{tabularx}
\end{center}
\end{propriete}

\begin{attention}[Ask, demand, require : trois verbes à ne pas confondre]
\begin{itemize}
\item \eng{to ask} = demander (poliment) : \eng{We asked the principal for permission.} « Nous avons demandé la permission au principal. »
\item \eng{to demand} = exiger (avec autorité, souvent agressif) : \eng{The workers demanded higher wages.} « Les ouvriers ont exigé des salaires plus élevés. »
\item \eng{to require} = nécessiter, requérir (formel) : \eng{The ceremony requires traditional costumes.} « La cérémonie requiert des costumes traditionnels. »
\end{itemize}
Ne dites jamais « I demand permission to speak » au principal : vous seriez très impoli ! Dites \eng{May I ask for permission to speak, Sir?}
\end{attention}

\begin{attention}[Prépositions piégées]
\begin{itemize}
\item \eng{to wait \textbf{for} someone} : attendre quelqu'un. \eng{We waited for the chief for an hour.}
\item \eng{to listen \textbf{to} the speech} : écouter le discours. \eng{Everybody listened to the principal's speech.}
\item \eng{to invite someone \textbf{to} a party} : inviter quelqu'un à une fête. \eng{She invited me to her birthday party.}
\item \eng{to congratulate someone \textbf{on} something} : féliciter quelqu'un de quelque chose. \eng{We congratulated the organisers on their success.}
\item \eng{to depend \textbf{on}} : dépendre de. \eng{The festival depends on the weather.}
\end{itemize}
\end{attention}

\coursec{Word families, pronunciation and spelling}

\textbf{Transition.} Pour enrichir rapidement votre vocabulaire, apprenez les mots par \emph{familles} : à partir d'un verbe, formez le nom, l'adjectif, la personne. C'est aussi une compétence exigée à l'examen (exercices de \emph{word formation}).

\courssub{Word families}

\begin{propriete}[Familles de mots du thème]
\begin{center}
\begin{tabularx}{\linewidth}{|l|l|l|X|}
\hline
\rowcolor{popblueL} \textbf{Verbe} & \textbf{Nom (action)} & \textbf{Personne} & \textbf{Adjectif} \\
\hline
to celebrate & celebration & celebrant & celebrated, celebratory \\
\hline
to organise & organisation & organiser & organised \\
\hline
to invite & invitation & invitee, guest & invited \\
\hline
to decorate & decoration & decorator & decorative \\
\hline
to perform & performance & performer & performing \\
\hline
to participate & participation & participant & participating \\
\hline
to entertain & entertainment & entertainer & entertaining \\
\hline
to gather & gathering & — & gathered \\
\hline
\end{tabularx}
\end{center}
\end{propriete}

\begin{exemplebox}[Une famille en action : celebrate]
\begin{itemize}
\item Verbe : \eng{We celebrate Christmas with the whole family.} « Nous fêtons Noël avec toute la famille. »
\item Nom : \eng{The celebration lasted three days.} « La fête a duré trois jours. »
\item Adjectif : \eng{It was a celebrated event in the village.} « C'était un événement célèbre dans le village. »
\end{itemize}
\end{exemplebox}

\begin{attention}[Orthographe : les suffixes qui changent tout]
\begin{itemize}
\item \eng{celebrate} $\rightarrow$ \eng{celebr\textbf{a}tion} (et non « celebretion ») ; de même \eng{invite} $\rightarrow$ \eng{invit\textbf{a}tion}, \eng{decorate} $\rightarrow$ \eng{decor\textbf{a}tion}, \eng{organise} $\rightarrow$ \eng{organis\textbf{a}tion}.
\item Le \emph{e} final muet disparaît devant \emph{-ation} et \emph{-ing} : \eng{celebrate} $\rightarrow$ \eng{celebrating}, \eng{decorate} $\rightarrow$ \eng{decorating}.
\item \eng{neighbour} garde le \emph{u} en anglais britannique (\eng{neighbor} en américain) ; de même \eng{colour}, \eng{honour}, \eng{favourite}.
\end{itemize}
\end{attention}

\courssub{Pronunciation: word stress in word families}

\begin{propriete}[L'accent tonique se déplace]
En anglais, l'accent tonique change souvent de place dans une famille de mots, ce qui modifie la prononciation :
\begin{itemize}
\item \eng{CElebrate} « sé-li-breit » $\rightarrow$ \eng{celeBRAtion} « sé-li-bré-cheune »
\item \eng{PHOtograph} « fo-to-graf » $\rightarrow$ \eng{phoTOgraphy} « feu-to-gra-fi »
\item \eng{INvite} (verbe) « in-vaït » $\rightarrow$ \eng{inviTAtion} « in-vi-té-cheune »
\item \eng{PARticipate} « par-ti-si-peit » $\rightarrow$ \eng{particiPAtion} « par-ti-si-pé-cheune »
\end{itemize}
Règle pratique : dans les noms en \emph{-tion}, l'accent tonique tombe toujours sur la syllabe qui précède \emph{-tion}.
\end{propriete}

\begin{exoresolu}[Word formation]
Énoncé. Complétez chaque phrase avec la forme correcte du mot entre parenthèses.

\begin{enumerate}
\item The \dots\ of the hall took the pupils two hours. (decorate)
\item All the \dots\ received a certificate at the end of the cultural week. (participate)
\item The \dots\ of the festival thanked everybody for coming. (organise)
\item We sent an \dots\ card to the principal. (invite)
\item The dancers gave a wonderful \dots\ in front of the chief. (perform)
\end{enumerate}

\tcblower
Solution détaillée.

\begin{enumerate}
\item \eng{decoration} — nom de l'action ; \eng{the decoration of the hall} « la décoration de la salle ».
\item \eng{participants} — nom de personne au pluriel (\eng{all} + pluriel) : « tous les participants ».
\item \eng{organisers} — nom de personne ; le pluriel est naturel ici (un comité), mais \eng{organiser} au singulier est aussi acceptable si le contexte l'impose.
\item \eng{invitation} — nom : \eng{an invitation card} « une carte d'invitation ».
\item \eng{performance} — nom de l'action : \eng{to give a performance} « donner un spectacle ».
\end{enumerate}

Piège à éviter : à la question 2, « participations » est un calque du français ; la personne qui participe est \eng{a participant}.
\end{exoresolu}

\coursec{Using the vocabulary in context}

\textbf{Transition.} Dernière étape : réemployer tout ce lexique dans des situations réelles de communication — un dialogue, une lettre d'invitation, un court paragraphe. C'est exactement ce que l'on vous demandera à l'examen (compréhension et composition).

\courssub{A model dialogue}

\begin{textebox}[Asking the principal for permission]
Student: Good morning, Sir. May I come in?

Principal: Good morning, Etienne. Yes, come in. What can I do for you?

Student: Sir, the students of Form Five are organising the end-of-year celebration. We would like to hold the ceremony on the school field on Friday, the twentieth of December.

Principal: That sounds interesting. Who is sponsoring the event?

Student: A local businessman has agreed to pay for the tents and the sound system, Sir. Our parents will prepare the food.

Principal: Very good. And who are your guests of honour?

Student: The chief of our neighbourhood and the elders, Sir. And we would be honoured if you could give the opening speech.

Principal: I accept with pleasure. You have my permission to use the field. Well done, and good luck with the organisation!

Student: Thank you very much, Sir. We are looking forward to seeing you at the celebration.
\end{textebox}

\textbf{Glossaire.} \eng{May I come in?} « Puis-je entrer ? » ; \eng{to hold the ceremony} « tenir la cérémonie » ; \eng{guests of honour} « invités d'honneur » ; \eng{I accept with pleasure} « j'accepte avec plaisir » ; \eng{well done} « bravo ».

\courssub{A model letter of invitation}

\begin{textebox}[Inviting the guest of honour]
Government High School, Bamenda

10th December 2024

Dear Sir,

The students of Form Five are organising the end-of-year celebration of our school. The ceremony will take place on the school field on Friday, 20th December, at nine o'clock in the morning.

We would be deeply honoured if you could attend the celebration as our guest of honour and give the opening speech. Traditional dances, songs and a shared meal are on the programme.

We look forward to hearing from you.

Yours faithfully,

Etienne Ndi

Class prefect, Form Five
\end{textebox}

\begin{methode}[Rédiger une lettre d'invitation formelle]
\begin{enumerate}
\item \textbf{En-tête} : nom de l'école ou adresse, puis la date en anglais (\eng{10th December 2024}).
\item \textbf{Formule d'appel} : \eng{Dear Sir,} / \eng{Dear Madam,} (virgule, pas deux-points).
\item \textbf{Paragraphe 1} : présenter l'événement (\eng{are organising}, \eng{will take place}).
\item \textbf{Paragraphe 2} : formuler l'invitation avec une formule de politesse (\eng{We would be honoured if you could ...}) et donner le programme.
\item \textbf{Paragraphe 3} : formule d'attente (\eng{We look forward to hearing from you.}).
\item \textbf{Formule de politesse finale} : \eng{Yours faithfully,} (quand on ne connaît pas le nom du destinataire), puis signature et qualité.
\end{enumerate}
\end{methode}

\begin{attention}[Erreurs fréquentes dans la lettre d'invitation]
\begin{itemize}
\item Pas de contractions dans une lettre formelle : écrivez \eng{we would}, \eng{we are}, pas « we'd », « we're ».
\item \eng{We look forward to hearing from you} : \eng{hearing} avec \emph{-ing} (préposition \eng{to}), jamais « to hear ».
\item La date anglaise : \eng{20th December} ou \eng{December 20th} ; ne mélangez pas les formats.
\item \eng{Yours faithfully} s'emploie avec \eng{Dear Sir/Madam} ; \eng{Yours sincerely} s'emploie quand on connaît le nom du destinataire (\eng{Dear Mr Tabi}).
\end{itemize}
\end{attention}

\begin{exoresolu}[Find and correct the mistakes]
Énoncé. Chaque phrase contient une erreur de vocabulaire ou de collocation. Corrigez-la.

\begin{enumerate}
\item We assisted to the harvest festival last Saturday.
\item The principal pronounced a long speech at the ceremony.
\item I demand you to attend our celebration, Sir.
\item We look forward to see you at the party.
\item The sponsor paid the drinks and the tents.
\end{enumerate}

\tcblower
Solution détaillée.

\begin{enumerate}
\item \eng{We \textbf{attended} the harvest festival last Saturday.} — \eng{attend} (sans \eng{to}) = « assister à » ; \eng{assist} signifie « aider ».
\item \eng{The principal \textbf{gave} a long speech at the ceremony.} — collocation : \eng{to give a speech}.
\item \eng{I \textbf{would like to invite} you to attend our celebration, Sir.} — \eng{demand} = « exiger », beaucoup trop brutal pour s'adresser à une autorité.
\item \eng{We look forward to \textbf{seeing} you at the party.} — après \eng{look forward to}, verbe en \emph{-ing}.
\item \eng{The sponsor paid \textbf{for} the drinks and the tents.} — \eng{to pay for something} ; on paie \emph{pour} quelque chose.
\end{enumerate}
\end{exoresolu}

\begin{exercice}[Guided writing]
Rédigez une courte lettre (10 à 12 lignes) au principal de votre lycée pour l'inviter, en tant que \eng{guest of honour}, à la fête culturelle de votre classe. Utilisez au moins cinq éléments du lexique de la leçon (\eng{celebration}, \eng{take place}, \eng{guest of honour}, \eng{give a speech}, \eng{look forward to}, \eng{attend} ...) et respectez la structure de la lettre formelle étudiée dans la méthode.
\end{exercice}

\begin{exercice}[Speaking practice]
Par binômes : l'élève A est un élève qui organise la fête de fin d'année du quartier ; l'élève B est un \eng{elder} du quartier. A invite B, lui explique le programme (danses, repas partagé, discours du chef) et répond à ses questions. Puis inversez les rôles avec un nouvel événement : un mariage ou une fête d'anniversaire.
\end{exercice}

\begin{aretenir}[L'essentiel de la leçon]
\begin{itemize}
\item \textbf{Lexique} : \eng{festival, ceremony, celebration, gathering, anniversary} ; \eng{host, guest, organiser, chief, elder, guest of honour} ; \eng{organise, plan, schedule, host, sponsor, invite} ; \eng{perform, parade, share a meal, give a speech, offer gifts}.
\item \textbf{Registre} : adaptez votre langue à l'interlocuteur — familier avec les \eng{mates}, soutenu avec les \eng{adults} et les \eng{school authorities} (\eng{May I ...?}, \eng{We would be honoured if ...}).
\item \textbf{Collocations} : \eng{hold a ceremony}, \eng{take place}, \eng{attend a party}, \eng{give a speech}, \eng{throw a party}, \eng{look forward to + V-ing}.
\item \textbf{Faux amis} : \eng{assist} (aider) $\neq$ \eng{attend} (assister à) ; \eng{demand} (exiger) $\neq$ \eng{ask} (demander) ; \eng{actually} (en fait) $\neq$ \eng{currently} (actuellement).
\item \textbf{Familles de mots} : \eng{celebrate / celebration}, \eng{invite / invitation}, \eng{participate / participant} ; l'accent tonique précède toujours le suffixe \emph{-tion}.
\end{itemize}
\end{aretenir}

\coursec{Interacting with different people: register and politeness}

Dans la section précédente, nous avons découvert le vocabulaire essentiel des fêtes et célébrations communautaires : \eng{celebration}, \eng{festival}, \eng{feast}, \eng{parade}, \eng{ceremony}\ldots Mais savoir nommer les choses ne suffit pas. Lors d'une fête de quartier à Yaoundé ou d'une \eng{school prize-giving day} à Buea, on ne parle pas de la même manière à son meilleur ami, à la voisine, au chef du village ou au proviseur. En anglais, cette adaptation du langage à l'interlocuteur s'appelle le \cle{register} (le registre de langue). C'est l'objet de cette section : apprendre à choisir les bons mots selon la personne à qui l'on s'adresse.

\courssub{Observing the language in context}

Commençons par observer quatre situations de communication authentiques. Lisez chaque dialogue et demandez-vous : qui parle à qui ? Quels mots changent d'une situation à l'autre ?

\begin{textebox}[Four dialogues at the Nkol-Eton community festival]
\textbf{Dialogue 1 — two mates:}\\
— “Hey Paul! What's up? Are you coming to the party tonight?”\\
— “Sure! It's gonna be fun. Come and join us at six, don't be late!”\\[4pt]
\textbf{Dialogue 2 — two neighbours:}\\
— “Good morning, Ma Grace. We're having a small celebration at our house on Saturday. Would you like to join us?”\\
— “That's very kind of you. Can I help you with the cooking?”\\
— “Oh, that would be lovely, thank you!”\\[4pt]
\textbf{Dialogue 3 — a student and an elder:}\\
— “Good afternoon, Sir. We would be honoured by your presence at our end-of-year festival. May I offer you a seat?”\\
— “Thank you, young man. I am pleased to be here.”\\[4pt]
\textbf{Dialogue 4 — a student and the principal:}\\
— “Good morning, Sir. I am writing to request permission to use the school hall for our cultural day. We would be grateful if you could attend.”\\
— “Very well. Bring me a formal letter by Friday.”
\end{textebox}

\begin{center}
\begin{tabularx}{\linewidth}{|X|X|X|}\hline
\rowcolor{popblueL} \textbf{English} & \textbf{Français} & \textbf{Observation} \\ \hline
What's up? & Quoi de neuf ? / Ça va ? & très familier, entre amis \\ \hline
It's gonna be fun. & Ça va être génial. & \eng{gonna} = \eng{going to} oral \\ \hline
Would you like to join us? & Voudriez-vous vous joindre à nous ? & poli, pour une invitation \\ \hline
We would be honoured by your presence. & Nous serions honorés de votre présence. & très formel \\ \hline
May I offer you a seat? & Puis-je vous offrir une place assise ? & respectueux, envers un ancien \\ \hline
I am writing to request permission. & Je vous écris pour demander l'autorisation. & style de lettre officielle \\ \hline
\end{tabularx}
\end{center}

On observe que pour la \emph{même action} (inviter quelqu'un à une fête), l'anglais utilise des mots et des structures complètement différents selon le destinataire. C'est exactement ce que nous allons systématiser.

\courssub{Talking to mates: informal language}

Avec ses camarades de classe ou ses amis du quartier, on utilise un langage \cle{informal} (familier) : salutations courtes, contractions, expressions de la vie quotidienne.

\begin{definition}[Informal greetings and invitations]
Pour saluer et inviter un camarade :
\begin{itemize}
\item \eng{Hi!} / \eng{Hello!} — Salut !
\item \eng{Hey!} — Hé ! (très familier, pour attirer l'attention)
\item \eng{What's up?} — Quoi de neuf ? (réponse typique : \eng{Not much!})
\item \eng{How's it going?} — Comment ça va ?
\item \eng{Come and join us!} — Viens nous rejoindre !
\item \eng{Are you coming to the party?} — Tu viens à la fête ?
\item \eng{It's gonna be fun!} — Ça va être génial ! (\eng{gonna} est la forme orale de \eng{going to} ; à l'écrit formel, on écrit toujours \eng{going to})
\item \eng{See you there!} — On se voit là-bas !
\end{itemize}
\end{definition}

\begin{exemplebox}[Talking to a mate]
\eng{“Hey Paul, are you coming to the party?”}\\
« Hé Paul, tu viens à la fête ? »\\[4pt]
\eng{“Hi Sandra! We're meeting at the stadium at four. Don't miss it!”}\\
« Salut Sandra ! On se retrouve au stade à quatre heures. Ne le rate pas ! »\\[4pt]
\eng{“What's up, man? Bring your drums to the festival!”}\\
« Quoi de neuf, mon pote ? Apporte tes tambours au festival ! »
\end{exemplebox}

\courssub{Talking to neighbours: friendly but respectful}

Avec les voisins, le ton est \cle{neutral} : chaleureux mais respectueux. On évite le langage trop familier, mais on reste simple et naturel. Au Cameroun anglophone, on s'adresse souvent à une voisine âgée en disant \eng{Ma} + prénom (\eng{Ma Grace}), ou \eng{Pa} + prénom pour un homme âgé : c'est une marque d'affection respectueuse.

\begin{exemplebox}[Talking to a neighbour]
\eng{“Good morning, Ma Grace. Would you like to join us for the celebration?”}\\
« Bonjour, Ma Grace. Aimeriez-vous vous joindre à nous pour la fête ? »\\[4pt]
\eng{“Can I help you with the cooking?”}\\
« Puis-je vous aider avec la cuisine ? »\\[4pt]
\eng{“Thank you for the plantains, Pa Njoh. They were delicious.”}\\
« Merci pour les plantains, Pa Njoh. Ils étaient délicieux. »
\end{exemplebox}

Remarquez l'emploi de \eng{“Would you like\ldots?”} et \eng{“Can I help you\ldots?”} : ce sont des formules \emph{polies} mais pas cérémonieuses, parfaites entre voisins.

\courssub{Talking to adults and elders: polite language}

Envers les adultes et les anciens (\eng{elders}), la politesse devient obligatoire. On utilise \eng{Sir} / \eng{Madam}, les auxiliaires de politesse \eng{would}, \eng{could}, \eng{may}, et des tournures plus longues.

\begin{propriete}[Les auxiliaires de la politesse]
\begin{itemize}
\item \eng{“May I\ldots?”} = Puis-je\ldots ? (demande de permission très polie) : \eng{“May I offer you a seat?”} « Puis-je vous offrir une place ? »
\item \eng{“Could you\ldots?”} = Pourriez-vous\ldots ? : \eng{“Could you open the ceremony, please?”} « Pourriez-vous ouvrir la cérémonie, s'il vous plaît ? »
\item \eng{“We would be honoured\ldots”} = Nous serions honorés\ldots : \eng{“We would be honoured by your presence.”} « Nous serions honorés de votre présence. »
\item \eng{“Would you mind\ldots?”} = Cela vous dérangerait-il de\ldots ? : \eng{“Would you mind saying a few words?”} « Auriez-vous l'amabilité de dire quelques mots ? »
\end{itemize}
\end{propriete}

\begin{exemplebox}[Talking to an elder]
\eng{“Good afternoon, Sir. We would be honoured by your presence at our festival.”}\\
« Bonjour, Monsieur. Nous serions honorés de votre présence à notre festival. »\\[4pt]
\eng{“May I offer you a seat, Madam?”}\\
« Puis-je vous offrir une place, Madame ? »\\[4pt]
\eng{“Please, Sir, could you bless the ceremony?”}\\
« S'il vous plaît, Monsieur, pourriez-vous bénir la cérémonie ? »
\end{exemplebox}

\courssub{Talking to school authorities: formal language}

Avec le proviseur (\eng{principal}, \eng{headmaster}), le principal ou les enseignants, on emploie le registre \cle{formal}, surtout à l'écrit (lettres officielles, demandes d'autorisation). Ce registre a ses propres formules figées.

\begin{definition}[Formal expressions for school authorities]
\begin{itemize}
\item \eng{“Dear Sir,”} / \eng{“Dear Madam,”} — Monsieur, / Madame, (en-tête de lettre)
\item \eng{“I am writing to request permission to\ldots”} — Je vous écris pour demander l'autorisation de\ldots
\item \eng{“We would be grateful if you could attend.”} — Nous vous serions reconnaissants de bien vouloir y assister.
\item \eng{“I would like to invite you to\ldots”} — J'aimerais vous inviter à\ldots
\item \eng{“Thank you for your kind attention.”} — Je vous remercie de votre aimable attention.
\item \eng{“Yours faithfully,”} — formule de politesse finale (quand on ne connaît pas le nom du destinataire)
\end{itemize}
\end{definition}

\begin{exemplebox}[A formal request]
\eng{“Dear Sir, I am writing to request permission to use the school hall for our cultural day on 15th May. We would be grateful if you could attend. Yours faithfully, Amina Bello.”}\\
« Monsieur, je vous écris pour demander l'autorisation d'utiliser la salle des fêtes de l'école pour notre journée culturelle du 15 mai. Nous vous serions reconnaissants de bien vouloir y assister. Veuillez agréer, Monsieur, l'expression de mes salutations distinguées. Amina Bello. »\\[6pt]
Comparez avec la version informelle :\\[4pt]
\eng{“Dear Sir, we would be delighted if you could attend our celebration.”}\\
« Cher Monsieur, nous serions ravis que vous assistiez à notre célébration. »
\end{exemplebox}

\courssub{The formality scale: a visual summary}

\begin{popfigure}
\begin{center}
\begin{tikzpicture}[scale=1]
\draw[-{Stealth[length=3mm]}, very thick, popline] (0,0) -- (12,0);
\node[below, popmuted] at (0.2,-0.15) {\small moins formel};
\node[below, popmuted] at (11.8,-0.15) {\small plus formel};
\node[draw, rounded corners, fill=poppinkL, align=center, minimum width=3.2cm, minimum height=2.4cm] at (2,1.9) {\textbf{\small INFORMAL}\\[2pt] \small mates, friends\\ \emph{\small “Hey! What's up?”}\\ \emph{\small “Come and join us!”}};
\node[draw, rounded corners, fill=popgoldL, align=center, minimum width=3.2cm, minimum height=2.4cm] at (6,1.9) {\textbf{\small NEUTRAL}\\[2pt] \small neighbours\\ \emph{\small “Good morning, Ma Grace.”}\\ \emph{\small “Would you like to join us?”}};
\node[draw, rounded corners, fill=popblueL, align=center, minimum width=3.4cm, minimum height=2.4cm] at (10,1.9) {\textbf{\small FORMAL}\\[2pt] \small elders, authorities\\ \emph{\small “We would be honoured\ldots”}\\ \emph{\small “May I offer you a seat?”}};
\fill[poppink] (2,0) circle (3pt);
\fill[popgold] (6,0) circle (3pt);
\fill[popblue] (10,0) circle (3pt);
\end{tikzpicture}
\end{center}
\legende{L'échelle de la formalité : plus l'interlocuteur est éloigné ou hiérarchiquement supérieur, plus le langage doit être formel.}
\end{popfigure}

\begin{propriete}[Le choix du vocabulaire dépend du destinataire]
En anglais comme en français, le \cle{register} (registre de langue) dépend de trois facteurs : l'\textbf{âge} de l'interlocuteur, son \textbf{statut social} et le \textbf{degré de familiarité}. On dit \eng{“Hi mate!”} à un camarade, mais \eng{“Good morning, Madam”} à une autorité. Utiliser le mauvais registre peut paraître impoli ou ridicule : dire \eng{“Hey! What's up?”} au proviseur serait aussi choquant que de dire « Salut la miss ! » à la principale du lycée.
\end{propriete}

\courssub{Comparative table: one message, four registers}

Voici le tableau de référence de cette leçon. Il montre comment exprimer les mêmes fonctions (saluer, inviter, remercier, prendre congé) selon l'interlocuteur.

\begin{center}
\begin{tabularx}{\linewidth}{|l|X|X|X|X|}\hline
\rowcolor{popblueL} \textbf{Fonction} & \textbf{Mate (informal)} & \textbf{Neighbour (neutral)} & \textbf{Adult / elder (polite)} & \textbf{Authority (formal)} \\ \hline
Salutation & Hi! Hey! & Good morning, Ma Grace. & Good afternoon, Sir. & Dear Sir, / Dear Madam, \\ \hline
Invitation & Come and join us! & Would you like to join us? & We would be honoured by your presence. & We would be grateful if you could attend. \\ \hline
Demande & Can you bring the drums? & Can I help you with the cooking? & Could you open the ceremony, please? & I am writing to request permission to\ldots \\ \hline
Remerciement & Thanks a lot! & Thank you very much. & Thank you for your kindness. & We really appreciate your support. \\ \hline
Congé & See you! Bye! & Goodbye, have a nice day. & Goodbye, Sir. It was a pleasure. & Yours faithfully, / Thank you for your attention. \\ \hline
\end{tabularx}
\end{center}

\begin{aretenir}[Formules de félicitation et de remerciement]
Quel que soit le registre, ces formules sont indispensables pendant une fête :
\begin{itemize}
\item \eng{Congratulations!} — Félicitations ! (à un lauréat, aux mariés)
\item \eng{Well done!} — Bravo ! Bien joué !
\item \eng{Thank you for coming.} — Merci d'être venu(e).
\item \eng{We really appreciate your support.} — Nous apprécions vraiment votre soutien. (formel)
\item \eng{It was a great pleasure to have you with us.} — Ce fut un grand plaisir de vous avoir parmi nous.
\end{itemize}
\end{aretenir}

\courssub{How to adapt your speech: the method}

\begin{methode}[Adapter son discours à l'interlocuteur en quatre étapes]
\begin{enumerate}
\item \textbf{Identifier l'interlocuteur} : est-ce un ami, un voisin, un adulte, une autorité ? Quel est son âge et son statut ?
\item \textbf{Choisir la salutation adaptée} : \eng{“Hey!”} pour un ami, \eng{“Good morning, Ma\ldots”} pour un voisin, \eng{“Good afternoon, Sir.”} pour un ancien, \eng{“Dear Sir,”} dans une lettre officielle.
\item \textbf{Employer les auxiliaires de politesse} : remplacer \eng{can} par \eng{could}, \eng{will} par \eng{would}, et utiliser \eng{may} pour demander la permission. Plus la phrase est longue et indirecte, plus elle est polie.
\item \textbf{Terminer par une formule de courtoisie} : \eng{“Thank you for coming.”}, \eng{“We really appreciate your support.”}, ou \eng{“Yours faithfully,”} à l'écrit.
\end{enumerate}
\end{methode}

\begin{exoresolu}[Choisir le bon registre]
\textbf{Énoncé.} Pour chaque situation, choisissez la phrase la plus adaptée (a, b ou c).\\[3pt]
1. You invite your best friend to the school party.\\
a) \eng{“We would be honoured by your presence.”} \quad b) \eng{“Hey, are you coming to the party?”} \quad c) \eng{“Dear Sir, I am writing to invite you\ldots”}\\[3pt]
2. You ask the headmaster for the school hall.\\
a) \eng{“Give me the hall, please.”} \quad b) \eng{“What's up, Sir? Can I take the hall?”} \quad c) \eng{“I am writing to request permission to use the school hall.”}\\[3pt]
3. You offer a seat to an old neighbour at the festival.\\
a) \eng{“May I offer you a seat, Madam?”} \quad b) \eng{“Sit down here!”} \quad c) \eng{“Hey, take this chair!”}
\tcblower
\textbf{Solution détaillée.}\\
1. Réponse \textbf{b} : un meilleur ami appelle le registre \emph{informal}. La phrase a) est trop cérémonieuse (on la réserve à un ancien) ; la phrase c) est une formule de lettre officielle, absurde entre amis.\\[3pt]
2. Réponse \textbf{c} : s'adresser au proviseur exige le registre \emph{formal}, avec la tournure \eng{“I am writing to request permission to\ldots”}. La phrase a) est un ordre direct (impoli) ; la phrase b) mélange \eng{“What's up”} (trop familier) et \eng{Sir} (formel) : registre incohérent.\\[3pt]
3. Réponse \textbf{a} : \eng{“May I offer you a seat?”} est la formule polie par excellence envers une personne âgée. Les phrases b) et c) sont des impératifs secs, inacceptables avec une voisine âgée.
\end{exoresolu}

\courssub{Common traps for francophone learners}

\begin{attention}[Ne traduisez jamais mot à mot !]
\textbf{Piège n\textsuperscript{o} 1 :} ne traduisez jamais « Tu peux venir ? » par \eng{“You can come?”} en vous adressant à un adulte ou à une autorité. Cette question directe sonne comme un ordre ou un manque de respect. Préférez \eng{“Would you like to come?”} ou \eng{“Could you join us?”}\\[4pt]
\textbf{Piège n\textsuperscript{o} 2 :} le tutoiement n'existe pas en anglais : \eng{you} sert pour tout le monde. La politesse ne passe donc pas par le pronom (comme « vous » en français) mais par les \emph{mots} : \eng{Sir}, \eng{Madam}, \eng{please}, \eng{would}, \eng{could}, \eng{may}.\\[4pt]
\textbf{Piège n\textsuperscript{o} 3 :} attention aux intonations montantes : \eng{“You can come?”} avec la voix qui monte reste familier. La forme interrogative correcte est \eng{“Can you come?”} (inversion), et sa version polie \eng{“Could you come?”}
\end{attention}

\begin{savaistu}[“Monsieur” tout seul n'existe pas en anglais !]
En français, on peut interpeller quelqu'un par « Monsieur ! » ou « Madame ! » tout court. En anglais, c'est impossible : on utilise \eng{Sir} ou \eng{Madam} (sans article, avec une majuscule en début de phrase), ou bien \eng{Mr} / \eng{Mrs} / \eng{Ms} \emph{suivis obligatoirement d'un nom de famille} : \eng{Mr Etoundi}, \eng{Mrs Ngono}. Dire \eng{Mister!} seul pour héler quelqu'un est très familier, voire impoli. Entre amis, on utilise simplement le prénom : \eng{“John, come here!”} Au Cameroun anglophone, les élèves disent respectueusement \eng{Sir} à un enseignant homme et \eng{Madam} ou \eng{Ma} à une enseignante — une habitude héritée de la tradition scolaire britannique.
\end{savaistu}

\begin{exercice}[Your turn: adapt the message]
Réécrivez le message « Viens à notre fête samedi ! » en anglais, en l'adaptant à chaque destinataire : 1) \eng{your best friend} ; 2) \eng{your neighbour Ma Grace} ; 3) \eng{the village chief} ; 4) \eng{your principal} (in a formal letter).
\end{exercice}

\begin{corrige}[Exercice « Your turn: adapt the message »]
1) \eng{“Hey! Come to our party on Saturday, it's gonna be fun!”}\\[2pt]
2) \eng{“Good morning, Ma Grace. Would you like to join us for our party on Saturday?”}\\[2pt]
3) \eng{“Good afternoon, Sir. We would be honoured by your presence at our celebration on Saturday.”}\\[2pt]
4) \eng{“Dear Sir, I am writing to invite you to our celebration on Saturday. We would be grateful if you could attend. Yours faithfully, \ldots”}
\end{corrige}

\begin{aretenir}[L'essentiel de la section]
\begin{itemize}
\item Le \cle{register} (registre) dépend de l'interlocuteur : \emph{informal} (mates), \emph{neutral} (neighbours), \emph{formal} (elders, authorities).
\item La politesse anglaise repose sur \eng{would}, \eng{could}, \eng{may}, \eng{please}, et non sur un pronom comme le « vous » français.
\item Salutations : \eng{“Hey!”} / \eng{“Good morning, Ma\ldots”} / \eng{“Good afternoon, Sir.”} / \eng{“Dear Sir,”}
\item Invitations : \eng{“Come and join us!”} / \eng{“Would you like to join us?”} / \eng{“We would be honoured by your presence.”} / \eng{“We would be grateful if you could attend.”}
\item Félicitations et remerciements : \eng{Congratulations!}, \eng{Well done!}, \eng{Thank you for coming.}, \eng{We really appreciate your support.}
\end{itemize}
\end{aretenir}

Maintenant que vous savez adapter votre langage à chaque interlocuteur, la section suivante vous montrera comment combiner ces mots naturellement grâce aux \emph{collocations} et aux expressions idiomatiques des fêtes communautaires.

\coursec{Collocations and idiomatic expressions}

Dans la section précédente, nous avons vu comment adapter notre langage (\emph{register}) selon que nous parlons à un camarade, un voisin, un adulte ou une autorité scolaire. Maîtriser la politesse ne suffit pas : il faut encore savoir \emph{quels mots vont ensemble} naturellement. En anglais, comme en français, les mots ne s’assemblent pas au hasard : on \eng{throws a party} (on organise une fête), mais on ne \eng{does a party} ; on \eng{marks an occasion} (on marque l’événement), mais on ne \eng{signs an occasion}. C’est l’objet des \textbf{collocations}. Parallèlement, les anglophones utilisent une multitude d’\textbf{expressions idiomatiques} (idioms) pour décrire la fête, l’ambiance et les relations sociales. Les connaître permet de passer d’un anglais « scolaire » à un anglais vivant, naturel et précis — celui qu’on entend dans les rues de Douala, de Yaoundé ou de Buea lors des fêtes de la jeunesse, du Ngondo ou de la rentrée scolaire.

\courssub{What is a collocation?}

\begin{definition}[Collocation]
Une \cle{collocation} est une association habituelle et prévisible de deux ou plusieurs mots qui vont ensemble « naturellement » dans la langue. Ce ne sont pas des règles grammaticales strictes, mais des habitudes lexicales : \eng{to attend a ceremony} (assister à une cérémonie), \eng{to have fun} (s’amuser), \eng{heavy rain} (pluie forte — pas \eng{thick rain}). En anglais, le verbe \textbf{détermine} souvent le nom qui suit (verb + noun), et inversement. Apprendre un mot isolé (\eng{party}) ne suffit pas ; il faut apprendre ses « partenaires » privilégiés (\eng{throw a party}, \eng{attend a party}, \eng{birthday party}, \eng{surprise party}).
\end{definition}

\begin{methode}[Comment mémoriser les collocations efficacement]
\begin{enumerate}
    \item \textbf{Apprenez le couple verbe + nom (verb + noun) en bloc}, jamais le verbe seul ni le nom seul. Écrivez : \eng{throw a party}, \eng{hold a celebration}, \eng{send an invitation}.
    \item \textbf{Contextualisez immédiatement} : écrivez une phrase personnelle vraie. Ex. : \eng{Last year, my family \textbf{held a celebration} for my grandmother's 80th birthday.}
    \item \textbf{Regroupez par thème} : fête (\eng{party/celebration}), événement (\eng{festival/event}), invitation (\eng{invitation}).
    \item \textbf{Révisez par « familles »} : tous les verbes qui vont avec \eng{party} (\eng{throw, hold, attend, enjoy, host, gatecrash}).
    \item \textbf{Attention aux faux amis structurels} : ne traduisez pas mot à mot depuis le français (\eng{do a party} $\leftarrow$ « faire une fête »).
\end{enumerate}
\end{methode}

\courssub{Core collocations: Party, Celebration, Festival, Event, Invitation}

Observons d’abord un court texte authentique : un message WhatsApp d’un élève de 1ère A de Buea à ses camarades pour l’anniversaire du lycée.

\begin{textebox}[WhatsApp group message — School Anniversary Planning]
Hey everyone! The Principal just announced we’ll \textbf{hold a celebration} for the school’s 25th anniversary next month. The Students’ Union will \textbf{organise a festival} of arts and culture on the esplanade. We need volunteers to \textbf{send invitations} to former students and local authorities. If you \textbf{accept the invitation} to help, meet me at the library at 3 pm. We’ll also \textbf{throw a party} for the teaching staff on Friday. \textbf{The more the merrier} — come along and \textbf{have a blast}!
\end{textebox}

\begin{center}
\begin{tabularx}{\linewidth}{|X|X|X|}
\hline
\rowcolor{popblueL} \textbf{Collocation (Verb + Noun)} & \textbf{Traduction française} & \textbf{Contexte / Niveau de langue} \\
\hline
\textbf{to throw a party} & organiser une fête (souvent informelle, entre amis) & Courant, neutre à familier. \eng{We’re throwing a party Saturday.} \\
\textbf{to hold a celebration} & organiser une célébration (plus formel, officiel) & Cérémonies officielles, anniversaires d’institution. \eng{The school held a celebration.} \\
\textbf{to organise a party / a festival} & organiser une fête / un festival & Standard, polyvalent. \eng{They organised a music festival.} \\
\textbf{to attend a party / a ceremony} & assister à une fête / une cérémonie & Formel. \eng{The Minister attended the ceremony.} \\
\textbf{to host a party / an event} & accueillir / héberger une fête (chez soi, ou en tant qu’organisateur principal) & Responsabilité de l’hôte. \eng{My uncle hosted the wedding party.} \\
\textbf{to enjoy oneself} & s’amuser, passer un bon moment & Réfléchi (oneself). \eng{Did you enjoy yourself at the Ngondo?} \\
\textbf{to have a great time / a good time} & passer un excellent / bon moment & Très courant, neutre. \eng{We had a great time at the get-together.} \\
\hline
\end{tabularx}
\end{center}

\begin{center}
\begin{tabularx}{\linewidth}{|X|X|X|}
\hline
\rowcolor{popgreenL} \textbf{Collocation (Verb + Noun)} & \textbf{Traduction française} & \textbf{Contexte / Niveau de langue} \\
\hline
\textbf{to organise a festival} & organiser un festival & Standard. \eng{The council organises a festival every December.} \\
\textbf{to take part in an event / a festival} & participer à un événement / un festival & Actif. \eng{Students took part in the cultural event.} \\
\textbf{to celebrate an anniversary} & fêter un anniversaire (date précise) & \eng{We celebrate our independence anniversary on May 20.} \\
\textbf{to mark an occasion / an anniversary} & marquer un événement / un anniversaire (par un acte symbolique) & Plus solennel. \eng{They planted a tree to mark the occasion.} \\
\textbf{to commemorate an event} & commémorer un événement (souvent historique, triste ou solennel) & Formel, historique. \eng{We commemorate the veterans.} \\
\hline
\end{tabularx}
\end{center}

\begin{center}
\begin{tabularx}{\linewidth}{|X|X|X|}
\hline
\rowcolor{poporangeL} \textbf{Collocation (Verb + Noun / Preposition)} & \textbf{Traduction française} & \textbf{Structure grammaticale} \\
\hline
\textbf{to send an invitation (to sb)} & envoyer une invitation (à qqn) & \eng{send} + objet + \eng{to} + destinataire \\
\textbf{to accept an invitation} & accepter une invitation & \eng{accept} + objet direct \\
\textbf{to decline an invitation} & refuser une invitation (poliment) & \eng{decline} + objet direct (plus formel que \eng{refuse}) \\
\textbf{to invite someone \textbf{to} + V / \textbf{to} + N} & inviter qqn à faire qqch / à qqch & \eng{invite} + objet + \eng{to} + infinitif / nom \\
\textbf{to be invited \textbf{to} + N} & être invité à qqch (passif) & \eng{be invited to} + événement \\
\textbf{to reply to an invitation / to RSVP} & répondre à une invitation & \eng{RSVP} = \eng{Répondez s’il vous plaît} (français dans l’anglais formel) \\
\hline
\end{tabularx}
\end{center}

\begin{exemplebox}[Exemples traduits — Collocations de base]
\eng{We are \textbf{throwing a party} for our school’s 50th anniversary.} \\
« Nous organisons une fête pour le 50e anniversaire de notre école. » \\[0.5em]
\eng{The Principal \textbf{held a celebration} to \textbf{mark the occasion}.} \\
« Le Proviseur a organisé une célébration pour marquer l’événement. » \\[0.5em]
\eng{Did you \textbf{receive the invitation}? You must \textbf{reply by Friday}.} \\
« As-tu reçu l’invitation ? Tu dois répondre avant vendredi. » \\[0.5em]
\eng{Many former students \textbf{accepted the invitation} to \textbf{attend the ceremony}.} \\
« De nombreux anciens élèves ont accepté l’invitation d’assister à la cérémonie. »
\end{exemplebox}

\begin{attention}[Piège majeur : « Do a party » n’existe pas !]
\emph{Erreur très fréquente chez les francophones :} \textbf{\eng{*We did a party yesterday.}} (calque de « Nous avons fait une fête hier. »)
\begin{itemize}
    \item \textbf{Correct :} \eng{We \textbf{threw} a party yesterday.} (verbe irrégulier : throw / threw / thrown)
    \item \textbf{Correct :} \eng{We \textbf{had} a party yesterday.}
    \item \textbf{Correct :} \eng{We \textbf{organised} a party yesterday.}
\end{itemize}
\eng{Do} ne se dit qu’avec \eng{exercise}, \eng{homework}, \eng{the dishes}, \eng{a favour}, \eng{business}… Jamais avec \eng{party}, \eng{festival}, \eng{celebration}, \eng{ceremony}.
\end{attention}

\begin{exoresolu}[Mettre le verbe entre parenthèses à la forme correcte et choisir la bonne collocation]
\textbf{Énoncé :} Complétez les phrases avec la collocation appropriée (forme verbale correcte).
\begin{enumerate}
    \item The Students’ Union \_\_\_\_\_\_\_\_ (decide) to \_\_\_\_\_\_\_\_ a big party for the end of year.
    \item All students are expected to \_\_\_\_\_\_\_\_ (attend) the flag-raising \_\_\_\_\_\_\_\_ (ceremony / celebration).
    \item We \_\_\_\_\_\_\_\_ (send) the invitations last week; so far, fifty people \_\_\_\_\_\_\_\_ (accept).
    \item To \_\_\_\_\_\_\_\_ (mark) the 50th anniversary, the school \_\_\_\_\_\_\_\_ (plant) a tree in the courtyard.
    \item I hope you \_\_\_\_\_\_\_\_ (enjoy) \_\_\_\_\_\_\_\_ at the festival tomorrow!
\end{enumerate}

\tcblower
\textbf{Solution détaillée :}
\begin{enumerate}
    \item \textbf{decided to \eng{throw} / \eng{organise} / \eng{hold}} a big party. (\eng{decide to} + infinitif ; \eng{throw/organise/hold a party} sont les collocations standards. \eng{Do a party} est incorrect — voir \textbf{Attention} ci-dessus.)
    \item \textbf{attend} the flag-raising \textbf{ceremony}. (\eng{attend a ceremony} = collocation formelle pour un événement officiel/solennel. \eng{Celebration} convient moins pour la levée des couleurs quotidienne.)
    \item \textbf{sent} the invitations last week; so far, fifty people \textbf{have accepted}. (\eng{send invitations} : passé simple \eng{sent} car action datée « last week » ; \eng{accept an invitation} : present perfect \eng{have accepted} car conséquence sur le présent « so far ».)
    \item \textbf{mark} the 50th anniversary, the school \textbf{planted} a tree. (\eng{to mark an occasion/anniversary} = but (infinitif de but) ; \eng{planted} = passé simple action ponctuelle passée.)
    \item \textbf{enjoy} \textbf{yourself} at the festival tomorrow! (\eng{enjoy oneself} = verbe pronominal obligatoire. \eng{Enjoy} seul est transitif direct : \eng{enjoy the festival}, mais \eng{enjoy at the festival} est incorrect.)
\end{enumerate}
\end{exoresolu}

\courssub{Idiomatic expressions for celebrations and socialising}

Les \cle{idiom}s (expressions idiomatiques) sont des phrases figées dont le sens global ne se déduit pas du sens des mots pris séparément. Ils sont omniprésents à l’oral et dans les messages informels. En voici cinq incontournables pour les fêtes et rencontres.

\begin{center}
\begin{tabularx}{\linewidth}{|X|X|X|}
\hline
\rowcolor{poppurpleL} \textbf{Idiom} & \textbf{Sens (français)} & \textbf{Contexte d’emploi + Exemple} \\
\hline
\textbf{the more the merrier} & plus on est de fous, plus on rit & Invitation ouverte, accueil chaleureux. \eng{« Can I bring my cousin? » — « Sure, \textbf{the more the merrier}! »} \\
\textbf{to have a blast} & s’éclater, passer un moment génial & Très courant chez les jeunes, familier positif. \eng{We \textbf{had a blast} at the concert last night!} \\
\textbf{to paint the town red} & faire la fête (sortir, boire, danser) toute la nuit & Imagé, un peu vieilli mais toujours utilisé. \eng{After exams, we’re going to \textbf{paint the town red}.} \\
\textbf{a get-together} (noun) & une réunion amicale, une petite fête entre connaissances & Neutre, moins formel que \eng{party}. \eng{We’re having a \textbf{get-together} at my place Sunday.} \\
\textbf{to break the ice} & briser la glace (détendre l’atmosphère au début d’une rencontre) & Social, professionnel. \eng{He told a joke to \textbf{break the ice} at the meeting.} \\
\hline
\end{tabularx}
\end{center}

\begin{exemplebox}[Exemples traduits — Idioms en situation]
\eng{Come along to the village feast — \textbf{the more the merrier}!} \\
« Viens à la fête du village — plus on est de fous, plus on rit ! » \\[0.5em]
\eng{We \textbf{had a blast} at the Ngondo festival in Douala!} \\
« On s’est vraiment éclatés au festival du Ngondo à Douala ! » \\[0.5em]
\eng{It was a nice informal \textbf{get-together} with the neighbours.} \\
« C’était une sympathique petite réunion entre voisins. »
\end{exemplebox}

\begin{savaistu}[« To have a blast » — un classique des jeunes]
L’expression \eng{\textbf{to have a blast}} est \emph{très} courante dans l’anglais parlé actuel (surtout Amérique du Nord, mais répandue partout via les séries/réseaux sociaux). À l’origine, \eng{a blast} = une explosion, un coup de vent violent. Au sens figuré : « une expérience intense et jouissive ».
\begin{itemize}
    \item \eng{We had a blast at the festival!} = On s’est éclatés / On a trop kiffé le festival !
    \item \eng{Have a blast!} = Éclate-toi bien ! (souhait avant une fête).
\end{itemize}
Attention : ne pas confondre avec \eng{to have a blast} (subir une explosion) au sens littéral militaire/minier. Le contexte lève l’ambiguïté.
\end{savaistu}

\begin{popfigure}
\begin{tikzpicture}[
    mindmap,
    concept color=popblue,
    level 1/.style={level distance=3.5cm, sibling angle=72},
    level 2/.style={level distance=2.5cm},
    every node/.style={font=\small, text=white},
    grow cyclic
]
\node[concept, scale=1.2] {CELEBRATE}
    child[concept color=poporange] { node[concept] {an anniversary} }
    child[concept color=popgreen] { node[concept] {Christmas} }
    child[concept color=poppurple] { node[concept] {a victory} }
    child[concept color=poppink] { node[concept] {together} }
    child[concept color=popgold] { node[concept] {with friends} };
\end{tikzpicture}
\legende{Carte mentale lexicale : collocations verbales autour de \eng{celebrate}.}
\end{popfigure}

\courssub{Functional language: Inviting and responding (speaking/writing)}

Au-delà du vocabulaire, il faut maîtriser les \textbf{formules toutes faites} (chunks) pour inviter et répondre. Le choix dépend du \cle{register} vu en section 2.

\begin{center}
\begin{tabularx}{\linewidth}{|X|X|X|}
\hline
\rowcolor{popblueL} \textbf{Inviting someone} & \textbf{Niveau / Register} & \textbf{Traduction / Nuance} \\
\hline
\eng{Would you like to come to...?} & Standard, poli (adultes, autorités) & « Voulez-vous venir à… ? » \\
\eng{We would be honoured if you could attend...} & Très formel (autorités, invités d’honneur) & « Nous serions honorés si vous pouviez assister… » \\
\eng{How about coming to...?} & Semi-formel / Amical (voisins, collègues) & « Et si tu venais à… ? » (suggestion) \\
\eng{You’re welcome to join us for...} & Chaleureux, inclusif (camarades, voisins) & « Tu es le bienvenu pour te joindre à nous… » \\
\eng{I’d love to invite you to...} & Personnel, sincère (amis, connaissances) & « J’aimerais t’inviter à… » \\
\eng{Do you fancy coming to...?} & Familier (très proches, camarades) & « Tu as envie de venir à… ? » (UK) \\
\eng{Why don’t you come along to...?} & Amical, incitatif & « Pourquoi ne viendrais-tu pas à… ? » \\
\hline
\end{tabularx}
\end{center}

\begin{center}
\begin{tabularx}{\linewidth}{|X|X|X|}
\hline
\rowcolor{popgreenL} \textbf{Accepting an invitation} & \textbf{Niveau / Register} & \textbf{Traduction / Nuance} \\
\hline
\eng{I’d love to! / I’d love to come.} & Universel, enthousiaste & « J’adorerais ! / J’adorerais venir. » \\
\eng{That sounds great / wonderful!} & Enthousiaste, positif & « Ça a l’air super / génial ! » \\
\eng{Thank you for inviting me. I’ll be there.} & Poli, standard & « Merci de m’inviter. Je serai là. » \\
\eng{Count me in!} & Familier, engagement fort & « Compte sur moi ! / Je suis partant ! » \\
\eng{With pleasure.} & Formel, élégant & « Avec plaisir. » \\
\hline
\rowcolor{poporangeL} \textbf{Declining an invitation (politely)} & \textbf{Niveau / Register} & \textbf{Traduction / Nuance} \\
\hline
\eng{I’m afraid I can’t make it.} & Standard, poli (la formule passe-partout) & « Je crains de ne pas pouvoir venir. » \\
\eng{Thanks for inviting me, but I have a prior engagement.} & Poli, raison vague mais acceptable & « Merci de m’inviter, mais j’ai un empêchement. » \\
\eng{I’d love to, but I’m not available that day.} & Regretté, raison précise & « J’aimerais bien, mais je ne suis pas dispo ce jour. » \\
\eng{Sorry, I can’t come.} & Direct, entre proches & « Désolé, je ne peux pas venir. » \\
\hline
\end{tabularx}
\end{center}

\begin{attention}[Erreurs fréquentes — Inviter et répondre]
\begin{itemize}
    \item \textbf{\eng{*I invite you to come.}} $\rightarrow$ Trop direct, sonne comme un ordre ou une citation de théâtre. Préférez \eng{I’d like to invite you to...} ou \eng{Would you like to...?}
    \item \textbf{\eng{*I can’t come because I have a homework.}} $\rightarrow$ \eng{Homework} est \textbf{indénombrable} : \eng{I have homework} / \eng{I have some homework} / \eng{I have a lot of homework}. Jamais \eng{a homework}.
    \item \textbf{\eng{*Thanks for the invitation. I accept.}} $\rightarrow$ Trop formel/robotique à l’oral. \eng{I’d love to!} ou \eng{That sounds great!} sont naturels.
    \item \textbf{Préposition :} \eng{invite someone \textbf{TO} + verbe} (\eng{invite you to come}), \eng{invite someone \textbf{TO} + nom d’événement} (\eng{invite you to the party}). Pas \eng{invite someone FOR}.
\end{itemize}
\end{attention}

\courssub{Word families, pronunciation and spelling}

\begin{propriete}[Dérivation nominale et adjectivale (Word Formation)]
L’anglais construit des familles de mots par \textbf{affixation} (préfixes, suffixes). Connaître la racine permet de deviner le sens et la catégorie grammaticale.
\begin{center}
\begin{tabular}{|l|l|l|l|}
\hline
\rowcolor{popblueL} \textbf{Verbe} & \textbf{Nom (action/état)} & \textbf{Nom (agent)} & \textbf{Adjectif} \\
\hline
\eng{celebrate} & \eng{celebration} & \eng{celebrant} & \eng{celebratory} \\
\eng{organise} & \eng{organisation} & \eng{organiser} & \eng{organisational} \\
\eng{invite} & \eng{invitation} & \eng{invitee} & \eng{inviting} \\
\eng{attend} & \eng{attendance} & \eng{attendee} & \eng{attentive} \\
\eng{festivate} (rare) & \eng{festival, festivity} & — & \eng{festive} \\
\eng{ceremonialise} (rare) & \eng{ceremony} & — & \eng{ceremonial} \\
\eng{commemorate} & \eng{commemoration} & — & \eng{commemorative} \\
\hline
\end{tabular}
\end{center}
\textbf{Règles clés :}
\begin{itemize}
    \item Verbe $\to$ Nom d’action : suffixe \textbf{-tion / -sion} (\eng{celebrate $\to$ celebration}, \eng{organise $\to$ organisation}). L’orthographe britannique privilégie \textbf{-ise} / \textbf{-isation} (pas \eng{-ize} / \eng{-ization}, sauf Oxford spelling).
    \item Verbe $\to$ Agent : suffixe \textbf{-er / -or / -ee} (\eng{organiser}, \eng{invitee} = celui qui reçoit l’invitation).
    \item Verbe/Adjectif $\to$ Adjectif : suffixe \textbf{-ive, -al, -ous} (\eng{celebratory}, \eng{ceremonial}, \eng{festive}).
\end{itemize}
\end{propriete}

\begin{exemplebox}[Familles de mots en phrases]
\eng{The \textbf{celebration} was \textbf{celebratory} in tone.} (La célébration était de ton festif.) \\
\eng{The \textbf{organiser} sent the \textbf{invitations} for the \textbf{festival}.} (L’organisateur a envoyé les invitations pour le festival.) \\
\eng{All \textbf{attendees} must confirm their \textbf{attendance}.} (Tous les participants doivent confirmer leur présence.) \\
\eng{It was a \textbf{festive} atmosphere during the \textbf{ceremony}.} (C’était une ambiance festive pendant la cérémonie.)
\end{exemplebox}

\begin{attention}[Pièges d’orthographe et de prononciation (Francophones)]
\begin{itemize}
    \item \textbf{\eng{Celebration}} : s’écrit avec \textbf{un seul \eng{l}} (contrairement à \eng{celebrate} qui en a un) et \textbf{deux \eng{t}} ? Non : \eng{c-e-l-e-b-r-a-t-i-o-n} (un \eng{l}, un \eng{t} avant \eng{ion}). Prononciation : « sè-lé-\textbf{bra}-ssion » (accent sur \eng{bra}).
    \item \textbf{\eng{Organisation}} : \textbf{\eng{s}} (britannique), pas \eng{z}. Prononciation : « or-ga-ni-\textbf{za}-ssion » (accent sur \eng{za}).
    \item \textbf{\eng{Invitation}} : \eng{invite} (verbe, accent sur 2e syllabe « in-\textbf{vaït} ») $\to$ \eng{invitation} (nom, accent sur 3e syllabe « in-vi-\textbf{ta}-ssion »). \textbf{Règle :} l’accent glisse vers le suffixe \eng{-tion}.
    \item \textbf{\eng{Attendance}} : \textbf{\eng{e}} après \eng{t} (\eng{attend} + \eng{ance}), pas \eng{a} (*attandance).
    \item \textbf{\eng{Commemorate / Commemoration}} : \textbf{double \eng{m}}, \textbf{simple \eng{r}}.
    \item \textbf{\eng{Festival}} : accent sur 1re syllabe « \textbf{fes}-ti-val » (pas « fes-ti-\textbf{val} »).
\end{itemize}
\end{attention}

\begin{exemplebox}[Guide de prononciation approximatif (français)]
\begin{itemize}
    \item \eng{celebrate} : « sè-lé-\textbf{braït} » \quad \eng{celebration} : « sè-lé-\textbf{bra}-ssion »
    \item \eng{organise} : « or-ga-\textbf{na-ïz} » \quad \eng{organisation} : « or-ga-ni-\textbf{za}-ssion »
    \item \eng{invite} : « in-\textbf{va-ït} » \quad \eng{invitation} : « in-vi-\textbf{ta}-ssion »
    \item \eng{attend} : « a-\textbf{tend} » \quad \eng{attendance} : « a-\textbf{ten}-danss »
    \item \eng{festival} : « \textbf{fes}-ti-val » \quad \eng{festive} : « \textbf{fes}-tiv »
    \item \eng{ceremony} : « sé-\textbf{ré}-mo-ni » \quad \eng{ceremonial} : « sé-ré-\textbf{mo}-ni-al »
\end{itemize}
\end{exemplebox}

\courssub{Using the vocabulary in context: a dialogue}

Le dialogue suivant met en scène deux élèves de 1ère A (Mireille et Jean) et leur professeur principal (M. Njoya) dans le bureau de ce dernier. Ils préparent la fête de fin d’année. Observez les collocations, idioms, registres de langue et formules d’invitation/réponse.

\begin{textebox}[Dialogue: Planning the End-of-Year Party — Lycée de Bafoussam]
\textbf{Mireille:} (knocking) \eng{Good afternoon, Sir. May we come in?}

\textbf{Mr Njoya:} \eng{Good afternoon. Come in, please. Take a seat. What can I do for you?}

\textbf{Jean:} \eng{Thank you, Sir. We’re here about the end-of-year \textbf{get-together}. The Students’ Union wants to \textbf{organise a party} for the final-year students next Friday.}

\textbf{Mireille:} \eng{We’d like to \textbf{invite you to attend}, Sir. And we hope you’ll \textbf{send an invitation} to the Principal and the Divisional Delegate as well.}

\textbf{Mr Njoya:} \eng{That’s a lovely initiative. \textbf{Who will host the event}?}

\textbf{Jean:} \eng{We’ll \textbf{hold the celebration} in the school hall. The Parents’ Association has offered to \textbf{provide refreshments}.}

\textbf{Mireille:} \eng{We’ve also planned some games to \textbf{break the ice} between the French-speaking and English-speaking sections. \textbf{The more the merrier}, right?}

\textbf{Mr Njoya:} \eng{Absolutely. Social cohesion matters. \textbf{I’d be delighted to come}. I’ll \textbf{accept the invitation} officially once I receive the written note.}

\textbf{Jean:} \eng{Great! We’ll \textbf{send the invitation} tomorrow.}

\textbf{Mireille:} \eng{And Sir, don’t worry — we’ll make sure everyone \textbf{has a blast}!}

\textbf{Mr Njoya:} \eng{(smiling) I’m sure you will. Now, go back to class. You have a test in ten minutes!}

\textbf{Both:} \eng{Oh no! \textbf{We have to run}. Thanks, Sir! See you Friday!}
\end{textebox}

\begin{center}
\begin{tabularx}{\linewidth}{|l|X|}
\hline
\rowcolor{popmuted} \textbf{Mot / Expression} & \textbf{Glossaire (nature, traduction)} \\
\hline
\eng{get-together} & (noun) réunion amicale, petite fête entre connaissances \\
\eng{organise a party} & (collocation) organiser une fête \\
\eng{invite someone to attend} & (collocation) inviter qqn à assister (à un événement) \\
\eng{send an invitation} & (collocation) envoyer une invitation \\
\eng{host an event} & (collocation) accueillir / être l’hôte principal d’un événement \\
\eng{hold a celebration} & (collocation) organiser une célébration (formel) \\
\eng{provide refreshments} & (collocation) fournir les rafraîchissements / le goûter \\
\eng{break the ice} & (idiom) briser la glace, détendre l’atmosphère \\
\eng{the more the merrier} & (idiom) plus on est de fous, plus on rit \\
\eng{I’d be delighted to} & (formule polie) j’en serais ravi / j’aimerais beaucoup \\
\eng{have a blast} & (idiom, familier) s’éclater, passer un moment génial \\
\eng{have to run} & (expression familière) devoir filer / partir en vitesse \\
\hline
\end{tabularx}
\end{center}

\begin{experience}[Jeu de rôle : « L’invitation officielle »]
\textbf{Objectif :} Réemployer le lexique (collocations, idioms, formules) et adapter le \cle{register} à l’interlocuteur.
\textbf{Consigne :} Par groupes de 3. Un élève joue le \textbf{représentant des élèves}, un autre le \textbf{proviseur / autorité}, le troisième un \textbf{voisin / parent d’élève}.
\begin{enumerate}
    \item Le représentant doit \textbf{inviter} les deux autres à la fête de fin d’année (date, lieu, tenue).
    \item Le proviseur \textbf{accepte} formellement (formule écrite/orale).
    \item Le voisin \textbf{refuse} poliment (empêchement).
    \item Utilisez au minimum : 3 collocations (\eng{hold a celebration}, \eng{send invitation}, \eng{attend ceremony}…), 1 idiom (\eng{the more the merrier}, \eng{have a blast}…), 1 mot de la famille \eng{celebrate/organise/invite}.
\end{enumerate}
\textbf{Débriefing :} La classe note les erreurs de registre (tutoiement inapproprié, \eng{do a party}, \eng{invite for}…) et les bonnes formulations.
\end{experience}

\courssub{Consolidation exercises}

\begin{exercice}[Entraînement : Collocations, Idioms et Formules d’invitation]
\textbf{A. Complétez avec la collocation correcte (forme verbale appropriée).}
\begin{enumerate}
    \item The association decided to \_\_\_\_\_\_\_\_ a festival to \_\_\_\_\_\_\_\_ the town’s founding anniversary.
    \item You must \_\_\_\_\_\_\_\_ the invitations at least two weeks before the ceremony.
    \item Hundreds of people \_\_\_\_\_\_\_\_ the Independence Day celebration at the stadium yesterday.
    \item We \_\_\_\_\_\_\_\_ a great time at the cultural week. We really \_\_\_\_\_\_\_\_ ourselves.
    \item The Head Teacher \_\_\_\_\_\_\_\_ us to the prize-giving ceremony.
\end{enumerate}

\textbf{B. Remplacez la partie soulignée par un idiom de la leçon.}
\begin{enumerate}
    \item \eng{Everyone is welcome to come to the party.} $\rightarrow$ \eng{Come on, \_\_\_\_\_\_\_\_!}
    \item \eng{We had an amazing time at the concert.} $\rightarrow$ \eng{We \_\_\_\_\_\_\_\_ at the concert.}
    \item \eng{They went out and partied all night after the exam.} $\rightarrow$ \eng{They \_\_\_\_\_\_\_\_.}
    \item \eng{He told a joke to make everyone relax at the start of the meeting.} $\rightarrow$ \eng{He told a joke to \_\_\_\_\_\_\_\_.}
    \item \eng{We are having a small friendly meeting at my house.} $\rightarrow$ \eng{We are having a \_\_\_\_\_\_\_\_ at my house.}
\end{enumerate}

\textbf{C. Réécrivez les phrases suivantes en adaptant le registre (formel / informel) indiqué.}
\begin{enumerate}
    \item (Formel — à un inspecteur) \eng{Hey, come to our party Friday!} $\rightarrow$ \eng{\_\_\_\_\_\_\_\_}
    \item (Informel — à un camarade) \eng{I would be honoured if you could attend.} $\rightarrow$ \eng{\_\_\_\_\_\_\_\_}
    \item (Formel — réponse à une invitation officielle) \eng{Yeah, sure, I’ll be there.} $\rightarrow$ \eng{\_\_\_\_\_\_\_\_}
    \item (Informel — refus à un ami) \eng{I regret to inform you that I am unable to attend.} $\rightarrow$ \eng{\_\_\_\_\_\_\_\_}
\end{enumerate}

\textbf{D. Formation de mots (Word Formation). Complétez le tableau.}
\begin{center}
\begin{tabular}{|l|l|l|l|}
\hline
\rowcolor{popblueL} \textbf{Verbe} & \textbf{Nom (action)} & \textbf{Nom (personne)} & \textbf{Adjectif} \\
\hline
\eng{celebrate} & \eng{celebration} & \eng{\_\_\_\_\_\_\_\_} & \eng{\_\_\_\_\_\_\_\_} \\
\hline
\eng{organise} & \eng{\_\_\_\_\_\_\_\_} & \eng{organiser} & \eng{\_\_\_\_\_\_\_\_} \\
\hline
\eng{invite} & \eng{\_\_\_\_\_\_\_\_} & \eng{\_\_\_\_\_\_\_\_} & \eng{inviting} \\
\hline
\eng{attend} & \eng{\_\_\_\_\_\_\_\_} & \eng{attendee} & \eng{\_\_\_\_\_\_\_\_} \\
\hline
\eng{commemorate} & \eng{\_\_\_\_\_\_\_\_} & — & \eng{\_\_\_\_\_\_\_\_} \\
\hline
\end{tabular}
\end{center}
\end{exercice}

\begin{corrige}[Exercice — Collocations, Idioms, Registre et Word Formation]
\textbf{A. Collocations}
\begin{enumerate}
    \item \textbf{organise / hold} a festival to \textbf{mark / celebrate} the town’s founding anniversary. (\eng{mark/celebrate an anniversary})
    \item \textbf{send} the invitations at least two weeks before the ceremony. (\eng{send invitations})
    \item \textbf{attended} the Independence Day celebration at the stadium yesterday. (\eng{attend a celebration/ceremony} — passé simple)
    \item \textbf{had} a great time at the cultural week. We really \textbf{enjoyed} ourselves. (\eng{have a great time}, \eng{enjoy oneself} — pronominal !)
    \item \textbf{invited} us \textbf{to} the prize-giving ceremony. (\eng{invite someone to + event})
\end{enumerate}

\textbf{B. Idioms}
\begin{enumerate}
    \item \eng{Come on, \textbf{the more the merrier}!}
    \item \eng{We \textbf{had a blast} at the concert.}
    \item \eng{They \textbf{painted the town red}.}
    \item \eng{He told a joke to \textbf{break the ice}.}
    \item \eng{We are having a \textbf{get-together} at my house.}
\end{enumerate}

\textbf{C. Adaptation du registre (propositions) }
\begin{enumerate}
    \item \eng{We would be delighted if you could attend our celebration on Friday. / We cordially invite you to our end-of-year celebration on Friday.}
    \item \eng{Hey, wanna come to the party Friday? / Do you fancy coming to the get-together Friday?}
    \item \eng{Thank you for the invitation. I would be honoured to attend. / I accept with pleasure.}
    \item \eng{Sorry, I can’t make it. / Thanks for asking, but I’m busy Friday.}
\end{enumerate}

\textbf{D. Word Formation}
\begin{center}
\begin{tabular}{|l|l|l|l|}
\hline
\rowcolor{popblueL} \textbf{Verbe} & \textbf{Nom (action)} & \textbf{Nom (personne)} & \textbf{Adjectif} \\
\hline
\eng{celebrate} & \eng{celebration} & \eng{celebrant} & \eng{celebratory} \\
\hline
\eng{organise} & \eng{organisation} & \eng{organiser} & \eng{organisational} \\
\hline
\eng{invite} & \eng{invitation} & \eng{invitee} & \eng{inviting} \\
\hline
\eng{attend} & \eng{attendance} & \eng{attendee} & \eng{attentive} \\
\hline
\eng{commemorate} & \eng{commemoration} & — & \eng{commemorative} \\
\hline
\end{tabular}
\end{center}
\end{corrige}

\coursec{False friends and common traps for francophones}

After mastering collocations and idiomatic expressions to sound natural during community celebrations, we must now tackle a silent killer of marks in examinations: \cle{false friends} (faux amis). These words look reassuringly familiar because they resemble French words, but their meanings have drifted apart over centuries. Using them incorrectly — saying \eng{I assisted the ceremony} for « J'ai assisté à la cérémonie » or \eng{I demand you to come} for « Je vous invite à venir » — instantly marks a candidate as a francophone learner and, in the Baccalauréat, counts as a lexical error. This section trains you to spot them, understand the semantic gap, and replace them with the correct English equivalents.

\courssub{Observation: when familiar words deceive us}

\begin{textebox}[A misunderstanding at the school festival]
The Principal of Government Bilingual High School Yaoundé stood at the podium during the End-of-Year Cultural Festival. He smiled at the crowd of students, parents, and neighbours. \eng{``We are delighted that so many parents could attend today,''} he began. \eng{``Your presence honours us. A special thank you to the Parents' Association who helped organise the refreshments. We requested your support last month, and you responded generously.''}

Near the stage, a French-speaking student whispered to his friend: \eng{``The Principal said many parents \textbf{assisted} the festival.''} His friend, an Anglophone student from Bamenda, frowned. \eng{``No, he said they \textbf{attended}. To \textbf{assist} means to help. The parents \textbf{attended} the festival; the Parents' Association \textbf{assisted} with the organisation.''}

Later, reading a student's report on the event, the English teacher underlined a sentence in red: \eng{``The students \textbf{demanded} the Principal to open the buffet.''} She wrote in the margin: \eng{``\textbf{Demand} = \emph{exiger}. You mean \textbf{asked} or \textbf{requested}. This changes the tone completely!''}
\end{textebox}

\begin{exemplebox}[Analyse des erreurs du texte]
\begin{itemize}
    \item \eng{\textbf{Assisted} the festival} $\rightarrow$ Faux ami. Le Principal a dit \eng{attend} (être présent). \eng{Assist} = aider.
    \item \eng{\textbf{Demanded} the Principal} $\rightarrow$ Faux ami. \eng{Demand} = exiger (ton agressif/autoritaire). Le sens voulu est \eng{asked} ou \eng{requested} (demander poliment).
    \item \eng{\textbf{Ceremony}} $\rightarrow$ Attention : en anglais, \eng{ceremony} désigne la partie formelle (discours, remise de prix), pas la fête dansante qui suit.
\end{itemize}
\end{exemplebox}

\courssub{The Big Five false friends: form, meaning, and correct equivalents}

The following table summarises the five most dangerous false friends for this topic. Study the \cle{real English meaning} and the \cle{correct French equivalent} carefully. The \cle{trap} column shows the erroneous French thought process.

\begin{popfigure}
\legende{Tableau récapitulatif des 5 faux amis majeurs — Mots anglais, sens français trompeur, sens réel et équivalents corrects}
\begin{center}
\begin{tabularx}{\linewidth}{|>{\bfseries}p{2.8cm}|X|X|}\hline
\rowcolor{popblue} \textcolor{white}{\textbf{English Word}} & \textcolor{white}{\textbf{Faux sens français (Piège)}} & \textcolor{white}{\textbf{Sens réel en anglais + Équivalent correct}} \\ \hline
to \textbf{assist} & assister (à un événement) & \textbf{to help / to aid} — \textit{aider} \\ 
& & \eng{Attend} / \eng{be present at} = \textit{assister à} \\ \hline
\textbf{eventually} & éventuellement & \textbf{finally / in the end} — \textit{finalement, au bout du compte} \\
& & \eng{Possibly} / \eng{if necessary} = \textit{éventuellement} \\ \hline
a \textbf{ceremony} & la fête / la cérémonie (au sens large) & \textbf{formal ritual / official act} — \textit{cérémonie formelle (remise de diplômes, mariage civil)} \\
& & \eng{A party / a celebration / a festival} = \textit{la fête, la célébration} \\ \hline
to \textbf{demand} & demander (poliment) & \textbf{to require forcefully / to claim as a right} — \textit{exiger, réclamer} \\
& & \eng{Ask (for)} / \eng{request} = \textit{demander (poliment)} \\ \hline
an \textbf{occasion} & une opportunité / une chance & \textbf{a particular time / a special event} — \textit{une occasion, un événement spécial} \\
& & \eng{An opportunity} / \eng{a chance} = \textit{une opportunité, une chance} \\ \hline
\end{tabularx}
\end{center}
\end{popfigure}

\courssub{Detailed analysis of each false friend}

\begin{propriete}[Faux ami 1 : \textbf{to assist} vs \textit{assister à}]
\eng{To assist} est un verbe transitif direct : il signifie « apporter de l'aide à quelqu'un » (\eng{help}). Il ne s'utilise \textbf{jamais} avec le sens de « être présent à un événement ».
\begin{itemize}
    \item \textbf{Forme} : \eng{assist + personne/objet} (pas de préposition \eng{to} après le verbe).
    \item \textbf{Emploi} : \eng{The volunteers assisted the organisers.} (Les bénévoles ont aidé les organisateurs.)
    \item \textbf{Correction} : Pour « assister à » (être spectateur), on utilise \eng{attend} (formel/neutre) ou \eng{be present at} (très formel). \eng{Go to} est courant à l'oral.
\end{itemize}
\begin{center}
\begin{tabularx}{\linewidth}{|X|X|X|}\hline
\rowcolor{popblueL} \textbf{Phrase incorrecte (Calque)} & \textbf{Phrase correcte} & \textbf{Traduction} \\ \hline
\eng{I assisted the wedding.} & \eng{I \textbf{attended} the wedding.} & J'ai assisté au mariage. \\ \hline
\eng{Did you assist the match?} & \eng{Did you \textbf{attend} the match?} & As-tu assisté au match ? \\ \hline
\eng{Many parents assisted.} & \eng{Many parents \textbf{attended}.} & Beaucoup de parents ont assisté. \\ \hline
\end{tabularx}
\end{center}
\end{propriete}

\begin{propriete}[Faux ami 2 : \textbf{eventually} vs \textit{éventuellement}]
\eng{Eventually} est un adverbe de temps qui marque l'aboutissement d'un processus long ou incertain : « finalement, au bout du compte, fin par ». Il n'exprime \textbf{pas} la possibilité conditionnelle.
\begin{itemize}
    \item \textbf{Emploi} : \eng{We walked for hours and eventually found the village.} (Nous avons marché des heures et avons \textbf{finalement} trouvé le village.)
    \item \textbf{Correction} : Pour « éventuellement » (si nécessaire, peut-être), on utilise \eng{possibly}, \eng{if necessary}, \eng{if need be}, ou la structure \eng{might}.
\end{itemize}
\begin{center}
\begin{tabularx}{\linewidth}{|X|X|X|}\hline
\rowcolor{poporangeL} \textbf{Phrase incorrecte} & \textbf{Phrase correcte} & \textbf{Traduction} \\ \hline
\eng{Eventually, we will invite the chief.} & \eng{\textbf{Possibly} / \textbf{If necessary}, we will invite the chief.} & Éventuellement, nous inviterons le chef. \\ \hline
\eng{You can eventually come tomorrow.} & \eng{You \textbf{might} come tomorrow.} / \eng{You can come tomorrow \textbf{if you wish}.} & Tu peux éventuellement venir demain. \\ \hline
\end{tabularx}
\end{center}
\end{propriete}

\begin{propriete}[Faux ami 3 : \textbf{a ceremony} vs \textit{la fête / la cérémonie}]
En anglais, \eng{ceremony} a un sens \textbf{restreint et formel} : acte solennel, rituel officiel (remise de diplômes, mariage civil, investiture, funérailles). Il ne désigne \textbf{pas} la partie festive (musique, danse, repas).
\begin{itemize}
    \item \textbf{Emploi correct} : \eng{The graduation ceremony starts at 10 a.m.} (La cérémonie de remise de diplômes commence à 10h.)
    \item \textbf{Pour la fête} : \eng{a party}, \eng{a celebration}, \eng{a festival}, \eng{a get-together} (informel), \eng{a reception} (réception formelle après cérémonie).
\end{itemize}
\end{propriete}

\begin{exemplebox}[Distinction claire]
\eng{The \textbf{ceremony} was short, but the \textbf{celebration} lasted all night.} \\
« La \textbf{cérémonie} (officielle) a été courte, mais la \textbf{fête} a duré toute la nuit. »
\end{exemplebox}


\begin{propriete}[Faux ami 4 : \textbf{to demand} vs \textit{demander}]
\eng{To demand} implique une relation de force, d'autorité ou de droit : « exiger », « réclamer fermement ». Il est souvent suivi de \eng{that} + subjonctif (base verbale) ou d'un nom. L'utiliser pour une invitation ou une demande polie est \textbf{impoli et agressif}.
\begin{itemize}
    \item \textbf{Emploi correct} : \eng{The workers demanded higher wages.} (Les ouvriers ont exigé des salaires plus élevés.) \eng{He demanded that she leave immediately.} (Il a exigé qu'elle parte immédiatement.)
    \item \textbf{Correction} : Pour « demander » (politesse), \eng{ask (for)} (courant), \eng{request} (formel, écrit), \eng{invite} (pour inviter).
\end{itemize}
\begin{center}
\begin{tabularx}{\linewidth}{|X|X|X|}\hline
\rowcolor{poppurpleL} \textbf{Phrase incorrecte (Agressive)} & \textbf{Phrase correcte (Polie)} & \textbf{Contexte} \\ \hline
\eng{I demand you to come.} & \eng{I would like to \textbf{invite} you.} & Invitation. \\ \hline
\eng{We demanded permission.} & \eng{We \textbf{requested} / \textbf{asked for} permission.} & Demande au proviseur. \\ \hline
\eng{He demanded a drink.} & \eng{He \textbf{asked for} a drink.} & Au bar / à la fête. \\ \hline
\end{tabularx}
\end{center}
\end{propriete}

\begin{propriete}[Faux ami 5 : \textbf{an occasion} vs \textit{une opportunité}]
\eng{An occasion} = un moment particulier, un événement marquant (\eng{a special occasion}). \eng{An opportunity} = une chance favorable de faire quelque chose (\eng{a chance}, \eng{an opening}). La confusion est fréquente car en français « occasion » couvre les deux sens.
\begin{itemize}
    \item \eng{The festival was a joyful \textbf{occasion}.} (Le festival fut une \textbf{belle occasion} / un bel événement.)
    \item \eng{This is a great \textbf{opportunity} to learn English.} (C'est une excellente \textbf{opportunité} d'apprendre l'anglais.)
    \item \eng{On this \textbf{occasion}, I would like to thank everyone.} (En cette \textbf{occasion}, je tiens à remercier tous.)
\end{itemize}
\end{propriete}

\courssub{Calques structurels à éradiquer}

Au-delà des mots isolés, les francophones traduisent souvent mot à mot des structures françaises entières. Voici les trois calques les plus fréquents sur ce thème et leurs corrections idiomatiques.

\begin{attention}[Les 3 calques interdits — « Red Flag » pour le correcteur du Bac]
\begin{enumerate}
    \item \textbf{\eng{To make a party}} $\rightarrow$ \textcolor{red}{FAUX}. \eng{Make} ne colle pas avec \eng{party}.
    \begin{itemize}
        \item \textbf{Correct} : \eng{To \textbf{throw} a party} (organiser, familier), \eng{To \textbf{have} a party} (neutre), \eng{To \textbf{organise/host} a party} (formel).
        \item \eng{We are \textbf{throwing} a party for the National Day.}
    \end{itemize}
    \item \textbf{\eng{To assist the festival / the ceremony}} $\rightarrow$ \textcolor{red}{FAUX}. Confusion \eng{assist} / \eng{attend}.
    \begin{itemize}
        \item \textbf{Correct} : \eng{To \textbf{attend} the festival}, \eng{To \textbf{go to} the festival}.
        \item \eng{Did you \textbf{attend} the Ngondo festival last year?}
    \end{itemize}
    \item \textbf{\eng{I demand you to come / I demand you something}} $\rightarrow$ \textcolor{red}{FAUX}. Double erreur : sens (\eng{demand} = exiger) + structure (\eng{demand} ne prend pas l'infinitif \eng{to come}).
    \begin{itemize}
        \item \textbf{Correct (Invitation)} : \eng{I would like to \textbf{invite} you to...}, \eng{We'd love you to \textbf{come}...}
        \item \textbf{Correct (Demande formelle)} : \eng{I \textbf{request} that you attend...}, \eng{I \textbf{ask} you to...}
        \item \eng{The Principal \textbf{requested} that all students \textbf{attend} the ceremony.}
    \end{itemize}
\end{enumerate}
\end{attention}

\courssub{Methodology: how to avoid false friends and calques}

\begin{methode}[Stratégie « Direct English Mapping » pour l'écrit et l'oral]
\emph{Objectif :} Construire la phrase directement en anglais sans passer par la traduction mentale mot à mot du français.

\begin{enumerate}
    \item \textbf{Identifiez l'intention communicative} (inviter, remercier, demander la permission, décrire un événement, s'excuser).
    \item \textbf{Activez le « bloc lexical » (chunk) anglais} associé à cette intention, appris par cœur comme une formule unique.
    \begin{itemize}
        \item Intention « Inviter » $\rightarrow$ Bloc : \eng{invite someone to + V-ing / noun}
        \item Intention « Demander poliment » $\rightarrow$ Bloc : \eng{ask someone for something / request something}
        \item Intention « Être présent » $\rightarrow$ Bloc : \eng{attend an event}
        \item Intention « Aider » $\rightarrow$ Bloc : \eng{assist someone / help someone}
    \end{itemize}
    \item \textbf{Vérifiez la collocation} : le verbe appelle-t-il une préposition ? (\eng{attend} $\emptyset$, \eng{ask for}, \eng{invite to}).
    \item \textbf{Relisez en vous demandant} : « Est-ce que ce mot ressemble à un mot français ? Si oui, ai-je vérifié qu'il a le \emph{même} sens ? »
\end{enumerate}
\end{methode}

\courssub{Word families and pronunciation traps}

Some false friends belong to word families where the noun/adjective/verb shifts meaning. Mastering the family prevents errors in derivation exercises.

\begin{definition}[Familles de mots à surveiller]
\begin{center}
\begin{tabular}{|l|l|l|l|}\hline
\rowcolor{popgreenL} \textbf{Verbe} & \textbf{Nom} & \textbf{Adjectif/Participe} & \textbf{Piège francophone} \\ \hline
\eng{Attend} & \eng{Attendance} (fréquentation) & \eng{Attentive} (attentif) & \eng{Assistance} = aide (pas présence) \\ \hline
\eng{Assist} & \eng{Assistance} (aide, secours) & \eng{Assistant} (adjoint, assistant) & \eng{Assistant} $\neq$ spectateur \\ \hline
\eng{Request} & \eng{A request} (une demande) & \eng{Requested} (demandé) & \eng{Demand} = exigence \\ \hline
\eng{Demand} & \eng{A demand} (une exigence) & \eng{Demanding} (exigeant) & \eng{Demanding} $\neq$ demandant \\ \hline
\eng{Celebrate} & \eng{Celebration} (célébration, fête) & \eng{Celebratory} (festif) & \eng{Ceremony} $\neq$ celebration \\ \hline
\end{tabular}
\end{center}
\end{definition}

\begin{aretenir}[Prononciation : attention aux accents toniques !]
\begin{itemize}
    \item \eng{at\emph{tend}} (accent sur 2\textsuperscript{e} syllabe) — « a-tènd »
    \item \eng{as\emph{sist}} (accent sur 2\textsuperscript{e}) — « a-sist »
    \item \eng{de\emph{mand}} (accent sur 2\textsuperscript{e}) — « di-mande »
    \item \eng{re\emph{quest}} (accent sur 2\textsuperscript{e}) — « ri-kouest »
    \item \eng{\emph{cer}e\emph{mo}ny} (accent sur 1\textsuperscript{re} et 3\textsuperscript{e} secondaire) — « sè-re-mo-ni »
    \item \eng{oc\emph{ca}sion} (accent sur 2\textsuperscript{e}) — « o-ca-jène »
    \item \eng{op\emph{por}tu\emph{ni}ty} (accent sur 3\textsuperscript{e}) — « o-por-tiou-ni-ti »
    \item \eng{eve\emph{ntual}ly} (accent sur 2\textsuperscript{e}) — « i-vèn-tchouè-li »
\end{itemize}
\emph{Astuce :} En anglais, les verbes à 2 syllabes d'origine latine accentuent souvent la 2\textsuperscript{e} syllabe (\eng{at\emph{tend}}, \eng{as\emph{sist}}, \eng{de\emph{mand}}, \eng{re\emph{quest}}), alors que les noms correspondants accentuent la 1\textsuperscript{re} (\eng{\emph{at}tendance}, \eng{\emph{as}}sistant, \eng{\emph{de}}mand, \eng{\emph{re}}quest — bien que \eng{request} garde l'accent sur la 2\textsuperscript{e} comme nom).
\end{aretenir}

\courssub{Resolved exercise: spotting and correcting errors}

\begin{exoresolu}[Correction d'un paragraphe d'élève — 8 erreurs à trouver]
\textbf{Énoncé :} Corrigez les erreurs lexicales (faux amis, calques) dans le paragraphe suivant écrit par un élève de 1ère A. Réécrivez la version correcte.

\eng{``Last weekend, my school organised a big ceremony for the National Day. Many parents and neighbours assisted the ceremony. The Principal demanded the students to come early. Eventually, the ceremony was very beautiful. We made a big party after the official speeches. It was a good opportunity to meet my old friends. I assisted my younger brother to find a seat. The Principal demanded a minute of silence for the ancestors.''}

\tcblower
\textbf{Solution détaillée et commentée :}

\begin{enumerate}
    \item \eng{organised a big \textbf{ceremony}} $\rightarrow$ \eng{organised a big \textbf{celebration} / \textbf{festival}}. \emph{Contexte : fête nationale, ambiance festive.}
    \item \eng{Many parents and neighbours \textbf{assisted} the ceremony} $\rightarrow$ \eng{\textbf{attended} the celebration}. \emph{Faux ami \#1.}
    \item \eng{The Principal \textbf{demanded} the students to come early} $\rightarrow$ \eng{The Principal \textbf{asked/requested} the students to come early}. \emph{Faux ami \#4 + structure. \eng{Demand} ne prend pas \eng{to + infinitif}.}
    \item \eng{\textbf{Eventually}, the ceremony was very beautiful} $\rightarrow$ \eng{\textbf{In the end} / \textbf{Finally}, the celebration was very beautiful}. \emph{Faux ami \#2 : \eng{Eventually} implique une difficulté ou une attente longue. Ici, simple conclusion.}
    \item \eng{We \textbf{made} a big party} $\rightarrow$ \eng{We \textbf{threw} / \textbf{had} a big party}. \emph{Calque \#1.}
    \item \eng{It was a good \textbf{opportunity} to meet...} $\rightarrow$ \textbf{Correct}. \eng{Opportunity} = chance favorable. Bien utilisé.
    \item \eng{I \textbf{assisted} my younger brother} $\rightarrow$ \eng{I \textbf{helped} / \textbf{assisted} my younger brother}. \emph{Ici \eng{assist} EST correct (sens « aider »), mais \eng{helped} est plus naturel à l'oral.}
    \item \eng{The Principal \textbf{demanded} a minute of silence} $\rightarrow$ \eng{The Principal \textbf{requested} / \textbf{called for} / \textbf{observed} a minute of silence}. \emph{Faux ami \#4. \eng{Demanded} sonne comme un ordre autoritaire. \eng{Observed a minute of silence} est l'expression idiomatique standard.}
\end{enumerate}

\textbf{Version corrigée :}
\eng{``Last weekend, my school organised a big \textbf{celebration} for the National Day. Many parents and neighbours \textbf{attended} the \textbf{celebration}. The Principal \textbf{asked} the students to come early. \textbf{In the end}, the \textbf{celebration} was very beautiful. We \textbf{had} a big party after the official speeches. It was a good opportunity to meet my old friends. I \textbf{helped} my younger brother to find a seat. The Principal \textbf{called for} a minute of silence for the ancestors.''}
\end{exoresolu}

\courssub{Practice exercises}

\begin{exercice}[Choix du mot juste — Complétez avec le mot approprié]
Choisissez entre les deux mots proposés (attention aux faux amis) et mettez le verbe à la forme correcte si nécessaire.
\begin{enumerate}
    \item Over 500 people \_\_\_\_\_\_\_\_\_\_ (attended / assisted) the cultural festival in Buea last Saturday.
    \item The organisers \_\_\_\_\_\_\_\_\_\_ (requested / demanded) permission from the Divisional Officer two weeks ago.
    \item \_\_\_\_\_\_\_\_\_\_ (Eventually / Possibly), we might organise a second event if the weather is bad.
    \item The graduation \_\_\_\_\_\_\_\_\_\_ (ceremony / celebration) lasts two hours, but the \_\_\_\_\_\_\_\_\_\_ (ceremony / celebration) continues all night.
    \item Please don't \_\_\_\_\_\_\_\_\_\_ (make / throw) too much noise; the neighbours are sleeping.
    \item I would like to \_\_\_\_\_\_\_\_\_\_ (invite / demand) you to our end-of-year party.
    \item Winning the cup was a great \_\_\_\_\_\_\_\_\_\_ (occasion / opportunity) for the team to travel abroad.
    \item The volunteers \_\_\_\_\_\_\_\_\_\_ (attended / assisted) the elderly people during the march.
\end{enumerate}

\end{exercice}

\begin{exercice}[Traduction piège — Évitez les calques]
Traduisez en anglais naturel (pas mot à mot).
\begin{enumerate}
    \item Beaucoup d'anciens élèves ont assisté à la cérémonie de remise des prix.
    \item Le proviseur a demandé aux élèves de respecter le règlement.
    \item Nous avons organisé une grande fête pour la fête de la Jeunesse.
    \item Éventuellement, nous pourrons inviter le maire.
    \item C'était une occasion unique de danser le makossa.
    \item J'ai aidé ma mère à préparer le ndolé pour la fête.
    \item Finalement, tout le monde s'est bien amusé.
    \item Les élèves ont exigé de meilleurs repas à la cantine. (Ici, \eng{demand} est correct !)
\end{enumerate}
\end{exercice}

\courssub{Cultural insight and exam tip}

\begin{savaistu}[Le faux ami au Baccalauréat : erreur lexicale sanctionnée]
Au Cameroun, comme dans tous les pays francophones évalués sur l'anglais langue étrangère, l'utilisation d'un faux ami dans une production écrite (\emph{Composition}, \emph{Guided Writing}, \emph{Summary}) est classée comme une \textbf{erreur lexicale majeure} (Lexical Error). Elle fait perdre des points sur le critère « Vocabulary Range and Accuracy ».
\begin{itemize}
    \item \eng{``The students assisted the match''} au lieu de \eng{attended} $\rightarrow$ \textbf{0 point} pour le mot, confusion de sens.
    \item \eng{``I demand you to come''} $\rightarrow$ \textbf{Erreur de registre + structure} : double pénalité.
\end{itemize}
\emph{Stratégie gagnante :} Apprenez les paires \textbf{Faux ami / Vrai équivalent} comme du vocabulaire à part entière. Révisez ce tableau la veille de l'épreuve. Chaque correction évitée est un point gagné sur 20.
\end{savaistu}

\begin{corrige}[Exercice 1 — Choix du mot juste]
\begin{enumerate}
    \item \textbf{attended} (assister à = être présent). \eng{Assisted} = aidé.
    \item \textbf{requested} (demander poliment/formellement). \eng{Demanded} = exigé.
    \item \textbf{Possibly} (éventuellement = si possible). \eng{Eventually} = finalement.
    \item \textbf{ceremony} (partie officielle) / \textbf{celebration} (partie festive).
    \item \textbf{make} (collocation : \eng{make noise} — ici \eng{make} est correct ! Piège inversé : \eng{make a party} est faux, mais \eng{make noise} est juste).
    \item \textbf{invite} (inviter). \eng{Demand} = exiger.
    \item \textbf{opportunity} (chance favorable). \eng{Occasion} = événement/moment.
    \item \textbf{assisted} (aidé). Ici le sens est « porter assistance », donc \eng{assist} est le bon verbe.
\end{enumerate}
\end{corrige}

\begin{corrige}[Exercice 2 — Traduction piège]
\begin{enumerate}
    \item Many former students \textbf{attended} the prize-giving \textbf{ceremony}.
    \item The Principal \textbf{asked} / \textbf{requested} the students to respect the rules.
    \item We \textbf{threw} / \textbf{had} a big \textbf{party} for Youth Day.
    \item \textbf{Possibly} / \textbf{If necessary}, we could invite the Mayor.
    \item It was a unique \textbf{opportunity} to dance makossa.
    \item I \textbf{helped} / \textbf{assisted} my mother prepare ndolé for the party.
    \item \textbf{Eventually} / \textbf{In the end}, everyone enjoyed themselves / had a good time.
    \item The students \textbf{demanded} better meals at the canteen. (\emph{Correct : sens d'exigence/revendication}).
\end{enumerate}
\end{corrige}

\coursec{Familles de mots, prononciation et orthographe}

\courssub{Transition et aperçu}

Dans la section précédente, nous avons identifié les \cle{faux amis} et les pièges de traduction qui guettent les francophones lorsqu'ils parlent de fêtes et de cérémonies. Nous savons désormais distinguer \eng{celebration} de \eng{celebrity}, ou \eng{festival} de \eng{feast}. Pour aller plus loin et gagner en aisance — tant à l'oral qu'à l'écrit — il est indispensable de maîtriser la \cle{dérivation} : la capacité à construire, à partir d'une racine, le verbe, le nom, l'adjectif, voire l'adverbe ou le nom d'agent. C'est l'objet de cette section. Nous observerons cinq familles lexicales centrales pour le thème des célébrations communautaires, nous fixerons les règles de prononciation et d'orthographe (notamment le \cle{-our} britannique vs \cle{-or} américain), et nous mettrons tout en pratique dans des phrases authentiques.

\courssub{La famille de \eng{celebrate} : verbe, nom, nom d'agent, adjectif}

Commençons par observer un court extrait d'un article de presse camerounais fictif décrivant la fête de la jeunesse à Buea.

\begin{textebox}{Youth Day in Buea: a community celebration}
Last February, the \textbf{celebration} of Youth Day brought the whole \textbf{neighbourhood} together. Students from Government High School \textbf{celebrated} the event with traditional dances and songs. The \textbf{celebrant} of the mass, Father Peter, reminded the youth of their civic duties. It was a truly \textbf{celebrated} occasion, reported even in the national newspapers. The \textbf{festivities} lasted three days, creating a \textbf{festive} atmosphere in every street.
\end{textebox}

\begin{definition}[Glossaire]
\textbf{celebration} (n.) : célébration, fête \quad \textbf{celebrated} (v. past / adj.) : a célébré / célèbre \quad \textbf{celebrant} (n.) : célébrant (celui qui officie ou qui fête) \quad \textbf{festivities} (n. pl.) : festivités, réjouissances \quad \textbf{festive} (adj.) : festif, joyeux \quad \textbf{civic duties} (n. pl.) : devoirs civiques
\end{definition}

\begin{propriete}[Famille de \eng{celebrate} — Formes et emplois]
\begin{center}
\begin{tabularx}{\linewidth}{|X|X|X|X|}
\hline
\rowcolor{popblueL}
\textbf{Forme} & \textbf{Nature} & \textbf{Sens principal} & \textbf{Exemple traduit} \\
\hline
\eng{celebrate} & verbe (régulier) & célébrer, fêter & \eng{We celebrate Christmas on 25th December.} (Nous fêtons Noël le 25 décembre.) \\
\eng{celebration} & nom (comptable) & célébration, fête & \eng{The celebration was a great success.} (La fête a été un grand succès.) \\
\eng{celebrant} & nom (comptable) & célébrant, participant & \eng{The celebrants wore traditional attire.} (Les participants portaient des tenues traditionnelles.) \\
\eng{celebrated} & adj. (participe passé) & célèbre, renommé & \eng{She is a celebrated author.} (C'est une auteure célèbre.) \\
\hline
\end{tabularx}
\end{center}
\emph{Remarque :} \eng{celebrated} comme adjectif ne signifie pas « célébré » (passif) mais « célèbre, illustre ». C'est un \cle{faux ami morphologique} fréquent.
\end{propriete}

\begin{attention}[Piège : \eng{celebrated} (adj.) $\neq$ « célébré »]
Ne confondez pas :
\begin{itemize}
\item \eng{The hero was celebrated by the crowd.} (verbe passif : Le héros a été acclamé par la foule.)
\item \eng{A celebrated hero.} (adjectif : Un héros \textbf{célèbre / illustre}.)
\end{itemize}
En français, on dirait « un héros célébré » pour le sens passif, mais « un héros célèbre » pour l'adjectif. En anglais, l'écrit est identique, seul le contexte (ou l'absence d'agent \eng{by}) disambiguïse.
\end{attention}

\begin{exemplebox}[Autres dérivés utiles]
\begin{itemize}
\item \eng{celebratory} (adj.) : qui marque une célébration — \eng{A celebratory dinner.} (Un dîner de célébration.)
\item \eng{celebrator} (n.) : celui qui fête (synonyme de \eng{celebrant} mais moins formel).
\end{itemize}
\end{exemplebox}

\begin{popfigure}
\begin{tikzpicture}[
  mindmap,
  grow cyclic,
  every node/.style={concept, execute at begin node=\hskip0pt, align=center},
  root concept/.append style={concept color=popblue, minimum size=2.5cm, text=white, font=\bfseries\large},
  level 1 concept/.append style={concept color=popblue!70, sibling angle=90, minimum size=2cm, text=white, font=\bfseries\normalsize},
  level 2 concept/.append style={concept color=popgreen!70, sibling angle=45, minimum size=1.8cm, text=white, font=\small},
  level 3/.style={sibling angle=30}
]
\node[root concept] {celebr-}
  child[clockwise from=30] { node[level 1 concept, align=center]{celebrate (v.)\\célébrer}
    child { node[level 2 concept, align=center]{celebrated (adj.)\\célèbre} }
    child { node[level 2 concept, align=center]{celebrating (part.)\\en train de fêter} }
  }
  child { node[level 1 concept, align=center]{celebration (n.)\\célébration}
    child { node[level 2 concept, align=center]{celebrations (pl.)\\fêtes} }
  }
  child { node[level 1 concept, align=center]{celebrant (n.)\\célébrant / participant}
    child { node[level 2 concept] {celebrants (pl.)} }
  }
  child { node[level 1 concept, align=center]{celebratory (adj.)\\de célébration} };
\end{tikzpicture}
\legende{Arbre dérivé de la racine \eng{celebr-} : verbe, noms, adjectifs et participes.}
\end{popfigure}

\courssub{La famille de \eng{festival} : nom, adjectif, pluriel festivities}

\begin{textebox}{The Ngondo Festival}
The \textbf{festival} of Ngondo is a major cultural event in Douala. During the \textbf{festivities}, the Sawa people gather on the banks of the Wouri River. The atmosphere is incredibly \textbf{festive}: colourful \textbf{outfits}, drumming, and traditional canoes racing on the water. It is a moment of unity and joy for the whole \textbf{community}.
\end{textebox}

\begin{definition}[Glossaire]
\textbf{festival} (n.) : festival \quad \textbf{festivities} (n. pl.) : festivités, réjouissances \quad \textbf{festive} (adj.) : festif, joyeux \quad \textbf{outfits} (n. pl.) : tenues \quad \textbf{Sawa people} : peuple Sawa \quad \textbf{banks} (n. pl.) : rives \quad \textbf{canoes} (n. pl.) : pirogues
\end{definition}

\begin{propriete}[Famille de \eng{festival} — Formes et emplois]
\begin{center}
\begin{tabularx}{\linewidth}{|X|X|X|X|}
\hline
\rowcolor{poporangeL}
\textbf{Forme} & \textbf{Nature} & \textbf{Sens principal} & \textbf{Exemple traduit} \\
\hline
\eng{festival} & nom (comptable) & festival, fête organisée & \eng{The Yaoundé Jazz Festival attracts many tourists.} (Le Festival de Jazz de Yaoundé attire beaucoup de touristes.) \\
\eng{festive} & adjectif & festif, joyeux, de fête & \eng{The town looks festive with all these lights.} (La ville a l'air festive avec toutes ces lumières.) \\
\eng{festivity} & nom (comptable) & réjouissance, fête (rare au singulier) & \eng{The program includes various festivities.} (Le programme comprend diverses réjouissances.) \\
\eng{festivities} & nom (pluriel usuel) & festivités, ensemble des réjouissances & \eng{The festivities kick off at 8 p.m.} (Les festivités commencent à 20 h.) \\
\hline
\end{tabularx}
\end{center}
\emph{Point de grammaire :} \eng{festivity} s'emploie presque toujours au pluriel \eng{festivities} pour désigner l'ensemble des animations d'un événement.
\end{propriete}

\begin{exemplebox}[Collocations avec \eng{festival / festive}]
\begin{itemize}
\item \eng{music festival / cultural festival / film festival}
\item \eng{festive season / festive atmosphere / festive mood / festive spirit}
\item \eng{to get into the festive spirit} (entrer dans l'esprit de la fête)
\end{itemize}
\end{exemplebox}

\courssub{La famille de \eng{invite} : verbe, nom, adjectif, invitee/guest}

\begin{textebox}{Invitation card for a wedding}
Mr and Mrs Nkeng have the pleasure of \textbf{inviting} you to the wedding \textbf{celebration} of their daughter, Alice. The ceremony will take place at the Presbyterian Church, Bastos, on Saturday 15th June at 10 a.m. An \textbf{inviting} reception will follow at the City Hall. Kindly RSVP by 1st June. We look forward to welcoming every \textbf{invitee} (or \textbf{guest}) personally.
\end{textebox}

\begin{definition}[Glossaire]
\textbf{inviting} (part. présent / adj.) : invitant / attrayant, accueillant \quad \textbf{invitation} (n.) : invitation \quad \textbf{invitee} (n.) : invité (personne conviée) \quad \textbf{guest} (n.) : invité, hôte \quad \textbf{RSVP} (abbr. fr.) : \emph{Répondez s'il vous plaît} \quad \textbf{Kindly} (adv.) : veuillez, aimablement
\end{definition}

\begin{propriete}[Famille de \eng{invite} — Formes, prononciation et emplois]
\begin{center}
\begin{tabularx}{\linewidth}{|X|X|X|X|X|}
\hline
\rowcolor{popgreenL}
\textbf{Forme} & \textbf{Nature} & \textbf{Prononciation (approx. fr.)} & \textbf{Sens} & \textbf{Exemple traduit} \\
\hline
\eng{invite} & verbe & « in-vaïte » & inviter & \eng{I invite you to my party.} (Je t'invite à ma fête.) \\
\eng{invitation} & nom & « in-vi-ta-ssion » & invitation & \eng{Did you receive the invitation?} (As-tu reçu l'invitation ?) \\
\eng{inviting} & adj. / part. & « in-vaï-ting » & attrayant, accueillant & \eng{An inviting smell of grilled fish.} (Une odeur alléchante de poisson grillé.) \\
\eng{invitee} & nom & « in-vi-taï » & l'invité (celui qui reçoit l'invitation) & \eng{All invitees must show their card.} (Tous les invités doivent montrer leur carte.) \\
\eng{guest} & nom & « gueste » & l'invité, l'hôte & \eng{We have guests for dinner.} (Nous avons des invités pour le dîner.) \\
\hline
\end{tabularx}
\end{center}
\emph{Prononciation clé :} Le verbe \eng{invite} se termine par le son /aɪt/ (« aïte »), alors que le nom \eng{invitation} porte l'accent sur la troisième syllabe et se termine par /ʃən/ (« ssion » → « cheune » en français). Ne dites pas *« in-vi-te » pour le verbe.
\end{propriete}

\begin{attention}[Faux ami / Calque : \eng{invité} $\neq$ \eng{invited} (seul)]
\begin{itemize}
\item \eng{The invited guests} (Les invités) : \eng{invited} est un adjectif (participe passé) qui qualifie \eng{guests}.
\item \eng{The invitees} (Les invités) : nom d'agent en \cle{-ee} (comme \eng{employee}, \eng{trainee}).
\item \textbf{Erreur fréquente :} *« The inviteds ». \textbf{Interdit.} On dit \eng{the guests}, \eng{the invitees} ou \eng{the invited guests}.
\end{itemize}
\end{attention}

\courssub{La famille de \eng{neighbour} : orthographe UK/US, neighbourhood, neighbouring}

\begin{textebox}{Neighbourhood life in Yaoundé}
In my \textbf{neighbourhood} in Yaoundé, \textbf{neighbours} know each other well. During the end-of-year \textbf{festivities}, they organise a big street party. The \textbf{neighbouring} streets are also invited. It is a tradition that strengthens community bonds. Note the British spelling: \textbf{neighbour} with a \textbf{u}.
\end{textebox}

\begin{definition}[Glossaire]
\textbf{neighbourhood} (n.) : quartier, voisinage \quad \textbf{neighbouring} (adj.) : voisin, adjacent \quad \textbf{bonds} (n. pl.) : liens \quad \textbf{strengthens} (v.) : renforce
\end{definition}

\begin{propriete}[Orthographe britannique vs américaine : le \cle{-our} / \cle{-or}]
\begin{center}
\begin{tabularx}{\linewidth}{|X|X|X|}
\hline
\rowcolor{poppurpleL}
\textbf{Britannique (UK / Cameroun / Commonwealth)} & \textbf{Américain (US)} & \textbf{Français} \\
\hline
\eng{neighbour} & \eng{neighbor} & voisin \\
\eng{neighbourhood} & \eng{neighborhood} & quartier \\
\eng{neighbourly} (adj.) & \eng{neighborly} & aimable, bon voisin \\
\eng{honour} & \eng{honor} & honneur \\
\eng{favourite} & \eng{favorite} & favori, préféré \\
\eng{colour} & \eng{color} & couleur \\
\hline
\end{tabularx}
\end{center}
\textbf{Règle pour le Bac et les examens officiels au Cameroun :} Les deux orthographes sont acceptées, \textbf{mais vous devez être cohérent} dans toute votre copie. Si vous choisissez l'orthographe britannique (recommandée dans le système éducatif camerounais), gardez-la pour \eng{neighbour}, \eng{colour}, \eng{honour}, \eng{favourite}, etc.
\end{propriete}

\begin{exemplebox}[Collocations avec \eng{neighbour / neighbourhood}]
\begin{itemize}
\item \eng{next-door neighbour} (voisin d'à côté / de palier)
\item \eng{neighbourhood watch} (comité de vigilance de quartier)
\item \eng{friendly neighbourhood} (quartier convivial)
\item \eng{neighbouring countries / villages / schools} (pays / villages / écoles voisins)
\end{itemize}
\end{exemplebox}

\courssub{La famille de \eng{organise} : verbe, nom, nom d'agent, adjectif}

\begin{textebox}{Organising the School Cultural Week}
The \textbf{organisation} of the Cultural Week is led by a committee of teachers and students. The main \textbf{organiser}, Mr. Atangana, has \textbf{organised} everything down to the last detail. The programme is well \textbf{organised}: dance competitions, poetry slams, and a grand finale. Everyone appreciates his \textbf{organisational} skills.
\end{textebox}

\begin{definition}[Glossaire]
\textbf{organisation} (n.) : organisation \quad \textbf{organiser} (n.) : organisateur \quad \textbf{organised} (adj.) : organisé, ordonné \quad \textbf{organisational} (adj.) : organisationnel \quad \textbf{down to the last detail} (expr.) : dans les moindres détails \quad \textbf{poetry slam} : concours de slam / poésie orale
\end{definition}

\begin{propriete}[Famille de \eng{organise} — Dérivation en \cle{-tion} / \cle{-sation}]
\begin{center}
\begin{tabularx}{\linewidth}{|X|X|X|X|}
\hline
\rowcolor{poppinkL}
\textbf{Verbe} & \textbf{Nom en \eng{-tion / -sation}} & \textbf{Nom d'agent (\eng{-r / -er})} & \textbf{Adjectif (part. passé / \eng{-al})} \\
\hline
\eng{organise} & \eng{organisation} & \eng{organiser} & \eng{organised} / \eng{organisational} \\
\eng{invite} & \eng{invitation} & \eng{inviter} (rare) & \eng{invited} / \eng{inviting} \\
\eng{celebrate} & \eng{celebration} & \eng{celebrant} & \eng{celebrated} / \eng{celebratory} \\
\eng{decorate} & \eng{decoration} & \eng{decorator} & \eng{decorated} / \eng{decorative} \\
\eng{prepare} & \eng{preparation} & \eng{preparer} (rare) & \eng{prepared} / \eng{preparatory} \\
\hline
\end{tabularx}
\end{center}
\emph{Règle de prononciation :} Tous les noms en \cle{-tion} / \cle{-sation} se terminent par le son /ʃən/ (« cheune » en approximation française) : \eng{organisation} → « or-ga-ni-za-cheune », \eng{celebration} → « sé-lé-bra-cheune », \eng{invitation} → « in-vi-ta-cheune ». L'accent tombe sur la syllabe précédant \eng{-tion}.
\end{propriete}

\begin{methode}[Enrichir son vocabulaire par la dérivation (5 étapes)]
\begin{enumerate}
\item \textbf{Identifiez} un mot-clé du thème (ex. \eng{decorate}).
\item \textbf{Dérivez} systématiquement : verbe → nom (\eng{-tion}) → nom d'agent (\eng{-r/-or}) → adjectif (\eng{-ed/-ing/-al/-ive}).
\item \textbf{Vérifiez} la prononciation (accent tonique, son \eng{-tion}) et l'orthographe (UK/US).
\item \textbf{Rédigez} une phrase personnelle pour \textbf{chaque forme} (contexte : fête, école, quartier).
\item \textbf{Révisez} par familles (carte mentale) avant l'évaluation.
\end{enumerate}
\end{methode}

\begin{savaistu}[Le suffixe \eng{-ee} : le destinataire de l'action]
En anglais, le suffixe \eng{-ee} (venant du français passé composé \emph{-é}) désigne \textbf{la personne qui subit ou reçoit l'action} :
\begin{itemize}
\item \eng{invite} $\rightarrow$ \eng{invitee} (celui qui est invité)
\item \eng{employ} $\rightarrow$ \eng{employee} (celui qui est employé = salarié)
\item \eng{train} $\rightarrow$ \eng{trainee} (stagiaire)
\item \eng{interview} $\rightarrow$ \eng{interviewee} (la personne interviewée)
\item \eng{referee} (arbitre / personne qui donne une référence) — attention, sens spécialisé.
\end{itemize}
C'est un outil puissant pour former des noms d'agent \textbf{passifs} sans utiliser de phrase relative.
\end{savaistu}

\courssub{Atelier prononciation et orthographe}

\begin{propriete}[Récapitulatif : sons difficiles et pièges orthographiques]
\begin{center}
\begin{tabularx}{\linewidth}{|X|X|X|}
\hline
\rowcolor{popgoldL}
\textbf{Mot / Famille} & \textbf{Piège / Difficulté} & \textbf{Conseil / Astuce} \\
\hline
\eng{celebration} & Prononciation de \eng{-tion} & Dites « sé-lé-bra-\textbf{cheune} », jamais « sé-lé-bra-\textbf{tion} ». \\
\eng{invitation} & Accent tonique & Sur la syllabe avant \eng{-tion} : « in-vi-\textbf{TA}-cheune ». \\
\eng{invite} (v.) vs \eng{invitation} (n.) & Voyelle finale & Verbe = /aɪ/ (« aï ») ; Nom = /eɪ/ (« é ») dans la 2e syllabe. \\
\eng{neighbour} & \textbf{u} muet après \textbf{g} & UK : \eng{neighbour} (8 lettres) ; US : \eng{neighbor} (7 lettres). Mnémo : « \textbf{U} are my neighbour » (UK). \\
\eng{festival} / \eng{festive} & Accent & \eng{FES-ti-val} (1re) vs \eng{FES-tive} (1re). Pas d'accent sur \eng{-val}. \\
\eng{organised} & \textbf{s} vs \textbf{z} & UK : \eng{organi\textbf{s}ed} ; US : \eng{organi\textbf{z}ed}. Cohérence exigée. \\
\hline
\end{tabularx}
\end{center}
\end{propriete}

\begin{exoresolu}[Exercice résolu : Dérivation et choix de forme]
\textbf{Énoncé :} Complétez les phrases avec la forme correcte du mot entre parenthèses.
\begin{enumerate}
\item The \_\_\_\_\_\_\_\_\_\_ of the national day parade was impressive. (organise)
\item My \_\_\_\_\_\_\_\_\_\_ baked a delicious cake for the street party. (neighbour)
\item We received a formal \_\_\_\_\_\_\_\_\_\_ for the wedding. (invite)
\item The atmosphere in the stadium was very \_\_\_\_\_\_\_\_\_\_. (festival)
\item He is a \_\_\_\_\_\_\_\_\_\_ musician known across Africa. (celebrate)
\item The \_\_\_\_\_\_\_\_\_\_ lasted until dawn. (festivity)
\item She has excellent \_\_\_\_\_\_\_\_\_\_ skills. (organise)
\end{enumerate}

\tcblower
\textbf{Solution détaillée :}
\begin{enumerate}
\item \eng{organisation} (nom sujet de \eng{was}) — \emph{L'organisation du défilé...}
\item \eng{neighbour} (nom singulier après \eng{My}) — \emph{Mon voisin...} (orthographe UK)
\item \eng{invitation} (nom après \eng{a formal}) — \emph{Une invitation formelle...}
\item \eng{festive} (adjectif après \eng{very}) — \emph{L'ambiance était très festive.}
\item \eng{celebrated} (adjectif attribut du sujet \eng{He}) — \emph{Il est un musicien célèbre.}
\item \eng{festivities} (nom pluriel sujet de \eng{lasted}) — \emph{Les festivités ont duré...}
\item \eng{organisational} (adjectif qualificatif de \eng{skills}) — \emph{Des compétences organisationnelles.}
\end{enumerate}
\end{exoresolu}

\courssub{Mise en contexte : rapport sur un festival communautaire}

Lisez le texte suivant, qui réemploie l'essentiel du vocabulaire de cette section. Repérez les formes dérivées, notez l'orthographe britannique, puis répondez aux questions.

\begin{textebox}{Report: The Mfoundi Community Festival}
Last weekend, the Mfoundi \textbf{neighbourhood} in Yaoundé hosted its annual cultural \textbf{festival}. The \textbf{organisation} had started three months earlier, led by a dynamic committee. The main \textbf{organiser}, Mrs. Essomba, had \textbf{invited} all residents, local \textbf{authorities}, and the \textbf{neighbouring} school principals. The \textbf{invitation} cards, designed by students, were very \textbf{inviting} with their bright colours.

On Saturday morning, the \textbf{celebration} began with a parade. \textbf{Celebrants} in traditional \eng{toghu} and \eng{ndop} fabrics danced to the rhythm of drums. The \textbf{festive} mood was contagious. The Divisional Officer, a \textbf{celebrated} administrator known for his support of youth projects, gave a short speech. He praised the \textbf{neighbourly} spirit of the community.

The \textbf{festivities} continued with a football tournament, a \eng{suya} grilling contest, and a talent show. Everything was perfectly \textbf{organised}. As the sun set, the \textbf{invitees} (or \textbf{guests}) shared a communal meal. It was a day of unity, joy, and shared \textbf{celebration}.
\end{textebox}

\begin{definition}[Glossaire complémentaire]
\textbf{hosted} (v.) : a accueilli / organisé \quad \textbf{annual} (adj.) : annuel \quad \textbf{dynamic} (adj.) : dynamique \quad \textbf{residents} (n. pl.) : habitants \quad \textbf{authorities} (n. pl.) : autorités \quad \textbf{principals} (n. pl.) : proviseurs / directeurs \quad \textbf{bright colours} : couleurs vives \quad \textbf{parade} (n.) : défilé \quad \textbf{toghu / ndop} : tissus traditionnels du Nord-Ouest \quad \textbf{contagious} (adj.) : contagieux \quad \textbf{Divisional Officer} : Sous-préfet \quad \textbf{praised} (v.) : a loué / félicité \quad \textbf{neighbourly} (adj.) : bon voisin, amical \quad \textbf{grilling contest} : concours de grillade \quad \textbf{communal meal} : repas communautaire \quad \textbf{shared} (adj.) : partagé
\end{definition}

\begin{exercice}[Questions de compréhension et de lexique]
\begin{enumerate}
\item Relevez dans le texte \textbf{sept} mots appartenant à la famille de \eng{organise}, \eng{invite}, \eng{celebrate}, \eng{festival} ou \eng{neighbour}. Indiquez pour chacun : la famille, la nature grammaticale et la traduction.
\item Trouvez dans le texte l'équivalent anglais de : « autorités locales », « tissus traditionnels », « concours de grillade », « repas communautaire ».
\item Pourquoi l'auteur écrit-il \eng{neighbourhood} et \eng{neighbourly} avec un \textbf{u} ? Quelle serait l'orthographe américaine ?
\item Transformez : \eng{The organisation was perfect.} $\rightarrow$ \eng{They \_\_\_\_\_\_\_\_\_\_ everything perfectly.} (verbe) \\
\eng{She is a celebrated artist.} $\rightarrow$ \eng{People \_\_\_\_\_\_\_\_\_\_ her art.} (verbe)
\item \textbf{Production écrite (30-40 mots) :} Imaginez que vous êtes le/la secrétaire du comité d'organisation. Écrivez un court message WhatsApp (en anglais) aux membres pour les remercier après le festival. Utilisez au moins 5 mots dérivés vus dans cette section.
\end{enumerate}
\end{exercice}

\begin{corrige}[Exercice : Report — Mfoundi Community Festival]
\begin{enumerate}
\item \textbf{Familles dérivées dans le texte :}
\begin{itemize}
\item \eng{organisation} (famille \eng{organise}, nom, organisation)
\item \eng{organiser} (famille \eng{organise}, nom d'agent, organisateur)
\item \eng{invited} (famille \eng{invite}, verbe passé / part. passé, a invité)
\item \eng{invitation} (famille \eng{invite}, nom, invitation)
\item \eng{inviting} (famille \eng{invite}, adj., attrayant)
\item \eng{celebration} (famille \eng{celebrate}, nom, célébration/fête — $\times$2 dans le texte)
\item \eng{celebrants} (famille \eng{celebrate}, nom pl., participants/célébrants)
\item \eng{festive} (famille \eng{festival}, adj., festif)
\item \eng{festivities} (famille \eng{festival}, nom pl., festivités)
\item \eng{celebrated} (famille \eng{celebrate}, adj., célèbre)
\item \eng{neighbourhood} (famille \eng{neighbour}, nom, quartier)
\item \eng{neighbouring} (famille \eng{neighbour}, adj., voisin)
\item \eng{neighbourly} (famille \eng{neighbour}, adj., bon voisin/amical)
\item \eng{organised} (famille \eng{organise}, adj., bien organisé)
\item \eng{invitees} (famille \eng{invite}, nom pl., invités)
\end{itemize}
\item \textbf{Équivalents :} local authorities / traditional fabrics (toghu and ndop) / grilling contest / communal meal.
\item \textbf{Orthographe :} Le \textbf{u} marque l'orthographe britannique (standard au Cameroun, au Royaume-Uni, dans le Commonwealth). US : \eng{neighborhood}, \eng{neighborly}.
\item \textbf{Transformation :}
\begin{itemize}
\item \eng{They \textbf{organised} everything perfectly.}
\item \eng{People \textbf{celebrate} her art.} (ou \eng{People \textbf{celebrated} her art.} au passé)
\end{itemize}
\item \textbf{Production écrite (modèle) :}
\eng{“Dear committee members, thank you for the amazing organisation of our festival! The celebration was a huge success. Your organisational skills made the festivities unforgettable. Special thanks to the organisers and all invitees who joined. Let's keep this neighbourly spirit alive! Mrs. Essomba”}
\end{enumerate}
\end{corrige}

\begin{aretenir}[Synthèse : les 5 familles clés pour “Community Celebrations”]
\begin{center}
\begin{tabularx}{\linewidth}{|X|X|X|X|X|}
\hline
\rowcolor{popblueL}
\textbf{Racine} & \textbf{Verbe} & \textbf{Nom (chose/événement)} & \textbf{Nom (personne)} & \textbf{Adjectif(s) clé(s)} \\
\hline
\eng{celebr-} & \eng{celebrate} & \eng{celebration} & \eng{celebrant} & \eng{celebrated} (célèbre), \eng{celebratory} \\
\eng{festiv-} & — & \eng{festival}, \eng{festivities} & — & \eng{festive} \\
\eng{invit-} & \eng{invite} & \eng{invitation} & \eng{invitee}, \eng{guest} & \eng{inviting} (attrayant), \eng{invited} \\
\eng{neighbour-} & — & \eng{neighbour}, \eng{neighbourhood} & \eng{neighbour} & \eng{neighbouring}, \eng{neighbourly} \\
\eng{organis-} & \eng{organise} & \eng{organisation} & \eng{organiser} & \eng{organised}, \eng{organisational} \\
\hline
\end{tabularx}
\end{center}
\textbf{Rappels sonores :} \eng{-tion} = « cheune » | \eng{invite} (v.) = « aïte » | \eng{neighbour} (UK) = avec \textbf{u} | Accent sur syllabe avant \eng{-tion}.
\end{aretenir}

\coursec{Community celebrations and festivals: core vocabulary}

Cette leçon ouvre le \textbf{Chapter 1 : Family and Social Life}. Nous abordons le lexique indispensable pour \cle{comprendre}, \cle{parler} et \cle{écrire} au sujet des célébrations communautaires (fêtes de fin d'année, fêtes traditionnelles, récoltes, journées culturelles scolaires) et pour \cle{interagir} avec les différents acteurs : camarades (\textit{mates}), voisins (\textit{neighbours}), aînés (\textit{elders}) et autorités scolaires (\textit{school authorities}). L'objectif est de maîtriser les \cle{champs lexicaux}, les \cle{collocations} (associations de mots naturelles), les \cle{expressions idiomatiques} et d'éviter les \cle{faux amis} fréquents chez les francophones.

\begin{definition}[Champ lexical : Les célébrations communautaires]
Le vocabulaire se regroupe en catégories fonctionnelles. Mémorisez les mots \textbf{par blocs thématiques} (collocations incluses) et non en listes isolées.
\end{definition}

\begin{center}
\begin{tabularx}{\linewidth}{|l|X|X|}\hline
\rowcolor{popblueL} \textbf{Catégorie} & \textbf{Mots-clés \& Collocations (UK English)} & \textbf{Traduction / Contexte} \\ \hline
\textbf{Personnes} & \textbf{mates} (informel), \textbf{peers}, \textbf{neighbours}, \textbf{elders}, \textbf{guest of honour}, \textbf{sponsors}, \textbf{volunteers}, \textbf{organisers}, \textbf{authorities} (Principal, Discipline Master, Mayor, Fon) & Camarades, pairs, voisins, aînés, invité d'honneur, sponsors/parrains, bénévoles, organisateurs, autorités \\ \hline
\textbf{Lieux} & \textbf{premises} (toujours pl.), \textbf{hall}, \textbf{community centre}, \textbf{quarter}, \textbf{village square}, \textbf{stadium} & Locaux/enceinte, salle, centre communautaire, quartier, place du village, stade \\ \hline
\textbf{Événements} & \textbf{festival}, \textbf{celebration}, \textbf{ceremony}, \textbf{cultural week}, \textbf{end-of-year party}, \textbf{bash} (argot), \textbf{feast}, \textbf{funfair}, \textbf{harvest festival} & Festival, célébration, cérémonie, semaine culturelle, fête de fin d'année, fête (fam.), festin, fête foraine, fête des récoltes \\ \hline
\textbf{Actions / Verbes} & \textbf{organise / arrange / secure}, \textbf{volunteer (to)}, \textbf{interact (with)}, \textbf{strengthen (bonds/ties)}, \textbf{sponsor}, \textbf{fund}, \textbf{attend}, \textbf{deliver a speech}, \textbf{propose a vote of thanks}, \textbf{grant permission} & Organiser, se porter volontaire, interagir, renforcer (liens), parrainer, financer, assister à, prononcer un discours, proposer une motion de remerciement, accorder la permission \\ \hline
\textbf{Nourriture / Boissons} & \textbf{sumptuous feast}, \textbf{refreshments}, \textbf{local dishes} (\textit{achu, eru, water fufu, ndolé}), \textbf{palm wine}, \textbf{soft drinks} & Festin somptueux, rafraîchissements, plats locaux, vin de palme, boissons gazeuses \\ \hline
\textbf{Notions abstraites} & \textbf{cultural heritage}, \textbf{unity}, \textbf{bonds of neighbourliness}, \textbf{highlights}, \textbf{permission}, \textbf{authorisation}, \textbf{speech}, \textbf{atmosphere} & Patrimoine culturel, unité, liens de bon voisinage, moments forts, permission, autorisation, discours, ambiance \\ \hline
\end{tabularx}
\end{center}

\begin{propriete}[Prononciation repère (approximation française)]
Contrairement au français, l'anglais a une accentuation tonique variable. Les syllabes en \textbf{gras} sont accentuées.
\begin{itemize}
\item \textbf{volun}\textbf{teer} « vo-lon-\textbf{tir} » (accent sur dernière syllabe)
\item \textbf{spon}\textbf{sor} « \textbf{spon}-seu »
\item \textbf{au}\textbf{thorities} « o-\textbf{to}-ri-tiz »
\item \textbf{prem}\textbf{ises} « \textbf{prem}-i-siz » (pluriel en \textit{-es} prononcé \textit{-iz})
\item \textbf{neigh}\textbf{bour} « \textbf{nei}-beu » (gh = /f/ ou muet, ici /ei/)
\item \textbf{her}\textbf{itage} « \textbf{he}-ri-tidj »
\item \textbf{high}\textbf{light} « \textbf{hai}-lait »
\end{itemize}
\end{propriete}

\begin{attention}[Orthographe : Pièges pour francophones]
\begin{itemize}
\item \textbf{Neighbour} (UK) vs \textit{Neighbor} (US) : gardez \textit{ou} (couleur, honneur, travail). Idem pour \textit{colour, honour, labour}.
\item \textbf{Organise / Organisation} (UK) : \textit{-ise} / \textit{-isation} (pas \textit{-ize} / \textit{-ization} bien que toléré, la norme UK éducative est \textit{-ise}).
\item \textbf{Premises} : toujours au \textbf{pluriel} (pas de \textit{premise} au sens de « locaux »). \textit{A premise} = une prémisse (logique).
\item \textbf{Sponsor} : pas de \textit{ss} (un seul \textit{s}).
\item \textbf{Feast} : \textit{ea} = /i:/ (comme \textit{beast}), pas /e/.
\end{itemize}
\end{attention}

\coursec{Collocations and idiomatic expressions}

Les \cle{collocations} sont des associations de mots « figées » par l'usage. Les \cle{idiomes} ont un sens global non déductible des mots isolés. Les maîtriser rend votre anglais naturel et précis (critère \textit{Vocabulary} au Bac).

\begin{exemplebox}[Collocations fortes du thème (à apprendre par cœur)]
\begin{itemize}
\item \textbf{Guest of honour} (invité d'honneur) — \textit{The Fon was the guest of honour.}
\item \textbf{Cultural heritage} (patrimoine culturel) — \textit{We must preserve our cultural heritage.}
\item \textbf{Sumptuous feast} (festin somptueux) — \textit{A sumptuous feast was served.}
\item \textbf{Vote of thanks} (motion/discours de remerciement) — \textit{The prefect proposed a vote of thanks.}
\item \textbf{Strengthen bonds / ties} (renforcer les liens) — \textit{Festivals strengthen community bonds.}
\item \textbf{Grant permission} (accorder la permission) — \textit{The Principal granted permission.}
\item \textbf{Deliver a speech} (prononcer un discours) — \textit{She delivered a moving speech.}
\item \textbf{Interact respectfully with} (interagir respectueusement avec) — \textit{Students must interact respectfully with elders.}
\item \textbf{Buzz with excitement / activity} (bourdonner d'excitation/d'activité) — \textit{The hall buzzed with excitement.}
\end{itemize}
\end{exemplebox}

\begin{exemplebox}[Expressions idiomatiques utiles (registre varié)]
\begin{center}
\begin{tabularx}{\linewidth}{|l|X|X|}\hline
\rowcolor{poporangeL} \textbf{Idiome} & \textbf{Sens (FR)} & \textbf{Registre / Exemple} \\ \hline
\textbf{Have a blast} & S'éclater, s'amuser follement & Informel. \eng{We had a blast at the party.} \\ \hline
\textbf{Have a whale of a time} & Passer un moment formidable & Neutre / Semi-formel. \eng{The children had a whale of a time at the funfair.} \\ \hline
\textbf{The icing on the cake} & La cerise sur le gâteau (point final parfait) & Neutre. \eng{The traditional dance was the icing on the cake.} \\ \hline
\textbf{Paint the town red} & Faire la fête toute la nuit / faire la bringue & Très informel, argot. \textbf{Évitez} dans un écrit scolaire formel. \\ \hline
\textbf{Make a day of it} & Profiter de la journée entière / en faire une journée complète & Neutre. \eng{Let's go early and make a day of it.} \\ \hline
\end{tabularx}
\end{center}
\end{exemplebox}

\coursec{Interacting with different people: register and politeness}

La compétence \textit{Interacting} exige d'adapter son \cle{registre de langue} (niveau de formalité) à son interlocuteur. Au Cameroun, le respect de la hiérarchie (âge, fonction) est culturellement crucial et se reflète dans l'anglais formel attendu aux examens.

\begin{textebox}[Dialogue 1: Between mates — Organising the bash (Informel)]
\textbf{A:} Hey guys, the end-of-year \textbf{bash} is next Friday. Who’s bringing the \textbf{soft drinks}?
\textbf{B:} I’ll \textbf{sort out} the sound system. We need a playlist that \textbf{slaps}.
\textbf{C:} Cool. I’ll ask my cousin to \textbf{bake a cake}. Let’s \textbf{make it lit}!
\textbf{A:} \textbf{Deal}. Don’t forget to invite the form teacher, though. She’s pretty \textbf{chill}.
\end{textebox}

\begin{exemplebox}[Dialogue 2: Class Prefect → Principal — Requesting the hall (Formel)]
\eng{Good morning, Sir. I am writing on behalf of Form Five to request permission to use the school hall for our end-of-year celebration.} « Bonjour Monsieur. J'écris au nom de la classe de 1ère pour demander l'autorisation d'utiliser la salle des fêtes pour notre célébration de fin d'année. »
\eng{We would be grateful if you could approve our request by Wednesday.} « Nous vous serions reconnaissants si vous pouviez approuver notre demande d'ici mercredi. »
\eng{We assure you that the premises will be left clean and tidy.} « Nous vous assurons que les lieux seront laissés propres et rangés. »
\eng{Thank you for your consideration. Yours faithfully, [Name], Class Prefect.} « Merci de votre considération. Cordialement, [Nom], Délégué de classe. »
\end{exemplebox}

\begin{textebox}[Dialogue 3: Student → Elder / Neighbour — Invitation (Semi-formel / Respectueux)]
\eng{Good afternoon, Papa. We are organising the quarter's cultural festival next Saturday. We would be honoured if you could attend as our guest of honour.} « Bonjour Papa. Nous organisons le festival culturel du quartier samedi prochain. Nous serions honorés si vous pouviez y assister en tant qu'invité d'honneur. »
\eng{Your presence would mean a lot to the community.} « Votre présence signifierait beaucoup pour la communauté. »
\eng{We have arranged transport for you. Please let us know if you can make it.} « Nous avons prévu le transport pour vous. Dites-nous si vous pouvez venir. »
\end{textebox}

\begin{propriete}[Adapter son registre : Règles clés pour le Bac]
\begin{itemize}
\item \textbf{Salutations :} Informel \eng{Hey / Hi guys / Hello everyone} vs Formel \eng{Good morning Sir/Madam, Dear Principal, Dear Sir/Madam}.
\item \textbf{Présentation de l'objectif :} Informel \eng{We wanna organise a party.} vs Formel \eng{I am writing to request permission to organise a celebration / an event.}
\item \textbf{Formulation de la demande :} Informel \eng{Can we use the hall? / Give us the key.} vs Formel \eng{May we use the hall? / We would be grateful if you could grant us access / authorise the use of the premises.}
\item \textbf{Engagement / Promesse :} Informel \eng{We'll clean up. Don't worry.} vs Formel \eng{We assure you the premises will be left clean and tidy. / We undertake to restore the premises to their original state.}
\item \textbf{Contractions :} Informel \eng{I'll, we'll, don't, let's, can't} (acceptables à l'oral, SMS, dialogue familier) vs Formel \textbf{Éviter les contractions} : \eng{I will, we will, do not, let us, cannot} (écrit administratif, lettre, composition, discours).
\item \textbf{Vocabulaire précis :} \cle{bash / do / party} (informel) vs \cle{celebration / event / function / ceremony} (formel). \cle{Sort out / fix / get} (informel) vs \cle{arrange / organise / secure / obtain} (formel).
\item \textbf{Marqueurs de politesse (formel) :} \cle{please}, \cle{kindly}, \cle{would be grateful}, \cle{could you please}, \cle{thank you for your time / consideration}.
\end{itemize}
\end{propriete}

\begin{center}
\begin{tabularx}{\linewidth}{|X|X|X|}\hline
\rowcolor{popblueL} \textbf{Fonction communicative} & \textbf{Informel (Mates / Peers)} & \textbf{Formel (Principal / Authorities / Elders)} \\ \hline
Salutation / Ouverture & Hey guys! / Hi everyone! & Good morning, Sir / Madam. Dear Principal, \\ \hline
Exposer le but & We wanna organise a party. & I am writing to request permission to organise a celebration. \\ \hline
Faire une demande & Can we use the hall? / Give us the key. & May we use the hall? We would be grateful if you could grant us access. \\ \hline
S'engager (nettoyage, discipline) & We'll clean up after. Don't worry. & We assure you the premises will be left clean and tidy. \\ \hline
Clôture / Formule de politesse & See you there! Cheers! / Later! & Thank you for your consideration. Yours faithfully, / Yours sincerely, \\ \hline
\end{tabularx}
\end{center}

\coursec{False friends and common traps for francophones}

Ces erreurs coûtent des points en \textit{Vocabulary} et \textit{Grammar} au Bac. La traduction littérale du français est le piège n°1.

\begin{definition}[Faux amis : Mots qui se ressemblent mais diffèrent en sens]
\begin{center}
\begin{tabularx}{\linewidth}{|l|l|X|l|}\hline
\rowcolor{popredL} \textbf{Mot FR} & \textbf{Faux ami EN} & \textbf{Sens réel du faux ami} & \textbf{Bon mot EN (contexte fête)} \\ \hline
Assister (à) & \eng{Assist} & Aider, secourir (\eng{The nurse assisted the doctor.}) & \eng{\textbf{Attend}} (\eng{We attended the festival.}) \\ \hline
Demander (poliment) & \eng{Demand} & Exiger avec autorité (\eng{The boss demanded an explanation.}) & \eng{\textbf{Ask / Request}} (\eng{We asked for permission.}) \\ \hline
Actuel / D'actualité & \eng{Actual} & Réel, effectif (\eng{The actual cost was higher.}) & \eng{\textbf{Current / Topical / Present}} (\eng{A current topic.}) \\ \hline
Passer (un moment) & \eng{Pass} (time) & Passer le temps (s'occuper) (\eng{We passed the time playing cards.}) & \eng{\textbf{Have a good time / Enjoy ourselves / Have a blast}} \\ \hline
Responsable de & \eng{Responsible of} & \textbf{Incorrect} (préposition erronée) & \eng{\textbf{Responsible FOR / In charge OF}} \\ \hline
Éventuellement & \eng{Eventually} & Finalement, après un délai (\eng{He eventually arrived.}) & \eng{\textbf{Possibly / If necessary}} \\ \hline
Sensible (émotion) & \eng{Sensible} & Raisonnable, sensé (\eng{A sensible decision.}) & \eng{\textbf{Sensitive}} (\eng{He is sensitive to criticism.}) \\ \hline
Blesser & \eng{Bless} & Bénir (\eng{God bless you.}) & \eng{\textbf{Injure / Hurt / Wound}} \\ \hline
Préservatif & \eng{Preservative} & Conservateur (alimentaire) (\eng{Without preservatives.}) & \eng{\textbf{Condom}} \\ \hline
Location (louer) & \eng{Location} & Lieu, emplacement (\eng{A filming location.}) & \eng{\textbf{Rental / Hire}} \\ \hline
\end{tabularx}
\end{center}
\end{definition}

\begin{attention}[Autres pièges structurels (Calques syntaxiques)]
\begin{itemize}
\item \textbf{Demander \textit{à} quelqu'un \textit{de} faire} : \eng{Ask someone \textbf{to} do sth} (pas \eng{ask someone \textbf{de} do} ni \eng{demand someone to do}).
\item \textbf{Permettre \textit{à} quelqu'un \textit{de} faire} : \eng{Allow / Enable someone \textbf{to} do sth} / \eng{Give someone permission \textbf{to} do sth}.
\item \textbf{S'intéresser \textit{à}} : \eng{Be interested \textbf{in}} (pas \eng{interested \textbf{to}} / \eng{interested \textbf{by}}).
\item \textbf{Participer \textit{à}} : \eng{Take part \textbf{in}} / \eng{Participate \textbf{in}} (pas \eng{participate \textbf{to}}).
\item \textbf{Insister \textit{sur}} : \eng{Insist \textbf{on}} + \textit{-ing} (\eng{He insisted on paying.}).
\end{itemize}
\end{attention}

\coursec{Word families, pronunciation and spelling}

Travailler les familles de mots (dérivation) permet de multiplier son lexique et de réussir les exercices de \textit{Word Formation} (très fréquents au Bac, Paper 1 \& 2).

\begin{propriete}[Dérivation nominale et adjectivale : Tableau de formation]
\begin{center}
\begin{tabular}{|l|l|l|l|l|}\hline
\rowcolor{poppurpleL} \textbf{Verbe} & \textbf{Nom (action/état)} & \textbf{Nom (agent)} & \textbf{Adjectif} & \textbf{Adverbe} \\ \hline
\eng{celebrate} & \eng{celebration} & \eng{celebrant} & \eng{celebratory} & \eng{celebratorily} \\ \hline
\eng{organise} & \eng{organisation} & \eng{organiser} & \eng{organised / organisational} & \eng{organisationally} \\ \hline
\eng{decorate} & \eng{decoration} & \eng{decorator} & \eng{decorative} & \eng{decoratively} \\ \hline
\eng{invite} & \eng{invitation} & \eng{invitee / inviter} & \eng{inviting} & \eng{invitingly} \\ \hline
\eng{participate} & \eng{participation} & \eng{participant} & \eng{participatory} & \eng{participatorily} \\ \hline
\eng{contribute} & \eng{contribution} & \eng{contributor} & \eng{contributory} & — \\ \hline
\eng{sponsor} & \eng{sponsorship} & \eng{sponsor} & \eng{sponsored} & — \\ \hline
\eng{interact} & \eng{interaction} & — & \eng{interactive} & \eng{interactively} \\ \hline
\eng{unite} & \eng{unity / union} & \eng{unifier} & \eng{united / unitary} & \eng{unitedly} \\ \hline
\eng{entertain} & \eng{entertainment} & \eng{entertainer} & \eng{entertaining} & \eng{entertainingly} \\ \hline
\end{tabular}
\end{center}
\end{propriete}

\begin{methode}[Règles d'orthographe et de prononciation pour ces familles]
\begin{enumerate}
\item \textbf{-ise / -isation (UK)} : \eng{organise, organisation, organiser}. Ne pas écrire \textit{organize} (US) dans les copies officielles MINESEC/GCE.
\item \textbf{Double consonne} : \eng{occur $\rightarrow$ occurrence}, \eng{refer $\rightarrow$ referral}, \eng{begin $\rightarrow$ beginning}. Règle : accent sur dernière syllabe + voyelle brève + consonne finale $\rightarrow$ double consonne avant \textit{-ence, -al, -ing, -ed}.
\item \textbf{-our / -or} : UK \eng{neighbour, colour, honour, labour, flavour} / US \eng{neighbor, color, honor, labor, flavor}. \textbf{Exigence UK au Cameroun}.
\item \textbf{-re / -er} : UK \eng{centre, theatre, metre, fibre} / US \eng{center, theater, meter, fiber}.
\item \textbf{Préfixes négatifs} : \eng{disorganised, disunity, non-participation, irresponsible, inactive}. Choisir le bon préfixe (\textit{un-, in-, im-, ir-, dis-, non-}).
\item \textbf{Accentuation (stress shift)} : \eng{ORganise (v) $\rightarrow$ organiSAtion (n)}, \eng{INvite (v) $\rightarrow$ inviTAtion (n)}, \eng{CELEbrate (v) $\rightarrow$ celeBRAtion (n)}. L'accent bouge souvent vers le suffixe \textit{-ation, -ity, -ic}.
\end{enumerate}
\end{methode}

\begin{exoresolu}[Word Formation : Complétez avec la forme correcte du mot entre parenthèses.]
\textbf{Énoncé :}
\begin{enumerate}
\item The \_\_\_\_\_\_\_\_\_\_\_\_ (organise) of the festival did a great job.
\item We need active \_\_\_\_\_\_\_\_\_\_\_\_ (participate) from all students.
\item The traditional dance was the \_\_\_\_\_\_\_\_\_\_\_\_ (high) of the evening.
\item It is \_\_\_\_\_\_\_\_\_\_\_\_ (respond) to arrive late when you are a guest of honour.
\item The \_\_\_\_\_\_\_\_\_\_\_\_ (celebrate) lasted three days.
\end{enumerate}
\tcblower
\textbf{Solution détaillée :}
\begin{enumerate}
\item \textbf{organisers} (nom agent, pluriel) — \textit{The organisers did a great job.}
\item \textbf{participation} (nom action, indénombrable) — \textit{Active participation.}
\item \textbf{highlight} (nom, formation par conversion/préfixe \textit{high} + \textit{light}) — \textit{The highlight of the evening.}
\item \textbf{irresponsible} (adj. négatif \textit{ir-} + \textit{responsible}) — \textit{It is irresponsible to arrive late.}
\item \textbf{celebration} (nom action, suffixe \textit{-tion}) — \textit{The celebration lasted three days.}
\end{enumerate}
\end{exoresolu}

\coursec{Using the vocabulary in context}

Après avoir exploré les familles de mots, la prononciation et l'orthographe, nous passons à l'étape cruciale : \textbf{la mise en contexte communicatif}. Connaître une liste de mots ne suffit pas ; il faut savoir les \cle{combiner} (collocations), adapter son \cle{registre de langue} selon son interlocuteur et éviter les \cle{faux amis}. Cette section vous entraîne par la lecture, l'analyse de dialogues, des exercices de consolidation et une production guidée évaluée.

\courssub{Reading comprehension: The End-of-Year Community Festival in Bamenda}

Observez le lexique cible dans un texte narratif cohérent, ancré dans le contexte camerounais. Identifiez les mots du champ lexical des fêtes, les collocations et le ton employé.

\begin{textebox}[The End-of-Year Community Festival in Bamenda]
Last Saturday, the quarter of Nkwen in Bamenda buzzed with excitement for the annual end-of-year community festival. The Fon’s representative, the \textbf{guest of honour}, delivered a heartfelt speech praising unity and \textbf{cultural heritage}. My \textbf{mates} and I \textbf{volunteered} to organise funfair games for the toddlers, while the women’s association prepared a \textbf{sumptuous feast} of achu and yellow soup. The \textbf{highlight} was the traditional dance troupe whose rhythmic steps thrilled the crowd. We \textbf{interacted respectfully} with \textbf{elders}, offering them seats and refreshments. Later, the youth leader proposed a \textbf{vote of thanks} to the \textbf{sponsors}. Everybody \textbf{had a blast}, strengthening the \textbf{bonds of neighbourliness}. It was a memorable day of sharing and celebration.
\end{textebox}

\begin{center}
\begin{tabularx}{\linewidth}{|l|X|X|}\hline
\rowcolor{popblueL} \textbf{English word / phrase} & \textbf{Part of speech} & \textbf{Traduction / Glossaire} \\ \hline
buzzed with excitement & verb phrase & bourdonnait d'excitation, était en effervescence \\ \hline
guest of honour & noun phrase & invité d'honneur \\ \hline
cultural heritage & noun phrase & patrimoine culturel \\ \hline
mates & noun (pl.) & copains, camarades (registre informel) \\ \hline
volunteered to & verb + inf. & s'est porté volontaire pour, a proposé spontanément \\ \hline
sumptuous feast & noun phrase (collocation) & festin somptueux, repas copieux \\ \hline
highlight & noun & moment fort, point culminant \\ \hline
interacted respectfully & verb + adv. & a interagi respectueusement \\ \hline
elders & noun (pl.) & aînés, personnes âgées (terme de respect) \\ \hline
vote of thanks & noun phrase (collocation) & motion de remerciement, discours de remerciement \\ \hline
sponsors & noun (pl.) & sponsors, parrains (financeurs) \\ \hline
had a blast & idiom (verb phrase) & s'est éclaté, s'est bien amusé (informel) \\ \hline
strengthening the bonds of neighbourliness & verb phrase & renforçant les liens de bon voisinage \\ \hline
\end{tabularx}
\end{center}

\begin{exoresolu}[Compréhension et repérage lexical]
\textbf{Énoncé :} Relisez le texte et répondez en anglais.
\begin{enumerate}
\item Who was the guest of honour at the festival?
\item What collocation describes the meal prepared by the women's association?
\item What did the mates organise for the toddlers?
\item Find the idiom meaning "enjoyed themselves immensely".
\item Which phrase means "renforcer les liens de bon voisinage"?
\end{enumerate}
\tcblower
\textbf{Solution détaillée :}
\begin{enumerate}
\item The Fon’s representative was the guest of honour.
\item A \textbf{sumptuous feast}.
\item They organised funfair games.
\item \textbf{Had a blast}.
\item \textbf{Strengthening the bonds of neighbourliness}.
\end{enumerate}
\end{exoresolu}

\courssub{Consolidation exercises: Gap-fill and Reformulation}

Appliquons le lexique ciblé (\textit{guest of honour, cultural heritage, volunteers, sumptuous feast, highlight, elders, vote of thanks, sponsors, have a blast, strengthen bonds, interact, neighbours, authorities, permission, premises}).

\begin{exoresolu}[Gap-fill : Complete the sentences with the correct word or phrase from the box.]
\textbf{Énoncé :} Complétez les phrases suivantes. Attention à la grammaire (temps, nombre).
\begin{enumerate}
\item The \_\_\_\_\_\_\_\_\_\_\_\_, a local businessman, cut the ribbon to open the new library.
\item During the festival, young people must \_\_\_\_\_\_\_\_\_\_\_\_ respectfully with the \_\_\_\_\_\_\_\_\_\_\_\_.
\item Organising such events helps \_\_\_\_\_\_\_\_\_\_\_\_ the \_\_\_\_\_\_\_\_\_\_\_\_ of neighbourliness in the quarter.
\item We need to write a formal letter to the school \_\_\_\_\_\_\_\_\_\_\_\_ to ask for \_\_\_\_\_\_\_\_\_\_\_\_ to use the hall.
\item The \_\_\_\_\_\_\_\_\_\_\_\_ of the evening was the traditional wrestling match.
\item Many local companies acted as \_\_\_\_\_\_\_\_\_\_\_\_ for the cultural week.
\item At the end of the ceremony, the prefect proposed a \_\_\_\_\_\_\_\_\_\_\_\_ to the organising committee.
\item Everyone \_\_\_\_\_\_\_\_\_\_\_\_ at the beach party last weekend; the music was amazing.
\item The women prepared a \_\_\_\_\_\_\_\_\_\_\_\_ of \textit{eru} and \textit{water fufu}.
\item Please ensure the \_\_\_\_\_\_\_\_\_\_\_\_ are locked after the event.
\end{enumerate}
\tcblower
\textbf{Solution détaillée :}
\begin{enumerate}
\item \textbf{guest of honour} (sujet, singulier)
\item \textbf{interact} (infinitif sans \textit{to} après modal \textit{must}) / \textbf{elders} (complément, pluriel, terme de respect)
\item \textbf{strengthen} (infinitif sans \textit{to} après \textit{helps}) / \textbf{bonds} (collocation \textit{strengthen bonds})
\item \textbf{authorities} (pluriel habituel pour l'administration scolaire) / \textbf{permission} (nom indénombrable, \textit{ask for permission})
\item \textbf{highlight} (sujet singulier, \textit{The highlight was...})
\item \textbf{sponsors} (pluriel, \textit{Many companies})
\item \textbf{vote of thanks} (collocation figée)
\item \textbf{had a blast} (prétérit de l'idiome \textit{have a blast})
\item \textbf{sumptuous feast} (collocation forte)
\item \textbf{premises} (toujours pluriel en anglais pour « les lieux / le bâtiment »)
\end{enumerate}
\end{exoresolu}

\begin{exercice}[Reformulation : Évitez les faux amis et calques du français]
\textbf{Consigne :} Les phrases ci-dessous contiennent des erreurs typiques de francophones. Réécrivez-les correctement en anglais naturel, en utilisant le vocabulaire de la leçon.
\begin{enumerate}
\item \textit{Français :} "Nous avons \textbf{assisté} à la fête hier." $\rightarrow$ \textit{Anglais incorrect :} "We \textbf{assisted} at the party yesterday."
\item \textit{Français :} "Le Principal a \textbf{demandé} aux élèves de nettoyer." $\rightarrow$ \textit{Anglais incorrect :} "The Principal \textbf{demanded} the students to clean."
\item \textit{Français :} "C'est un \textbf{événement} très \textbf{actuel}." $\rightarrow$ \textit{Anglais incorrect :} "It is a very \textbf{actual} \textbf{event}."
\item \textit{Français :} "Nous avons \textbf{passé} un bon moment." $\rightarrow$ \textit{Anglais incorrect :} "We \textbf{passed} a good moment."
\item \textit{Français :} "Il est \textbf{responsable} de l'organisation." $\rightarrow$ \textit{Anglais incorrect :} "He is \textbf{responsible} of the organisation."
\end{enumerate}
\end{exercice}

\begin{corrige}[Exercice : Reformulation]
\begin{enumerate}
\item \eng{We \textbf{attended} the party yesterday.} (Faux ami : \textit{assist} = aider ; \textit{attend} = assister à / être présent).
\item \eng{The Principal \textbf{asked} the students to clean.} / \eng{The Principal \textbf{requested that} the students clean.} (Faux ami : \textit{demand} = exiger avec autorité ; \textit{ask/request} = demander poliment. Structure : \textit{ask someone to do sth}).
\item \eng{It is a very \textbf{current} / \textbf{topical} event.} / \eng{It is a \textbf{major} event.} (Faux ami : \textit{actual} = réel, effectif ; \textit{current/topical} = d'actualité).
\item \eng{We \textbf{had a great time} / \textbf{enjoyed ourselves} / \textbf{had a blast}.} (Calque : \textit{pass time} = passer le temps ; \textit{have a good time} = passer un bon moment).
\item \eng{He is \textbf{responsible for} the organisation.} (Préposition : \textit{responsible FOR} / \textit{in charge OF}, pas \textit{of}).
\end{enumerate}
\end{corrige}

\courssub{Guided production: Writing about a village festival}

Voici la tâche finale : produire un court paragraphe (5 phrases) en réinvestissant activement le lexique. Suivez la méthode rigoureuse ci-dessous.

\begin{methode}[Réussir sa production écrite guidée]
\begin{enumerate}
\item \textbf{Lister} 5 à 7 mots/expressions du champ lexical (noms, verbes, collocations, idiomes) que vous comptez utiliser. Cochez-les au fur et à mesure.
\item \textbf{Identifier} le destinataire / le contexte : Écrivez-vous pour des amis (informel), pour un journal scolaire (standard / semi-formel), ou pour une autorité (formel) ? Cela dicte le registre.
\item \textbf{Rédiger} les 5 phrases en insérant \textbf{au minimum 3 collocations} (ex. \textit{guest of honour, sumptuous feast, vote of thanks, strengthen bonds, cultural heritage}) et \textbf{1 expression idiomatique} (ex. \textit{have a blast, have a whale of a time, the icing on the cake}).
\item \textbf{Relire} systématiquement en traquant : les \textbf{faux amis} (\textit{actually, assist, demand, pass, responsible of}), l'\textbf{orthographe} (\textit{premises, neighbours, sponsors, organisation}), la \textbf{ponctuation} et les \textbf{majuscules} (jours, mois, nationalités, \textit{I}).
\item \textbf{Vérifier} l'accord des temps (souvent \textit{past simple} pour un récit passé) et la cohérence du registre.
\end{enumerate}
\end{methode}

\begin{popfigure}
\begin{tikzpicture}[
    node distance=1.2cm and 1.5cm,
    every node/.style={font=\small, align=center},
    startstop/.style={rectangle, rounded corners=5pt, minimum width=2.5cm, minimum height=1cm, text centered, draw=popblue, fill=popblueL, thick},
    process/.style={rectangle, minimum width=3cm, minimum height=1cm, text centered, draw=poporange, fill=poporangeL, thick},
    decision/.style={diamond, aspect=2, minimum width=2.5cm, minimum height=1cm, text centered, draw=popgreen, fill=popgreenL, thick},
    arrow/.style={-Stealth, thick, popdark}
]
\node (start) [startstop, align=center]{DÉBUT\\Choisir l'événement};
\node (lexique) [process, below of=start, align=center]{Lister 5--7 mots\\collocations / idiomes};
\node (registre) [decision, below of=lexique, align=center]{Quel destinataire ?\\Mates / Adultes / Autorités};
\node (informel) [process, left of=registre, xshift=-2.5cm, align=center]{Registre INFORMEL\\Hey, bash, have a blast, sort out};
\node (formel) [process, right of=registre, xshift=2.5cm, align=center]{REGISTRE FORMEL\\Dear Sir, celebration, request, grateful};
\node (redaction) [process, below of=registre, yshift=-0.5cm, align=center]{Rédiger 5 phrases\\Insérer 3 colloc + 1 idiome};
\node (verif) [process, below of=redaction, align=center]{Relire \& Corriger\\Faux amis, temps, registre};
\node (end) [startstop, below of=verif, align=center]{FIN\\Production finale};

\draw [arrow] (start) -- (lexique);
\draw [arrow] (lexique) -- (registre);
\draw [arrow] (registre) -- node[left, pos=0.5, fill=popcream, rounded corners, inner sep=1pt] {Informel} (informel);
\draw [arrow] (registre) -- node[right, pos=0.5, fill=popcream, rounded corners, inner sep=1pt] {Formel} (formel);
\draw [arrow] (informel) |- (redaction);
\draw [arrow] (formel) |- (redaction);
\draw [arrow] (redaction) -- (verif);
\draw [arrow] (verif) -- (end);
\end{tikzpicture}
\legende{Organigramme de la production guidée : du choix lexical à la vérification finale.}
\end{popfigure}

\begin{exercice}[Mini-production guidée : My Village Festival]
\textbf{Consigne :} Rédigez un paragraphe de \textbf{exactement 5 phrases} pour décrire la fête annuelle de votre village ou quartier.
\begin{itemize}
\item Contexte : Vous écrivez pour le journal de l'école (registre \textbf{standard / semi-formel}).
\item Contraintes obligatoires :
  \begin{itemize}
  \item Utiliser au moins \textbf{3 collocations} distinctes (ex. \textit{guest of honour, cultural heritage, sumptuous feast, vote of thanks, strengthen bonds}).
  \item Utiliser au moins \textbf{1 expression idiomatique} (ex. \textit{have a blast, have a whale of a time, the icing on the cake}).
  \item Employer correctement \textbf{elders, sponsors, premises, permission, interact, highlight, volunteers}.
  \item Zéro faux ami (\textit{actually, assist, demand, pass, responsible of}).
  \end{itemize}
\end{itemize}
\textit{Aide au démarrage :} "Last December, my village organised its annual cultural festival..."
\end{exercice}

\courssub{Self-assessment checklist}

Avant de remettre votre production, utilisez cette grille d'auto-évaluation. Cochez \textbf{oui} seulement si vous êtes certain. Si vous cochez \textbf{non}, revenez à la case correspondante dans la méthode.

\begin{center}
\begin{tabularx}{\linewidth}{|X|c|c|X|}\hline
\rowcolor{popblueL} \textbf{Critère de vérification} & \textbf{Oui} & \textbf{Non} & \textbf{Action corrective / Référence leçon} \\ \hline
1. \textbf{Registre adapté} : Le ton convient au journal scolaire (standard, pas de \textit{Hey/guys/bash}, pas de \textit{Yours faithfully} trop formel). & $\square$ & $\square$ & Relire \textit{Propriété : Adapter son registre} \\ \hline
2. \textbf{Collocations (min 3)} : J'ai utilisé au moins 3 associations naturelles (ex. \textit{guest of honour, sumptuous feast, vote of thanks, strengthen bonds, cultural heritage}). & $\square$ & $\square$ & Relire \textit{Texte support} \& \textit{Section Collocations} \\ \hline
3. \textbf{Expression idiomatique (min 1)} : J'ai inséré un idiome correct (ex. \textit{had a blast, had a whale of a time}). & $\square$ & $\square$ & Relire \textit{Dialogue 1} \& \textit{Section Collocations} \\ \hline
4. \textbf{Vocabulaire cible} : Mots imposés employés correctement (\textit{elders, sponsors, premises, permission, interact, highlight, volunteers}). & $\square$ & $\square$ & Relire \textit{Exercice Gap-fill} \\ \hline
5. \textbf{Zéro faux ami} : Pas de \textit{actually} pour "actuellement", \textit{assist} pour "assister", \textit{demand} pour "demander", \textit{pass} pour "passer un moment", \textit{responsible of}. & $\square$ & $\square$ & Relire \textit{Exercice Reformulation} \& \textit{Section Faux amis} \\ \hline
6. \textbf{Grammaire \& Orthographe} : Past simple correct, \textit{premises} (pl.), \textit{neighbours} (UK), majuscules (\textit{I, December, Bamenda}). & $\square$ & $\square$ & Relire \textit{Section Word families, spelling} \\ \hline
7. \textbf{Cohérence \& Longueur} : Exactement 5 phrases, lien logique entre elles. & $\square$ & $\square$ & Compter les phrases (points finals). \\ \hline
\end{tabularx}
\end{center}

\begin{attention}[Pièges à éviter en production formelle / semi-formelle]
\begin{itemize}
\item \textbf{Contractions excessives :} N'écrivez pas \eng{We're gonna have a party} ou \eng{Don't forget} dans un courrier au Principal. Écrivez \eng{We are going to organise} / \eng{Please do not forget}.
\item \textbf{Interjections familières :} Bannissez \eng{Hey!}, \eng{Wow!}, \eng{Cool!}, \eng{OK} dans un écrit adressé à une autorité. Préférez \eng{Please note}, \eng{We would like to inform you}.
\item \textbf{Vocabulaire trop vague :} Évitez \eng{thing, stuff, get, do, make} (sauf collocations figées comme \eng{make a speech}). Préférez \eng{arrange, organise, secure, obtain, event, celebration, refreshments}.
\item \textbf{Traduction mot à mot :} \eng{We passed a good moment} $\rightarrow$ \eng{We had a great time}. \eng{The Principal accepted our demand} $\rightarrow$ \eng{The Principal granted our request}.
\end{itemize}
\end{attention}

\begin{savaistu}[Le saviez-vous ? Le Bac et les cérémonies]
Au Baccalauréat (Paper 2 - Writing), les sujets de \textbf{composition} ou de \textbf{lettre formelle} portent très fréquemment sur :
\begin{itemize}
\item La description d'une cérémonie traditionnelle (funérailles, mariage, fête des récoltes, fête nationale) : \textit{"Describe a traditional ceremony you attended..."}
\item La rédaction d'une \textbf{lettre d'invitation} adressée à une autorité (Principal, Mayor, Fon, Minister) : \textit{"Write a letter to the Principal inviting him to the school cultural week..."}
\item Un \textbf{discours} (speech) pour un \textit{vote of thanks} ou une cérémonie de remise de prix.
\end{itemize}
Maîtriser ce lexique (\textit{guest of honour, cultural heritage, sponsors, vote of thanks, strengthen bonds, premises, elders, interact respectfully}) et la distinction registres (formel/informel) vous donne un \textbf{avantage décisif} pour gagner des points en \textit{Vocabulary} et \textit{Register} aux grilles d'évaluation officielles.
\end{savaistu}

\newpage

\coursec{Activité d'intégration}

\coursec{Lesson 3 — Vocabulary: Interacting on Community Celebrations and Festivals}

\courssub{Objectifs de la leçon}
\noindent
À l'issue de cette leçon, l'élève sera capable de :
\begin{itemize}
    \item \textbf{Identifier et mobiliser} le vocabulaire thématique des fêtes communautaires (\eng{celebration}, \eng{festival}, \eng{gathering}, \eng{festivities}, \eng{venue}, \eng{refreshments}, \eng{entertainment}, \eng{potluck}, \eng{logistics}, \eng{contribution}, \eng{attendance}, \eng{RSVP}, \eng{chief guest}).
    \item \textbf{Adapter} son registre de langue selon l'interlocuteur : \textbf{informel} entre \eng{mates} (camarades) et \textbf{formel} avec les \eng{adults} et \eng{school authorities} (autorités scolaires).
    \item \textbf{Employer} correctement des \textbf{collocations} figées (\eng{extend an invitation}, \eng{grace the occasion}, \eng{logistical arrangements}, \eng{foster unity}, \eng{kindly confirm attendance}) et des \textbf{expressions idiomatiques} (\eng{bring a plate}, \eng{the more the merrier}, \eng{chip in}, \eng{have a blast}, \eng{sort out}, \eng{good shout}).
    \item \textbf{Éviter} les \textbf{faux amis} majeurs (\eng{actually}, \eng{assist}, \eng{entertainment}, \eng{eventually}, \eng{sensible}, \eng{pretend}, \eng{realise}) et les calques syntaxiques du français.
    \item \textbf{Maîtriser} l'orthographe, la prononciation et la dérivation des mots-clés (familles de mots).
    \item \textbf{Produire} un dialogue de planification (oral) et une lettre d'invitation formelle (écrit) ancrés dans le contexte camerounais.
\end{itemize}

\courssub{Situation déclenchante}
\noindent
Le contexte est le \textbf{Lycée Bilingue de Buea} (établissement type du système francophone/anglophone camerounais). La classe de \textbf{1ère A} organise son \eng{End-of-Year Class Festival} (fête de fin d'année) pour marquer la fin du premier cycle du secondaire. Les élèves, par l'intermédiaire de leur \eng{Organising Committee}, doivent gérer l'organisation interne (entre camarades) et la communication institutionnelle (vers l'administration).

\begin{textebox}[Documents déclencheurs : Contexte et échanges]
\textbf{Document A : Extrait du règlement intérieur — \eng{School Policy on Student Events}}
\begin{enumerate}
    \item Any \eng{student-organised event} (\eng{student-led event}) must have a \eng{faculty advisor} (\fr{enseignant référent}).
    \item A formal \eng{letter of invitation} must be submitted to the \eng{Principal's Office} at least two weeks before the \eng{event date}.
    \item The \eng{venue} (\fr{lieu}) must be booked via the \eng{School Bursar} (\fr{Intendant}).
    \item \eng{Refreshments} must comply with school health regulations (no alcohol, labelled allergens).
\end{enumerate}

\textbf{Document B : Échange WhatsApp entre deux délégués de classe (registre informel)}
\begin{verbatim}
Ngwa: Hey mate, the end-of-year bash is coming up. Have you sorted the venue?
Fon: Not yet. The main hall is booked for exams. Maybe the courtyard?
Ngwa: Good shout. We need to sort food too. "Bring a plate" style?
Fon: Yeah, potluck. The more the merrier! But we MUST invite the Principal formally.
Ngwa: Right. Formal letter. "Dear Sir..." No slang there. I'll draft it tonight.
Fon: Cool. Let's make it a blast!
\end{verbatim}
\end{textebox}

\courssub{Mise en œuvre : Activité d'intégration (Tâche complexe)}
\noindent
\textbf{Votre mission :} En groupe de 4 élèves (deux binômes), vous formez le \eng{Organising Committee}. Préparez l'événement sous ses deux facettes : l'organisation interne (oral) et la communication institutionnelle (écrit).

\begin{methode}[Consignes de l'activité — \eng{Task Instructions}]
\begin{enumerate}
    \item \textbf{Étape 1 — Planification orale (Binôme A) :} \textbf{Durée : 15 min préparation + 2 min présentation.}
    \begin{itemize}
        \item Imaginez un dialogue entre deux camarades (\eng{mates/classmates}) qui organisent la fête.
        \item \textbf{Contenu obligatoire :} Date, \eng{venue} (lieu), \eng{activities} (danse, sketch, match), \eng{refreshments}/\eng{potluck} (nourriture), \eng{budget}/\eng{contribution}.
        \item \textbf{Contraintes linguistiques :}
        \begin{itemize}
            \item Utilisez \textbf{au moins 5 mots/expressions du champ lexical} (voir Ressources).
            \item Intégrez \textbf{2 expressions idiomatiques minimum} (ex. \eng{to bring a plate}, \eng{the more the merrier}, \eng{to have a blast}, \eng{to chip in}).
            \item Respectez le \textbf{registre informel} : contractions (\eng{we'll, let's, wanna}), vocabulaire familier (\eng{mate, guys, cool, sorted, bash}), phrases courtes, ellipses.
        \end{itemize}
        \item Répétez pour une présentation fluide (prononciation, intonation, contact visuel).
    \end{itemize}

    \item \textbf{Étape 2 — Rédaction formelle (Binôme B) :} \textbf{Durée : 20 min rédaction.}
    \begin{itemize}
        \item Rédigez la \textbf{lettre formelle d'invitation} adressée au \eng{Principal}.
        \item \textbf{Structure imposée (Format lettre administrative britannique) :}
        \begin{enumerate}
            \item \eng{Sender's address} (Adresse de l'expéditeur : Lycée..., Buea) + \eng{Date} (format \eng{10th June 2025}).
            \item \eng{Recipient's address} (The Principal, Lycée Bilingue..., Buea).
            \item \eng{Subject line} (\eng{RE: Invitation to End-of-Year Class Festival}).
            \item \eng{Salutation} (\eng{Dear Sir,} ou \eng{Dear Mr [Nom],}).
            \item \textbf{Corps de la lettre (3 paragraphes distincts) :}
            \begin{itemize}
                \item \eng{Opening} : Objet précis, classe concernée, date/heure/lieu.
                \item \eng{Details} : Nature de l'événement (\eng{cultural display, student talents}), objectif (\eng{foster unity, celebrate achievements}), \textbf{3 collocations formelles minimum} (voir Ressources).
                \item \eng{Closing} : \eng{RSVP} (\eng{Kindly confirm your attendance}), formule de politesse soutenue.
            \end{itemize}
            \item \eng{Complimentary close} (\eng{Yours faithfully,}) + \eng{Signatures} (Class Prefect, Deputy Prefect).
        \end{enumerate}
        \item \textbf{Contraintes linguistiques :}
        \begin{itemize}
            \item \textbf{Zéro contraction} (\eng{we will, we would, cannot, it is}).
            \item Vocabulaire soutenu (\eng{request the honour of your presence, grace the occasion, logistical arrangements, anticipate, student-led initiative}).
            \item \textbf{Évitez les faux amis} (voir Ressources / Boîte Attention).
        \end{itemize}
    \end{itemize}

    \item \textbf{Étape 3 — Présentation et évaluation par les pairs :} \textbf{Durée : 15 min.}
    \begin{itemize}
        \item Chaque groupe présente : Binôme A joue le dialogue (debout, sans lire), Binôme B lit la lettre ou la projette.
        \item La classe utilise la \textbf{Grille de vérification} ci-dessous pour valider :
    \end{itemize}
\end{enumerate}

\end{methode}

\begin{definition}[Grille d'auto-évaluation et pair-évaluation — \eng{Peer Assessment Checklist}]
\begin{center}
\begin{tabularx}{\linewidth}{|X|c|}
\hline
\rowcolor{popblueL} \textbf{Critère de réussite} & \textbf{Validé ?} \\
\hline
Registre respecté (informel vs formel) & \square \\
\hline
5 mots du champ lexical + 2 idiomes (oral) / 3 collocations (écrit) & \square \\
\hline
Aucun faux ami détecté (\eng{actually, assist, entertainment, eventually, sensible}) & \square \\
\hline
Prononciation claire (oral) / Orthographe \& ponctuation soignées (écrit) & \square \\
\hline
Cohérence situationnelle (contexte camerounais réaliste : \eng{Bursar, Discipline Master, courtyard, eru}) & \square \\
\hline
\end{tabularx}
\end{center}
\end{definition}

\noindent
L'enseignant note les points forts et les corrections collectives au tableau.

\courssub{Ressources linguistiques à mobiliser}
\noindent
Voici la synthèse des outils lexicaux et grammaticaux vus en classe pour réussir la tâche.

\begin{definition}[Champ lexical de base — \eng{Core Vocabulary: Celebrations \& Festivals}]
\begin{center}
\begin{tabularx}{\linewidth}{|X|X|X|}
\hline
\rowcolor{popblueL} \textbf{Mot / Expression} & \textbf{Nature} & \textbf{Traduction / Contexte camerounais} \\
\hline
\eng{celebration} & n. & fête, célébration (ex. \eng{the celebration of Youth Day}) \\
\eng{festival} & n. & festival (souvent culturel/artistique, ex. \eng{Ngonso Festival}) \\
\eng{gathering} & n. & rassemblement, réunion (ex. \eng{a family gathering}) \\
\eng{festivities} & n. pl. & festivités (ensemble des activités festives) \\
\eng{venue} & n. & lieu de l'événement (ex. \eng{the school courtyard}, \eng{the main hall}) \\
\eng{refreshments} & n. pl. & rafraîchissements, collation (boissons, petits fours) \\
\eng{potluck} & n. & repas partagé (chacun apporte un plat : \eng{bring a plate}) \\
\eng{entertainment} & n. & divertissement, animation, spectacle (\textbf{pas} \fr{entretien}) \\
\eng{logistics} / \eng{logistical arrangements} & n. pl. & logistique, organisation matérielle \\
\eng{contribution} & n. & participation (financière ou en nature) \\
\eng{attendance} & n. & présence, participation (ex. \eng{confirm your attendance}) \\
\eng{RSVP} & abbr. & \eng{Répondez s'il vous plaît} (sur carton d'invitation) \\
\eng{chief guest} / \eng{guest of honour} & n. & invité d'honneur \\
\eng{to grace the occasion} & v. phr. & honorer de sa présence (très formel) \\
\eng{stakeholders} & n. pl. & parties prenantes (admin, parents, chefs traditionnels) \\
\hline
\end{tabularx}
\end{center}
\end{definition}

\begin{definition}[Collocations formelles — \eng{Formal Collocations (pour l'écrit officiel)}]
\noindent Ces associations de mots sont figées ; ne les traduisez pas mot à mot.
\begin{itemize}
    \item \eng{\textbf{extend an invitation} (to sb)} — adresser une invitation (à qqn) \textit{(pas \eng{make/give an invitation})}.
    \item \eng{\textbf{grace the occasion / the event with one's presence}} — honorer l'événement de sa présence.
    \item \eng{\textbf{logistical arrangements}} — dispositions logistiques / organisation matérielle.
    \item \eng{\textbf{foster unity / camaraderie / solidarity}} — favoriser l'unité / la camaraderie / la solidarité.
    \item \eng{\textbf{kindly confirm your attendance}} — veuillez confirmer votre présence (formule \eng{RSVP} standard).
    \item \eng{\textbf{anticipate a memorable event}} — s'attendre à / prévoir un événement mémorable.
    \item \eng{\textbf{student-led initiative}} — initiative portée par les élèves.
    \item \eng{\textbf{request the honour of your presence}} — solliciter l'honneur de votre présence (très protocolaire).
\end{itemize}
\end{definition}

\begin{attention}[Faux amis — \textbf{PIÈGES MAJEURS pour francophones}]
\noindent Ne confondez jamais le mot anglais et son « ami » français trompeur.
\begin{center}
\begin{tabularx}{\linewidth}{|X|X|X|}
\hline
\rowcolor{poporangeL} \textbf{Mot anglais (\eng{English})} & \textbf{Sens réel (\fr{Real Meaning})} & \textbf{Piège à éviter (\fr{False Friend})} \\
\hline
\eng{actually} & \fr{en fait, en réalité, vraiment} & \fr{actuellement} $\rightarrow$ \eng{currently, at present} \\
\eng{assist (sb)} & \fr{aider, assister qqn} & \fr{assister à} $\rightarrow$ \eng{attend, be present at} \\
\eng{entertainment} & \fr{divertissement, spectacle, animation} & \fr{entretien} $\rightarrow$ \eng{maintenance, interview} \\
\eng{eventually} & \fr{finalement, à la fin, après tout} & \fr{éventuellement} $\rightarrow$ \eng{possibly, potentially} \\
\eng{sensible} & \fr{raisonnable, sensé, judicieux} & \fr{sensible} $\rightarrow$ \eng{sensitive} \\
\eng{pretend} & \fr{faire semblant, simuler} & \fr{prétendre} $\rightarrow$ \eng{claim, assert} \\
\eng{realise} & \fr{se rendre compte, prendre conscience ; concrétiser} & \fr{réaliser (un film, un projet)} $\rightarrow$ \eng{achieve, shoot (film), carry out} \\
\eng{venue} & \fr{lieu (d'un événement)} & \fr{venue (venue d'arrivée)} $\rightarrow$ \eng{arrival} \\
\hline
\end{tabularx}
\end{center}
\end{attention}

\begin{definition}[Expressions idiomatiques — \eng{Idioms pour l'oral informel (Mates)}]
\noindent À utiliser \textbf{uniquement} entre pairs (\eng{mates, classmates}), jamais dans une lettre au Proviseur.
\begin{itemize}
    \item \eng{\textbf{Bring a plate}} — Apporter un plat (pour un \eng{potluck}). \eng{« It's a potluck, so bring a plate! »} (« C'est un repas partagé, apporte un plat ! »)
    \item \eng{\textbf{The more the merrier}} — Plus on est de fous, plus on rit. \eng{« Invite the whole class, the more the merrier! »}
    \item \eng{\textbf{Chip in}} — Contribuer (argent, effort, idée). \eng{« Everyone chips in 1 000 francs for drinks. »}
    \item \eng{\textbf{Have a blast}} — S'éclater, passer un moment génial. \eng{« We're gonna have a blast! »}
    \item \eng{\textbf{Sort out}} — Régler, organiser, trouver une solution. \eng{« Have you sorted out the venue? »}
    \item \eng{\textbf{Good shout}} — Bonne idée, bon plan (argot britannique/camerounais). \eng{« Courtyard? Good shout! »}
    \item \eng{\textbf{Show up}} — Se pointer, arriver (souvent inopinément). \eng{« Hope the Principal shows up. »}
\end{itemize}
\end{definition}

\begin{propriete}[Registres de langue — \eng{Register Shift: Informel vs Formel}]
\noindent Le passage d'un registre à l'autre est une compétence sociolinguistique clé (niveau B1/B2).
\begin{center}
\begin{tabularx}{\linewidth}{|X|X|}
\hline
\rowcolor{popgreenL} \textbf{Informel — \eng{Informal (Mates / Binôme A)}} & \textbf{Formel — \eng{Formal (Authorities / Binôme B)}} \\
\hline
Contractions : \eng{we'll, let's, it's, can't, gonna, wanna} & Formes pleines : \eng{we will, let us, it is, cannot, going to, want to} \\
Vocabulaire germanique/courant : \eng{party, bash, food, place, guys, kids, fix} & Vocabulaire roman/soutenu : \eng{festival, refreshments, venue, distinguished guests, arrange} \\
Verbes à particule (\eng{phrasal verbs}) : \eng{sort out, chip in, show up, set up, call off} & Verbes latins (\eng{latinate verbs}) : \eng{arrange, contribute, attend, organise, cancel} \\
Phrases elliptiques / hémicycles : \eng{Date sorted? Venue booked?} & Phrases complètes, syntaxe soignée : \eng{Has the date been finalised? Has the venue been booked?} \\
Marqueurs discursifs : \eng{Hey, mate, right?, yeah, cool, obviously, sweet} & Marqueurs discursifs : \eng{Dear Sir, furthermore, therefore, kindly, we would be grateful} \\
Pronoms : \eng{we, you, I} (familier) & Pronoms / Formules : \eng{we, you (formal), on behalf of..., the undersigned} \\
\hline
\end{tabularx}
\end{center}
\end{propriete}

\begin{propriete}[Familles de mots, orthographe et prononciation — \eng{Word Families, Spelling \& Pronunciation}]
\noindent Travaillez la dérivation pour enrichir votre lexique (préparation Bac : \eng{Word Formation}).
\begin{center}
\begin{tabular}{|l|l|l|l|l|}
\hline
\rowcolor{poppurpleL} \textbf{Verbe} & \textbf{Nom (action/état)} & \textbf{Nom (agent)} & \textbf{Adjectif} & \textbf{Adverbe} \\
\hline
\eng{celebrate} & \eng{celebration} & \eng{celebrant} & \eng{celebratory} & \eng{celebratorily} \\
\eng{organise} & \eng{organisation} & \eng{organiser} & \eng{organised} & \eng{organisedly} \\
\eng{invite} & \eng{invitation} & \eng{invitee} & \eng{inviting} & \eng{invitingly} \\
\eng{contribute} & \eng{contribution} & \eng{contributor} & \eng{contributory} & --- \\
\eng{attend} & \eng{attendance} & \eng{attendee} & \eng{attentive} & \eng{attentively} \\
\eng{entertain} & \eng{entertainment} & \eng{entertainer} & \eng{entertaining} & \eng{entertainingly} \\
\eng{festive} (adj) & \eng{festivity / festivities} & --- & \eng{festive} & \eng{festively} \\
\hline
\end{tabular}
\end{center}

\begin{itemize}
    \item \textbf{Orthographe (UK) :} \eng{organise} (pas \eng{organize}), \eng{realise}, \eng{finalise}, \eng{programme} (nom), \eng{honour}, \eng{favour}, \eng{labour}.
    \item \textbf{Prononciation clé (notation française approximative) :}
    \begin{itemize}
        \item \eng{venue} : « vé-niou » /"vEnju:/
        \item \eng{courtyard} : « courte-yarde » /"kO:tja:d/
        \item \eng{refreshments} : « re-frech-ments » /rI"freSm@nts/
        \item \eng{logistics} : « lo-jis-tics » /l@"dZIstIks/
        \item \eng{RSVP} : se lit lettre par lettre /"A: es vi: "pi:/
        \item \eng{eru} : « é-rou » (plat local) /"eru/
        \item \eng{jollof} : « djo-lof » /"dZ@l@f/
    \end{itemize}
\end{itemize}
\end{propriete}

\begin{methode}[Rédiger une lettre formelle d'invitation — \eng{Formal Invitation Letter Format (UK Standard)}]
\begin{enumerate}
    \item \textbf{En-tête expéditeur (haut droite) :} Adresse complète de l'école + Date (\eng{10th June 2025} — pas \eng{06/10/25}).
    \item \textbf{Destinataire (gauche, sous la date) :} Titre + Nom/Fonction + Adresse.
    \item \textbf{Objet (\eng{Subject / RE:}) :} Gras, centré ou aligné gauche. \eng{RE: Invitation to End-of-Year Class Festival — Form 1A}.
    \item \textbf{Salutation :} \eng{Dear Sir,} (si nom inconnu) ou \eng{Dear Mr [Nom],} (si connu). \textbf{Pas de prénom.}
    \item \textbf{Corps (3 paragraphes, saut de ligne entre chacun) :}
    \begin{itemize}
        \item \eng{Opening} : \eng{On behalf of..., we have the honour to extend an invitation to...}
        \item \eng{Details} : \eng{The event aims to... The festivities include... Logistical arrangements are finalised...}
        \item \eng{Closing} : \eng{We kindly request you to confirm your attendance by [date]... We anticipate...}
    \end{itemize}
    \item \textbf{Formule de clôture (\eng{Complimentary close}) :}
    \begin{itemize}
        \item \eng{Yours faithfully,} (si salutation \eng{Dear Sir/Madam,}).
        \item \eng{Yours sincerely,} (si salutation \eng{Dear Mr/Mrs [Nom],}).
    \end{itemize}
    \item \textbf{Signature :} Manuscrit + Nom imprimé en majuscules + Titre (\eng{Class Prefect, Form 1A}).
\end{enumerate}
\end{methode}

\courssub{Production modèle et analyse commentée}
\noindent Voici une production attendue pour un groupe performant (niveau A2+/B1). Les commentaires en \textbf{gras} expliquent les choix linguistiques.

\begin{exoresolu}[Production modèle — Groupe \eng{"Unity 1A"}]
\textbf{ÉTAPE 1 — DIALOGUE INFORMEL (Binôme A : Ngwa \& Fon)}
\eng{
\textbf{Ngwa:} Hey \eng{mate}, the end-of-year \eng{bash} is in two weeks. Have you \eng{sorted out} the \eng{venue} yet? \\
\textbf{Fon:} Not really. The main hall is taken for exams. I was thinking the \eng{courtyard} — open air, plenty of space. \\
\textbf{Ngwa:} \eng{Good shout}! That works. What about \eng{refreshments}? We said \eng{potluck}, right? \eng{Bring a plate}? \\
\textbf{Fon:} Exactly. \eng{The more the merrier}! Everyone brings a local dish — \eng{eru}, \eng{jollof}, \eng{puff-puff}... And we \eng{chip in} 500 francs for drinks. \\
\textbf{Ngwa:} Sweet. \eng{Entertainment}? We need a playlist and maybe a sketch. \\
\textbf{Fon:} I'll ask the drama club. Oh, and we \eng{must} invite the Principal. Formal letter, no slang. \\
\textbf{Ngwa:} \eng{Obviously}. I'll draft it tonight. We're gonna \eng{have a blast}!
}
\tcblower
\textbf{Commentaire linguistique (Oral / Binôme A) :}
\begin{itemize}
    \item \textbf{Registre informel authentique :} Contractions (\eng{we're, gonna, I'll}), vocabulaire familier (\eng{mate, bash, sweet, obviously, cool}), verbes à particule (\eng{sort out, chip in}).
    \item \textbf{Champ lexical (6 items) :} \eng{bash, venue, courtyard, refreshments, potluck, entertainment, Principal}.
    \item \textbf{Idiomes (5 utilisés) :} \eng{"Good shout"}, \eng{"The more the merrier"}, \eng{"Bring a plate"}, \eng{"Chip in"}, \eng{"Have a blast"} — densité excellente.
    \item \textbf{Ancrage camerounais :} \eng{eru, jollof, puff-puff} (plats locaux), \eng{courtyard} (cour de récréation typique), \eng{500 francs} (monnaie locale).
    \item \textbf{Prononciation cible :} \eng{venue} /"vEnju:/, \eng{courtyard} /"kO:tja:d/, \eng{eru} /"eru/, \eng{jollof} /"dZ@l@f/.
\end{itemize}

\textbf{ÉTAPE 2 — LETTRE FORMELLE (Binôme B : Class Prefect \& Deputy)}
\eng{
GOVERNMENT BILINGUAL HIGH SCHOOL BUEA \\
P.O. Box 45, Buea, Southwest Region \\
10th June 2025 \\[0.5em]

THE PRINCIPAL \\
GOVERNMENT BILINGUAL HIGH SCHOOL BUEA \\
P.O. Box 45, BUEA \\[0.5em]

\textbf{RE: INVITATION TO END-OF-YEAR CLASS FESTIVAL — FORM 1A} \\[0.5em]

Dear Sir, \\[0.5em]

On behalf of the students of Form 1A, we have the honour to \eng{extend a formal invitation} to you to \eng{grace the occasion} of our End-of-Year Class Festival, scheduled for \eng{Friday, 27th June 2025 at 2:00 p.m.} in the \eng{school courtyard}. \\[0.5em]

This \eng{student-led initiative} aims to \eng{foster unity} and celebrate our academic achievements in a joyful atmosphere. The \eng{festivities} will include a cultural display of traditional dances, a short drama sketch on school life, and a \eng{potluck} lunch featuring local delicacies. We have finalised all \eng{logistical arrangements} with the School Bursar and the Discipline Master to ensure a smooth and safe event. Your presence would be a great source of encouragement for the entire class. \\[0.5em]

We \eng{kindly request you to confirm your attendance} by \eng{Wednesday, 25th June} to facilitate the seating protocol. We \eng{anticipate a memorable event} and hope you will honour us with your presence. \\[0.5em]

Yours faithfully, \\[1.5em]

\textbf{[Signature]} \\
NGWA Samuel \\
Class Prefect, Form 1A \\[1em]

\textbf{[Signature]} \\
FON Grace \\
Deputy Class Prefect, Form 1A
}
\tcblower
\textbf{Commentaire linguistique (Écrit / Binôme B) :}
\begin{itemize}
    \item \textbf{Format lettre administrative parfait :} Adresses alignées (expéditeur droite, destinataire gauche), date format UK (\eng{10th June 2025}), objet \eng{RE:}, salutation \eng{Dear Sir,}, clôture \eng{Yours faithfully,} (cohérence \eng{Dear Sir} $\rightarrow$ \eng{Yours faithfully}).
    \item \textbf{Zéro contraction :} \eng{we have, it is, we will, cannot, we would} (implicite dans \eng{we would be grateful} si utilisé).
    \item \textbf{Collocations formelles (6 utilisées, objectif dépassé) :} \eng{extend a formal invitation}, \eng{grace the occasion}, \eng{student-led initiative}, \eng{foster unity}, \eng{logistical arrangements}, \eng{kindly request you to confirm your attendance}, \eng{anticipate a memorable event}.
    \item \textbf{Faux amis évités :} Pas de \eng{actually} pour « actuellement » (utilisation de \eng{scheduled for}), pas de \eng{assist} pour « assister » (utilisation de \eng{grace/attendance/presence}), pas de \eng{entertainment} pour « entretien ».
    \item \textbf{Vocabulaire précis et contextuel :} \eng{delicacies} (mets délicats), \eng{seating protocol} (protocole de placement), \eng{student-led initiative}, \eng{Bursar}, \eng{Discipline Master}.
    \item \textbf{Cohérence inter-textuelle :} La lettre reprend les infos du dialogue (date, lieu, potluck, drama club) mais transposées dans le registre approprié.
\end{itemize}
\end{exoresolu}

\courssub{Bilan et ancrage des acquis}
\noindent Cette activité d'intégration a permis de vérifier l'appropriation des compétences lexicales et sociolinguistiques de la Leçon 3.

\begin{aretenir}[Ce qu'il faut retenir de la Leçon 3]
\begin{itemize}
    \item \textbf{Le contexte dicte le registre :} On ne s'adresse pas à son \eng{mate} comme au \eng{Principal}. Le passage de l'informel (\eng{phrasal verbs, slang, contractions, ellipses}) au formel (\eng{latinate verbs, full forms, fixed polite formulas, passive voice}) est une compétence clé du niveau B1/B2 et un attendus fort du Baccalauréat (épreuve d'expression écrite).
    \item \textbf{Les collocations sont la marque du niveau avancé :} Dire \eng{extend an invitation} au lieu de \eng{make/give an invitation}, ou \eng{grace the occasion} au lieu de \eng{come to the party}, montre une maîtrise lexicale fine (compétence \eng{lexical resource}).
    \item \textbf{Les faux amis sont des pièges réels à l'écrit comme à l'oral :} Écrire \eng{"We actually prepare the party"} pour \fr{"Nous préparons actuellement la fête"} change radicalement le sens (\eng{"En fait, nous préparons la fête..."}). La vigilance est de mise sur \eng{actually, assist, entertainment, eventually, sensible, pretend, realise}.
    \item \textbf{L'anglais camerounais standard existe :} L'intégration de réalités locales (\eng{eru, jollof, puff-puff, courtyard, bendskin rhythm, bursar, discipline master, PTA, chief}) ancre l'apprentissage dans le vécu de l'élève tout en exigeant un anglais standard (UK norms) pour l'écrit formel.
    \item \textbf{La dérivation (Word Formation) est un levier de score au Bac :} Maîtriser \eng{organise $\rightarrow$ organisation/organiser}, \eng{invite $\rightarrow$ invitation/invitee}, \eng{contribute $\rightarrow$ contribution/contributor} permet de varier le lexique en composition.
\end{itemize}
\end{aretenir}

\begin{experience}[Prolongements possibles — \eng{Extension Tasks (Différenciation)}]
\begin{itemize}
    \item \textbf{Expression orale — \eng{Ceremonial Speech (Semi-formel) :}} Rédiger et prononcer le \eng{welcome speech} (1 min) du \eng{Class Prefect} le jour J pour accueillir le Proviseur. Registre \eng{semi-formel / ceremonial} : \eng{"Distinguished guests, ladies and gentlemen, on behalf of Form 1A..."}.
    \item \textbf{Compréhension de l'oral — \eng{Listening / Note-taking :}} Enregistrer les meilleurs dialogues des groupes. Diffuser en classe pour un exercice de \eng{note-taking} : relever date, lieu, 3 activités, 1 problème logistique mentionné, 2 idiomes entendus.
    \item \textbf{Projet interdisciplinaire (EMC / Vie scolaire) :} Organiser \textbf{vraiment} la fête de classe (avec accord de l'administration) en utilisant la lettre rédigée comme document officiel. Impliquer le \eng{Bursar} (Intendant) pour la réservation de la \eng{venue}.
    \item \textbf{Écriture créative — \eng{Article de journal scolaire :}} Rédiger un article (\eng{school magazine}) \textit{après} l'événement : \eng{"Form 1A End-of-Year Festival: A Resounding Success"} (passé, discours rapporté, vocabulaire des sentiments).
\end{itemize}
\end{experience}

\begin{savaistu}[Le saviez-vous ? — Culture scolaire anglophone au Cameroun]
Dans le système éducatif anglophone du Cameroun (héritage \eng{West Cameroon}), le \eng{Principal} est le chef d'établissement (équivalent du Proviseur francophone). Le \eng{Vice-Principal} gère souvent la discipline, l'emploi du temps et le remplacement des enseignants absents. Le \eng{Discipline Master} (ou \eng{Mistress}) est un poste clé, très respecté, chargé de l'ordre, de la ponctualité et du port de l'uniforme (\eng{school uniform}). Les grandes cérémonies officielles de fin d'année s'appellent \eng{Speech and Prize Giving Day} ; y sont invités les \eng{chiefs} (chefs traditionnels), le \eng{PTA} (\eng{Parent-Teacher Association}), le \eng{Divisional Officer} (Sous-préfet). Sur les cartons d'invitation officiels, la formule consacrée est invariablement : \eng{"We request the honour of your presence..."} ou \eng{"We invite you to grace the occasion..."}. L'expression \eng{"bring a plate"} pour les \eng{potlucks} est très courante dans les communautés anglophones (écoles, églises, quartiers).
\end{savaistu}

\newpage

\coursec{Méthodes et exercices résolus}

\coursec{Lesson 3 -- Vocabulary: Words and expressions related to interacting with mates, neighbours, adults and school authorities on community celebrations/festivals}

\courssub{Methods and Solved Exercises}

\begin{methode}[Identifier et classer le vocabulaire de base des célébrations communautaires]
\textbf{Objectif :} Maîtriser le lexique essentiel (noms, verbes, adjectifs) pour décrire une fête, un festival ou une cérémonie au Cameroun et l'organiser.
\begin{enumerate}
    \item \textbf{Lire / Écouter} le texte support en repérant les mots-clés : noms d'événements (\eng{festival, ceremony, celebration, gathering, holiday}), actions (\eng{organise, host, attend, participate, celebrate, decorate, invite}), acteurs (\eng{host, guest, participant, organiser, chief, principal, neighbour, mate}) et lieux/temps (\eng{venue, community hall, palace, school grounds, date, schedule}).
    \item \textbf{Classer} les mots par \textbf{champs lexicaux} dans un tableau (ex. : \textit{Preparation}, \textit{During the event}, \textit{People}, \textit{Atmosphere}).
    \item \textbf{Vérifier} la prononciation (accent tonique) et l'orthographe des mots transparents (\eng{ceremony / festival / decoration / invitation}).
    \item \textbf{Mémoriser} 5 à 7 collocations fortes : \eng{host a festival}, \eng{attend a ceremony}, \eng{deliver a speech}, \eng{grant permission}, \eng{make a toast}.
\end{enumerate}
\textbf{Réflexe :} Toujours associer le nom au verbe qui lui va naturellement (collocation).
\end{methode}

\begin{exoresolu}[Catégorisation et emploi du vocabulaire de base]
\textbf{Énoncé :} Le texte ci-dessous décrit la préparation de la fête de la Jeunesse dans un lycée de Buea. Complétez le tableau en classant les mots \textbf{surlignés} dans la bonne colonne. Ensuite, réécrivez les phrases 1 à 3 en remplaçant les parties soulignées par une \textbf{collocation} appropriée du tableau.

\begin{textebox}[School Magazine Article -- Youth Day Preparations]
The \textbf{organising committee} met yesterday at the \textbf{venue}, the school \textbf{grounds}. The \textbf{principal} \textbf{granted permission} for the cultural \textbf{display}. Students will \textbf{perform} traditional dances and \textbf{deliver speeches}. The \textbf{guests of honour} include the \textbf{Division Officer} and local \textbf{chiefs}. Everyone hopes to \textbf{enjoy} a colourful \textbf{atmosphere}.
\end{textebox}

\begin{center}
\begin{tabularx}{\linewidth}{|X|X|X|X|}
\hline
\rowcolor{popblueL} \textbf{People / Roles} & \textbf{Places} & \textbf{Actions / Verbs} & \textbf{Event Features} \\ \hline
 &  &  &  \\ \hline
 &  &  &  \\ \hline
 &  &  &  \\ \hline
\end{tabularx}
\end{center}

\textbf{Phrases à transformer :}
\begin{enumerate}
    \item The headmaster \underline{said yes} for the party. \eng{(formal register)}
    \item The students will \underline{do} traditional dances on stage.
    \item The VIPs will \underline{make} speeches during the ceremony.
\end{enumerate}

\tcblower
\textbf{Solution détaillée :}

\textbf{1. Classement du vocabulaire (lecture active) :}
\begin{center}
\begin{tabularx}{\linewidth}{|X|X|X|X|}
\hline
\rowcolor{popblueL} \textbf{People / Roles} & \textbf{Places} & \textbf{Actions / Verbs} & \textbf{Event Features} \\ \hline
organising committee & venue & granted permission & cultural display \\
principal & school grounds & perform & atmosphere \\
guests of honour &  & deliver speeches &  \\
Division Officer &  & enjoy &  \\
chiefs &  &  &  \\ \hline
\end{tabularx}
\end{center}
\textit{Justification :} \eng{Organising committee, principal, guests of honour, Division Officer, chiefs} désignent des personnes. \eng{Venue, school grounds} sont des lieux. \eng{Granted permission, perform, deliver speeches, enjoy} sont des actions (verbes ou locutions verbales). \eng{Cultural display, atmosphere} caractérisent l'événement.

\textbf{2. Transformation par collocations (registre formel) :}
\begin{enumerate}
    \item \eng{The Principal \textbf{granted permission} for the celebration.} \\
    \textit{Justification :} \eng{``Said yes''} est trop familier. \eng{Grant permission} est la collocation administrative standard (autoriser officiellement).
    \item \eng{The students will \textbf{perform} traditional dances on stage.} \\
    \textit{Justification :} \eng{Perform} (se produire, exécuter) est le verbe précis pour un spectacle/danse. \eng{Do} est trop vague.
    \item \eng{The VIPs will \textbf{deliver speeches} during the ceremony.} \\
    \textit{Justification :} \eng{Deliver a speech} est la collocation figée pour ``faire un discours'' (formel). \eng{Make a speech} est possible mais moins solennel ; \eng{make} seul est incorrect ici.
\end{enumerate}
\textbf{Conclusion :} Le choix du verbe précis (\eng{grant, perform, deliver}) marque le registre formel attendu avec les autorités.
\end{exoresolu}

\begin{methode}[Adapter son registre de langue selon l'interlocuteur (Mates vs Authorities)]
\textbf{Objectif :} Choisir les formules de politesse, les modaux et le vocabulaire adaptés à quatre cibles : \textbf{mates} (camarades), \textbf{neighbours} (voisins), \textbf{adults} (adultes/parents), \textbf{school authorities} (autorités scolaires/admin).
\begin{enumerate}
    \item \textbf{Identifier} l'interlocuteur et le contexte (informel, semi-formel, formel).
    \item \textbf{Sélectionner} la structure grammaticale :
    \begin{itemize}
        \item \textbf{Mates :} Impératif, \eng{Let's...}, \eng{Can I...?}, \eng{Wanna...?}, prénom, tutoiement (\eng{you}).
        \item \textbf{Neighbours / Adults :} \eng{Could you...?}, \eng{Would you mind...?}, \eng{Mr/Mrs + Nom}, \eng{please}, remerciements appuyés.
        \item \textbf{Authorities :} \eng{May I...?}, \eng{I would like to request...}, \eng{We would be honoured if...}, \eng{Sir/Madam/Title + Name}, phrases complètes, pas de contractions (\eng{do not} vs \eng{don't}).
    \end{itemize}
    \item \textbf{Adapter} le lexique : \eng{party/bash} (mates) $\rightarrow$ \eng{gathering/event} (neighbours) $\rightarrow$ \eng{ceremony/function} (authorities).
    \item \textbf{Relire} pour vérifier la cohérence (pas de mélange de registres dans une même phrase).
\end{enumerate}
\textbf{Structure à retenir :} \eng{Would it be possible for [us] to + V-infinitif ?} (très poli, authorities).
\end{methode}

\begin{exoresolu}[Réécriture selon le destinataire : Invitation à la fête de fin d'année]
\textbf{Énoncé :} Vous organisez la \eng{End-of-Year Party} de votre classe. Réécrivez le message principal \textbf{« Come to our party on Friday at 4 pm »} pour les quatre destinataires ci-dessous. Respectez le registre, la politesse et le vocabulaire adaptés.

\begin{enumerate}
    \item Votre meilleur ami (WhatsApp).
    \item Votre voisin, M. Fouda, père d'un élève, que vous vouvoyez.
    \item Le Proviseur (Principal) du lycée, M. Essomba.
    \item Le Délégué Départemental de l'Éducation Secondaire (Divisional Delegate), Mme Ngo.
\end{enumerate}

\tcblower
\textbf{Solution détaillée :}

\begin{center}
\begin{tabularx}{\linewidth}{|X|X|}
\hline
\rowcolor{popblueL} \textbf{Destinataire} & \textbf{Message rédigé (Registre \& Vocabulaire)} \\ \hline
\textbf{1. Mate (Ami)} & \eng{Hey! Don't forget our end-of-year bash on Friday at 4 pm sharp. See u there!} \\
\rowcolor{popcream} \textit{Analyse :} Familier (\eng{Hey, bash, u, sharp}). Impératif direct. \\ \hline
\textbf{2. Neighbour (M. Fouda)} & \eng{Good evening, Mr Fouda. We would be delighted if you could join us for the students' end-of-year gathering this Friday at 4 pm.} \\
\rowcolor{popcream} \textit{Analyse :} Semi-formel (\eng{delighted, could join, gathering}). Conditionnel \eng{would be delighted if you could}. \eng{Mr + Nom}. \\ \hline
\textbf{3. Authority (Principal Essomba)} & \eng{Dear Principal Essomba, We respectfully invite you to attend the End-of-Year Ceremony scheduled for Friday at 4 pm on the school grounds. Your presence would honour us.} \\
\rowcolor{popcream} \textit{Analyse :} Formel (\eng{respectfully invite, attend, ceremony, scheduled, honour us}). Pas de contractions. Vocabulaire précis (\eng{ceremony, grounds}). \\ \hline
\textbf{4. High Authority (Delegate Ngo)} & \eng{Madam Divisional Delegate, On behalf of the organising committee, I would like to request the honour of your presence at the End-of-Year Function this Friday at 4 pm. We would be grateful if you could grace the occasion with a keynote address.} \\
\rowcolor{popcream} \textit{Analyse :} Très formel / Administratif (\eng{On behalf of, request the honour, function, grace the occasion, keynote address}). Titre complet \eng{Madam Divisional Delegate}. Structure \eng{I would like to request...}. \\ \hline
\end{tabularx}
\end{center}

\textbf{Points clés corrigés :}
\begin{itemize}
    \item \eng{Party} $\rightarrow$ \eng{bash / gathering / ceremony / function} (échelle de formalité).
    \item \eng{Come} $\rightarrow$ \eng{join us / attend / grace the occasion}.
    \item Pronoms/Modaux : \eng{Can/Will} (mates) $\rightarrow$ \eng{Could/Would} (neighbours) $\rightarrow$ \eng{May/Might/Would be honoured} (authorities).
\end{itemize}
\end{exoresolu}

\begin{methode}[Maîtriser les collocations et expressions idiomatiques du thème]
\textbf{Objectif :} Utiliser des blocs de mots figés (collocations) et des idiomes pour parler de l'organisation, du déroulement et de l'ambiance des fêtes de façon naturelle.
\begin{enumerate}
    \item \textbf{Apprendre les collocations par paires Verbe + Nom} (ne pas apprendre le nom seul) :
    \eng{organise/host an event}, \eng{attend/join a celebration}, \eng{make/give a speech/toast}, \eng{grant/seek permission}, \eng{decorate the venue}, \eng{entertain the guests}.
    \item \textbf{Repérer les expressions idiomatiques} (sens non compositionnel) :
    \eng{have a blast} (s'éclater), \eng{the life of the party} (l'âme de la fête), \eng{pull out all the stops} (mettre les petits plats dans les grands), \eng{let one's hair down} (se détendre), \eng{a red-letter day} (un jour mémorable).
    \item \textbf{Contextualiser :} Écrire une phrase personnelle pour chaque collocation/idiome.
    \item \textbf{Attention aux prépositions :} \eng{invite someone \textbf{to} an event}, \eng{congratulate someone \textbf{on} an occasion}, \eng{participate \textbf{in} an activity}.
\end{enumerate}
\textbf{Réflexe Bac :} Dans une composition, remplacer \eng{``make a party''} (calque fr) par \eng{``host/throw a party''} ou \eng{``organise an event''}.
\end{methode}

\begin{exoresolu}[Complétion et correction de collocations / idiomes]
\textbf{Énoncé :}
\textbf{A.} Complétez les phrases avec le verbe qui colle au nom (collocation). Choisissez parmi : \eng{host, deliver, grant, pull out, make, attend}.
\begin{enumerate}
    \item The traditional ruler will \_\_\_\_\_\_ a grand \_\_\_\_\_\_ to celebrate the Ngondo festival.
    \item Students must \_\_\_\_\_\_ \_\_\_\_\_\_ from the Principal before leaving the campus.
    \item Our class decided to \_\_\_\_\_\_ \_\_\_\_\_\_ all the \_\_\_\_\_\_ for the cultural night.
    \item The Head Girl will \_\_\_\_\_\_ a \_\_\_\_\_\_ on behalf of the students.
    \item We hope many parents will \_\_\_\_\_\_ the \_\_\_\_\_\_.
\end{enumerate}

\textbf{B.} Corrigez les erreurs de collocation ou d'idiome dans les phrases suivantes (style composition).
\begin{enumerate}
    \item We \textbf{made a big party} for the National Day.
    \item The Minister \textbf{said a discourse} during the ceremony.
    \item Everyone \textbf{passed a good moment} at the festival.
    \item She was \textbf{the soul of the party}, dancing all night.
\end{enumerate}

\tcblower
\textbf{Solution détaillée :}

\textbf{A. Collocations Verbe + Nom (réponses attendues) :}
\begin{enumerate}
    \item The traditional ruler will \textbf{host a grand \eng{ceremony/celebration}}... (\eng{Host a ceremony} = organiser/présider officiellement).
    \item Students must \textbf{seek/obtain permission}... (\eng{Seek permission} = demander l'autorisation ; \eng{grant} = accorder, sujet = autorité).
    \item Our class decided to \textbf{pull out all the stops}... (\textbf{Idiome} : \eng{pull out all the stops} = ne rien négliger, tout faire pour réussir).
    \item The Head Girl will \textbf{deliver a \eng{speech/address}}... (\eng{Deliver a speech} = prononcer un discours).
    \item We hope many parents will \textbf{attend the \eng{ceremony/function/event}}... (\eng{Attend} = assister/être présent).
\end{enumerate}

\textbf{B. Correction d'erreurs fréquentes (Francophonismes) :}
\begin{enumerate}
    \item \eng{We \textbf{hosted/threw/organised a big party} for the National Day.}
    \textit{Explication :} \eng{Make a party} n'existe pas. \eng{Throw a party} (familier), \eng{host/organise} (standard/formel).
    \item \eng{The Minister \textbf{delivered/gave a speech/address} during the ceremony.}
    \textit{Explication :} \eng{Say a discourse} $\times$. \eng{Discourse} = discours au sens abstrait/linguistique. \eng{Speech/Address} = allocution. Verbe : \eng{deliver/give}.
    \item \eng{Everyone \textbf{had a great time / enjoyed themselves} at the festival.}
    \textit{Explication :} \eng{Pass a good moment} $\times$ (calque \textit{passer un bon moment}). \eng{Have a good time / Enjoy oneself}.
    \item \eng{She was \textbf{the life of the party}, dancing all night.}
    \textit{Explication :} \eng{Soul of the party} $\times$ (calque \textit{l'âme de la fête}). L'idiome anglais est \eng{the life of the party} (ou \eng{the life and soul of the party}).
\end{enumerate}
\end{exoresolu}

\begin{methode}[Éviter les faux-amis et pièges de traduction (Français $\rightarrow$ Anglais)]
\textbf{Objectif :} Identifier et corriger les erreurs lexicales classiques des francophones liées aux célébrations et à l'interaction sociale.
\begin{enumerate}
    \item \textbf{Lister les paires « Faux-amis / Vrais sens »} du thème :
    \begin{center}
    \begin{tabular}{|l|l|l|}
    \hline
    \rowcolor{poporangeL} \textbf{Mot anglais (Piège)} & \textbf{Sens réel (FR)} & \textbf{Mot anglais correct (FR visé)} \\ \hline
    \eng{Actually} & En fait / En réalité & \eng{Currently / At present} (Actuellement) \\ \hline
    \eng{Assist (to)} & Aider / Secourir & \eng{Attend} (Assister à un événement) \\ \hline
    \eng{Disguise} & Déguiser (transitif) & \eng{Dress up / Wear a costume} (Se déguiser) \\ \hline
    \eng{Costume} & Costume (homme) / Déguisement & \eng{Suit} (Costume homme) / \eng{Outfit} (Tenue) \\ \hline
    \eng{Festival} & Festival (souvent artistique) & \eng{Holiday / Feast / Celebration} (Fête nationale/religieuse) \\ \hline
    \eng{Ceremony} & Cérémonie (formel) & \eng{Party / Gathering} (Fête informelle) \\ \hline
    \eng{Authority} & Autorité (pouvoir) & \eng{Authorities} (pl. = administrations) \\ \hline
    \eng{Manifestation} & Manifestation (démonstration) & \eng{Demonstration / Protest / Rally} (Manifestation politique) \\ \hline
    \eng{Animation} & Animation (dessin/activité) & \eng{Entertainment / Activities / Hosting} (Animation de soirée) \\ \hline
    \end{tabular}
    \end{center}
    \item \textbf{Surveiller les prépositions pièges :} \eng{Invite \textbf{to}} (pas \eng{for}), \eng{Congratulate \textbf{on}} (pas \eng{for}), \eng{Participate \textbf{in}} (pas \eng{to}).
    \item \textbf{Tester la phrase} : « Est-ce que ce mot veut dire ce que je pense en français ? » $\rightarrow$ Dictionnaire / Glossaire.
\end{enumerate}
\end{methode}

\begin{exoresolu}[Chasse aux faux-amis et correction de calques]
\textbf{Énoncé :} Le paragraphe suivant a été écrit par un élève de 1ère A. Il contient \textbf{7 erreurs lexicales} (faux-amis, calques, mauvais choix de registre). Identifiez-les, soulignez-les et réécrivez le paragraphe correctement.

\begin{textebox}[Student's Draft]
``Next Saturday, our school will \textbf{organise a manifestation} for the \textbf{festival} of Tabaski. The Principal will \textbf{assist} the ceremony. He will \textbf{say a discourse} about peace. Many parents are \textbf{actually} preparing food. The \textbf{animation} will be done by a famous DJ. Students must wear their \textbf{costumes}. It will be a \textbf{sympathetic} day.''
\end{textebox}

\tcblower
\textbf{Solution détaillée :}

\textbf{Analyse erreur par erreur :}

\begin{center}
\begin{tabularx}{\linewidth}{|X|X|X|}
\hline
\rowcolor{poporangeL} \textbf{Erreur (Texte élève)} & \textbf{Type / Explication} & \textbf{Correction} \\ \hline
\eng{organise a \textbf{manifestation}} & \textbf{Faux-ami majeur.} \eng{Manifestation} = démo politique/protestation. Pour fête scolaire : \eng{event, celebration, function}. & \eng{organise an \textbf{event / celebration}} \\ \hline
\eng{\textbf{festival} of Tabaski} & \textbf{Mauvais choix lexical.} \eng{Festival} = événement artistique/culturel programmé (ex. Festival de Cannes). Fête religieuse = \eng{Feast / Holiday / Celebration}. & \eng{\textbf{Feast / Celebration} of Tabaski} \\ \hline
\eng{\textbf{assist} the ceremony} & \textbf{Faux-ami structurel.} \eng{Assist} = aider (qqn). \eng{Assist at} = assister (être présent) vieilli. Verbe standard : \eng{Attend}. & \eng{\textbf{attend} the ceremony} \\ \hline
\eng{\textbf{say a discourse}} & \textbf{Calque + Faux-ami.} \eng{Say} $\neq$ faire un discours. \eng{Discourse} = discours théorique/linguistique. Collocation : \eng{deliver/give a speech/address}. & \eng{\textbf{deliver an address / give a speech}} \\ \hline
\eng{are \textbf{actually} preparing} & \textbf{Faux-ami temporel.} \eng{Actually} = « en fait / en réalité » (opposition). « Actuellement » = \eng{Currently / At present / Right now}. & \eng{are \textbf{currently / right now} preparing} \\ \hline
\eng{The \textbf{animation} will be done} & \textbf{Calque / Faux-ami.} \eng{Animation} = dessin animé / action d'animer. Pour « animation musicale/soirée » : \eng{Entertainment / The entertainment / The programme}. & \eng{The \textbf{entertainment} will be provided / handled} \\ \hline
\eng{wear their \textbf{costumes}} & \textbf{Faux-ami / Registre.} \eng{Costume} = déguisement (Halloween) ou costume de scène. Tenue habillée = \eng{Outfit / Attire / Traditional wear / Best clothes}. & \eng{wear traditional \textbf{attire / outfits / best clothes}} \\ \hline
\eng{\textbf{sympathetic} day} & \textbf{Faux-ami adjectif.} \eng{Sympathetic} = compatissant, qui a de la sympathie (pour qqn). « Sympa / Agréable » = \eng{Nice / Pleasant / Enjoyable / Lovely}. & \eng{an \textbf{enjoyable / lovely} day} \\ \hline
\end{tabularx}
\end{center}

\textbf{Paragraphe corrigé (Registre formel/scolaire) :}
\begin{textebox}[Corrected Version]
``Next Saturday, our school will \textbf{host a celebration} for the \textbf{Feast} of Tabaski. The Principal will \textbf{attend} the ceremony. He will \textbf{deliver an address} about peace. Many parents are \textbf{currently} preparing food. The \textbf{entertainment} will be \textbf{provided} by a famous DJ. Students must wear traditional \textbf{attire}. It will be an \textbf{enjoyable} day.''
\end{textebox}
\textbf{Note :} \eng{Sympathetic} est un piège très fréquent au Bac (confusion avec \textit{sympathique}). \eng{Actually} / \eng{Currently} est un classique des sujets de traduction.
\end{exoresolu}

\begin{methode}[Construire des familles de mots, orthographe et prononciation]
\textbf{Objectif :} Dériver les mots (nom, verbe, adjectif, adverbe, personne) pour enrichir son lexique et éviter les fautes d'orthographe/prononciation.
\begin{enumerate}
    \item \textbf{Partir du verbe de base} (souvent le plus court) : \eng{Celebrate, Organise, Invite, Decorate, Congratulate, Participate}.
    \item \textbf{Appliquer les suffixes productifs} :
    \begin{itemize}
        \item \eng{-tion / -sion} $\rightarrow$ Nom d'action (\eng{celebration, organisation, invitation, decoration, congratulation, participation}). \textbf{Orthographe :} \eng{celebrate} $\rightarrow$ \eng{celebration} (disparition du \eng{e} muet).
        \item \eng{-er / -or} $\rightarrow$ Agent (\eng{organiser, participant, celebrant, host}).
        \item \eng{-ive / -ory / -al} $\rightarrow$ Adjectif (\eng{celebratory, festive, ceremonial, official}).
        \item \eng{-ly} $\rightarrow$ Adverbe (\eng{officially, festively}).
    \end{itemize}
    \item \textbf{Travailler la prononciation (accent tonique)} : Le suffixe \eng{-tion} attire l'accent sur la syllabe précédente (\eng{cel-e-BRA-tion}, \eng{or-gan-i-SA-tion}, \eng{in-vi-TA-tion}). Marquer l'accent sur le papier (\eng{cé-lé-BRÉ-chn}).
    \item \textbf{Vérifier les doubles consonnes} : \eng{Occasion} (2 c, 1 s), \eng{Committee} (2 m, 2 t, 2 e), \eng{Accommodation} (2 c, 2 m).
\end{enumerate}
\textbf{Outil :} Tableau de dérivation à compléter pour chaque nouvelle racine.
\end{methode}

\begin{exoresolu}[Dérivation, orthographe et accent tonique]
\textbf{Énoncé :}
\textbf{A.} Complétez le tableau de dérivation. Attention aux changements d'orthographe et à la catégorie grammaticale.
\begin{center}
\begin{tabular}{|l|l|l|l|l|l|}
\hline
\rowcolor{poppurpleL} \textbf{Verbe} & \textbf{Nom (Action/État)} & \textbf{Nom (Agent)} & \textbf{Adjectif} & \textbf{Adverbe} & \textbf{Traduction (Verbe)} \\ \hline
\eng{Celebrate} &  &  &  &  & Célébrer \\ \hline
\eng{Organise} &  &  &  &  & Organiser \\ \hline
\eng{Invite} &  &  &  &  & Inviter \\ \hline
\eng{Decorate} &  &  &  &  & Décorer \\ \hline
\eng{Congratulate} &  &  &  &  & Féliciter \\ \hline
\eng{Authorise} &  &  &  &  & Autoriser \\ \hline
\end{tabular}
\end{center}

\textbf{B.} Placez l'accent tonique (') sur le mot en gras dans les phrases ci-dessous (règle des suffixes \eng{-tion}, \eng{-ive}, \eng{-al}).
\begin{enumerate}
    \item The \textbf{celebration} was a huge success.
    \item It is a \textbf{celebratory} meal.
    \item We need official \textbf{authorisation}.
    \item The \textbf{organiser} spoke to the \textbf{participants}.
\end{enumerate}

\textbf{C.} Corrigez l'orthographe de ces mots souvent mal écrits au Bac :
\eng{occassion, comittee, accomodation, refering, begining, speach}.

\tcblower
\textbf{Solution détaillée :}

\textbf{A. Tableau de dérivation complet :}
\begin{center}
\begin{tabular}{|l|l|l|l|l|l|}
\hline
\rowcolor{poppurpleL} \textbf{Verbe} & \textbf{Nom (Action)} & \textbf{Nom (Agent)} & \textbf{Adjectif} & \textbf{Adverbe} & \textbf{Traduction} \\ \hline
\eng{Celebrate} & \eng{Celebration} & \eng{Celebrant} & \eng{Celebratory} & \eng{Celebratorily*} & Célébrer \\ \hline
\eng{Organise} & \eng{Organisation} & \eng{Organiser} & \eng{Organisational} & \eng{Organisationally} & Organiser \\ \hline
\eng{Invite} & \eng{Invitation} & \eng{Invitee / Host} & \eng{Inviting} & \eng{Invitingly} & Inviter \\ \hline
\eng{Decorate} & \eng{Decoration} & \eng{Decorator} & \eng{Decorative} & \eng{Decoratively} & Décorer \\ \hline
\eng{Congratulate} & \eng{Congratulation} & \eng{Congratulator} & \eng{Congratulatory} & \eng{Congratulatorily} & Féliciter \\ \hline
\eng{Authorise} & \eng{Authorisation} & \eng{Authoriser} & \eng{Authorised / Authoritative} & \eng{Authoritatively} & Autoriser \\ \hline
\end{tabular}
\end{center}
\textit{Notes :} \eng{Celebratorily} est rare. \eng{Inviting} (adj) = accueillant. \eng{Authoritative} = faisant autorité (différent de \eng{authorised} = autorisé).

\textbf{B. Accent tonique (Règle : suffixes \eng{-tion}, \eng{-sion}, \eng{-ic}, \eng{-ial}, \eng{-itive} attirent l'accent sur la syllabe qui précède immédiatement) :}
\begin{enumerate}
    \item The \eng{cel-e-\textbf{BRA}-tion} was a huge success. (Accent sur \textbf{BRA})
    \item It is a \eng{cel-e-\textbf{BRA}-to-ry} meal. (Accent sur \textbf{BRA} - règle \eng{-tory})
    \item We need of-fi-ci-\textbf{AU}-tho-ri-sa-tion. (Accent sur \textbf{AU} - règle \eng{-tion})
    \item The \textbf{OR}-gan-i-ser spoke to the par-\textbf{TI}-ci-pants. (\eng{Organiser} : accent sur radical \textbf{OR} ; \eng{Participants} : accent sur \textbf{TI} avant \eng{-pant}).
\end{enumerate}

\textbf{C. Corrections orthographe (Pièges Bac) :}
\begin{center}
\begin{tabular}{|l|l|l|}
\hline
\rowcolor{popredL} \textbf{Erreur élève} & \textbf{Correct} & \textbf{Mnémotechnique / Règle} \\ \hline
\eng{occassion} & \eng{Occasion} & 2 \textbf{c}, 1 \textbf{s} (comme \eng{occasionnel}) \\ \hline
\eng{comittee} & \eng{Committee} & Double \textbf{m}, double \textbf{t}, double \textbf{e} (Com-mit-tee) \\ \hline
\eng{accomodation} & \eng{Accommodation} & Double \textbf{c}, double \textbf{m} (Ac-com-mo-da-tion) \\ \hline
\eng{refering} & \eng{Referring} & Double \textbf{r} (verbe monosyllabe accentué \eng{refer} + \eng{ing}) \\ \hline
\eng{begining} & \eng{Beginning} & Double \textbf{n} (Begin + ning) \\ \hline
\eng{speach} & \eng{Speech} & \textbf{ee} son \eng{/i:/} (comme \eng{teach, reach}) pas \eng{ea} \\ \hline
\end{tabular}
\end{center}
\end{exoresolu}

\begin{methode}[Rédiger un dialogue ou un email formel/informel pour une invitation/remmerciement]
\textbf{Objectif :} Produire un écrit cohérent (email, lettre, dialogue) intégrant le vocabulaire du thème, le registre adapté et les conventions de politesse.
\begin{enumerate}
    \item \textbf{Analyser la consigne :} Qui écrit à qui ? (Registre). Quel but ? (Inviter, remercier, demander permission, s'excuser). Format ? (Email, Lettre formelle, SMS, Dialogue).
    \item \textbf{Structurer le message} :
    \begin{itemize}
        \item \textbf{Objet / Sujet} (clair : \eng{Invitation: End-of-Year Ceremony}).
        \item \textbf{Salutation} adaptée (\eng{Dear Sir,} / \eng{Hi John,}).
        \item \textbf{Contexte / Accroche} (\eng{I am writing to invite you...} / \eng{Thanks for the invite...}).
        \item \textbf{Détails clés :} \eng{What, When, Where, Dress code, RSVP}.
        \item \textbf{Formule de politesse de clôture} (\eng{Yours faithfully,} / \eng{Best regards,} / \eng{Cheers,}).
        \item \textbf{Signature}.
    \end{itemize}
    \item \textbf{Incorporer le lexique cible :} 3 collocations minimum, 1 idiome si informel, 0 faux-ami.
    \item \textbf{Relire (Check-list) :} Orthographe (dates, noms propres), Grammaire (modaux, prépositions), Registre (pas de contractions en formel), Ponctuation.
\end{enumerate}
\textbf{Modèle Email Formel (Authorities) :} \eng{Subject: Invitation – [Event Name] / Dear [Title] [Name], / On behalf of [Group], I have the honour to invite you to... / The event will take place on [Date] at [Time] in [Venue]. / We would be grateful if you could [action: deliver a speech / grace the occasion]. / Kindly RSVP by [Date]. / Yours faithfully, / [Name] [Function]}.
\end{methode}

\begin{exoresolu}[Rédaction d'un email formel d'invitation (Autorités)]
\textbf{Énoncé :} Vous êtes le Président du Bureau des Élèves (Student Council President) du Lycée Bilingue de Yaoundé. Rédigez un \textbf{email formel} au \textbf{Proviseur (Principal)} pour l'inviter à la \textbf{Cérémonie de Remise des Prix (Prize-Giving Ceremony)}.
\textbf{Contraintes :} Registre formel soutenu. Inclure : date (vendredi 14 juin), heure (10h), lieu (amphithéâtre), demande de \eng{keynote address} (allocution d'ouverture), \eng{RSVP} avant le 7 juin. Vocabulaire imposé : \eng{host, grant, honour, distinguished, attendance}.

\tcblower
\textbf{Solution détaillée :}

\begin{textebox}[Email Formel -- Modèle Bac]
\textbf{To:} principal.lyceebilingue.yaounde@edu.cm \\
\textbf{Cc:} student.council@edu.cm \\
\textbf{Subject:} Invitation – Annual Prize-Giving Ceremony (June 14)

\textbf{Dear Principal,}

On behalf of the Student Council and the Organising Committee, \textbf{I have the honour to invite you} to \textbf{host} the Annual Prize-Giving Ceremony of the Lycée Bilingue de Yaoundé.

The ceremony will \textbf{take place} on \textbf{Friday, June 14}, at \textbf{10:00 am} in the \textbf{school auditorium}. Your \textbf{distinguished attendance} would greatly \textbf{honour} the laureates and the entire school community.

We would be \textbf{grateful if you could grant} us the privilege of delivering the \textbf{keynote address} on the theme ``Excellence and Discipline''.

Kindly \textbf{RSVP} (confirm your attendance) by \textbf{Friday, June 7}, to allow us to finalise the protocol arrangements.

\textbf{Yours faithfully,}

\textbf{Ayuk Manfred} \\
President, Student Council \\
Lycée Bilingue de Yaoundé \\
+237 6XX XX XX XX
\end{textebox}

\textbf{Analyse des points forts (Grille d'évaluation) :}
\begin{itemize}
    \item \textbf{Registre :} \eng{I have the honour to, would be grateful if you could, kindly, Yours faithfully}. Pas de contractions.
    \item \textbf{Vocabulaire imposé intégré :} \eng{host (v), grant (v), honour (v/n), distinguished (adj), attendance (n)}.
    \item \textbf{Collocations :} \eng{take place, keynote address, finalise arrangements, on behalf of}.
    \item \textbf{Conventions email :} Objet clair, Cc pertinent, Salutation/Clôture formelles (\eng{Dear Principal / Yours faithfully}), Signature complète avec fonction.
    \item \textbf{Grammaire :} Conditionnel passé \eng{would be grateful if you could} (politesse), Passif \eng{be granted} évité au profit de l'actif \eng{grant us}.
\end{itemize}
\end{exoresolu}

\begin{attention}[Les 5 erreurs qui coûtent des points au Bac (Vocabulaire Fêtes \& Interaction)]
\begin{enumerate}
    \item \textbf{\eng{``Make a party / Make a ceremony''} (Calque \textit{Faire une fête}).}
    \eng{$\rightarrow$ CORRECTION :} \eng{Host / Throw (informel) / Organise a party.} \eng{Host / Hold / Conduct a ceremony.}
    \item \textbf{\eng{``Assist the ceremony / Assist to the festival''} (Faux-ami \textit{Assister}).}
    \eng{$\rightarrow$ CORRECTION :} \eng{Attend the ceremony / the festival.} (\eng{Assist} = aider quelqu'un : \eng{assist the elderly}).
    \item \textbf{\eng{``Actually''} pour dire \textit{Actuellement}.}
    \eng{$\rightarrow$ CORRECTION :} \eng{Currently / At present / Right now.} (\eng{Actually} = « En fait / En réalité »).
    \item \textbf{\eng{``Sympathetic''} pour dire \textit{Sympa / Agréable}.}
    \eng{$\rightarrow$ CORRECTION :} \eng{Nice / Pleasant / Enjoyable / Lovely / Friendly.} (\eng{Sympathetic} = compatissant, bienveillant).
    \item \textbf{\eng{``Costume''} pour dire \textit{Costume (homme) / Tenue traditionnelle}.}
    \eng{$\rightarrow$ CORRECTION :} \eng{Suit} (costume ville), \eng{Traditional attire / outfit / wear} (tenue trad.), \eng{Fancy dress / Costume} (déguisement seulement).
\end{enumerate}
\textbf{Astuce de relecture :} Surlignez dans votre copie tous les mots qui ressemblent au français (\eng{actually, assist, manifestation, animation, costume, sympathetic, festival, authority}). Vérifiez systématiquement leur sens anglais dans le contexte.
\end{attention}

\newpage

\coursec{Exercices}

\courssub{Je vérifie mes connaissances}

\begin{exercice}[QCM : Vocabulaire des célébrations]
Chaque phrase ne comporte qu'une seule réponse correcte. Choisis la bonne option (a, b ou c).
\begin{enumerate}
\item The school \eng{\_\_\_\_\_\_\_\_\_\_} a cultural festival every year to promote local traditions.
      \begin{itemize}
      \item[a)] organises
      \item[b)] assists
      \item[c)] attends
      \end{itemize}
\item The \eng{\_\_\_\_\_\_\_\_\_\_} of honour at the ceremony was the Minister of Youth Affairs.
      \begin{itemize}
      \item[a)] guest
      \item[b)] host
      \item[c)] visitor
      \end{itemize}
\item We need to \eng{\_\_\_\_\_\_\_\_\_\_} permission from the principal before putting up posters.
      \begin{itemize}
      \item[a)] demand
      \item[b)] request
      \item[c)] order
      \end{itemize}
\item The \eng{\_\_\_\_\_\_\_\_\_\_} lasted three days and included traditional dances and food fairs.
      \begin{itemize}
      \item[a)] celebration
      \item[b)] party
      \item[c)] festivities
      \end{itemize}
\item My neighbours \eng{\_\_\_\_\_\_\_\_\_\_} us to their house-warming party next Saturday.
      \begin{itemize}
      \item[a)] invited
      \item[b)] inviting
      \item[c)] invitation
      \end{itemize}
\end{enumerate}
\end{exercice}

\begin{exercice}[Vrai ou Faux justifié]
Détermine si les affirmations suivantes sont vraies (V) ou fausses (F) en te basant sur le lexique de la leçon. Justifie ta réponse en français par une phrase complète.
\begin{enumerate}
\item \eng{« I assisted the wedding ceremony yesterday. »} est la phrase correcte pour dire « J'ai assisté au mariage hier. »
\item \eng{« The more the merrier »} signifie « Plus on est de fous, plus on rit ».
\item Pour inviter le proviseur, on écrira : \eng{« Hey Principal, come to our party! »}.
\item \eng{« To break the ice »} veut dire « Briser la glace » au sens propre (faire un trou dans un lac gelé).
\item \eng{« Festivity »} est un synonyme de \eng{« celebration »} et s'utilise souvent au pluriel \eng{« festivities »}.
\end{enumerate}
\end{exercice}

\begin{exercice}[Vocabulaire : Registres de langue]
Classe les douze formules suivantes dans le tableau ci-dessous selon l'interlocuteur le plus approprié (Mate / Neighbour / Adult / School Authority). Certaines formules peuvent convenir à plusieurs catégories.
\begin{center}
\begin{tabular}{|l|}
\hline
\eng{Hey mate! / Good morning Sir/Madam / Hi there! / Excuse me, Mr. ... / What's up? / Dear Principal,} \\
\eng{Hello Mrs. ... / Yo! / Dear Mr. Mayor, / Hi neighbour! / Good afternoon everyone / Respectfully yours} \\
\hline
\end{tabular}
\end{center}
\begin{center}
\begin{tabularx}{\linewidth}{|X|X|X|X|}
\hline
\rowcolor{popblueL} \textbf{Mate (Camarade)} & \textbf{Neighbour (Voisin)} & \textbf{Adult (Adulte)} & \textbf{School Authority (Autorité scolaire)} \\
\hline
 &  &  &  \\
\hline
 &  &  &  \\
\hline
 &  &  &  \\
\hline
\end{tabularx}
\end{center}
\end{exercice}

\begin{exercice}[Conjugaison et formation de mots]
Complète les phrases en conjuguant le verbe entre parenthèses au temps qui convient ou en formant le mot approprié (nom, adjectif, participe).
\begin{enumerate}
\item The committee \eng{\_\_\_\_\_\_\_\_\_\_} (to organise) the end-of-year party since January.
\item We received a formal \eng{\_\_\_\_\_\_\_\_\_\_} (to invite) from the village chief.
\item The \eng{\_\_\_\_\_\_\_\_\_\_} (to celebrate) of the 20th May is a national holiday in Cameroon.
\item Look! The dancers \eng{\_\_\_\_\_\_\_\_\_\_} (to perform) on the stage right now.
\item If we \eng{\_\_\_\_\_\_\_\_\_\_} (to have) enough funds, we would invite a famous artist.
\end{enumerate}
\end{exercice}

\courssub{Je m'entraîne}

\begin{exercice}[Traduction : Du français vers l'anglais]
Traduis les phrases suivantes en anglais correct et naturel. Attention aux faux amis et aux collocations.
\begin{enumerate}
\item Nous avons assisté à une magnifique cérémonie traditionnelle à Bafoussam.
\item Le maire a donné son accord pour la fête de quartier.
\item Ils ont organisé une collecte de fonds pour l'école.
\item « Plus on est de fous, plus on rit ! » a dit mon ami.
\item Nous nous sommes bien amusés lors du festival de Ngondo.
\end{enumerate}
\end{exercice}

\begin{exercice}[Transformation de phrases : Éviter les calques]
Réécris les phrases suivantes en corrigeant les erreurs d'anglicisme ou de structure calquée sur le français. La phrase corrigée doit garder le même sens.
\begin{enumerate}
\item \eng{We made a party for the national day.}
\item \eng{I assisted the meeting with the principal.}
\item \eng{We demanded a permission to the director.}
\item \eng{The animation was very funny.}
\item \eng{He invited me to assist his birthday.}
\end{enumerate}
\end{exercice}

\begin{exercice}[Texte à trous : Festival camerounais]
Complète le texte suivant avec les mots de la banque ci-dessous. Il y a deux mots en trop.
\begin{textebox}[Bank of words]
celebration -- guests of honour -- attend -- festivities -- organise --
invitation -- breaking the ice -- requested -- blast -- authorities
\end{textebox}
\begin{textebox}[Gap-fill text]
Every year, our school \eng{\_\_\_\_\_\_\_\_\_\_} a big cultural \eng{\_\_\_\_\_\_\_\_\_\_} to mark the end of the academic year. The \eng{\_\_\_\_\_\_\_\_\_\_} usually last three days. This year, the \eng{\_\_\_\_\_\_\_\_\_\_} were the Regional Delegate of Education and a famous local musician. We \eng{\_\_\_\_\_\_\_\_\_\_} them formally two months ago. Many parents and neighbours came to \eng{\_\_\_\_\_\_\_\_\_\_} the event. At first, the students were shy, but the traditional dances helped with \eng{\_\_\_\_\_\_\_\_\_\_}. Everyone had a \eng{\_\_\_\_\_\_\_\_\_\_}. The local \eng{\_\_\_\_\_\_\_\_\_\_} praised the students for their discipline.
\end{textebox}
\end{exercice}

\begin{exercice}[Appariement : Expressions idiomatiques]
Associe chaque expression idiomatique anglais (1--8) à son équivalent français (A--H).
\begin{center}
\begin{tabularx}{\linewidth}{|X|X|}
\hline
\rowcolor{popblueL} \textbf{English} & \textbf{Français} \\
\hline
1. To have a blast & A. Mettre l'ambiance / Briser la gêne \\
2. To break the ice & B. S'amuser comme un fou \\
3. The more the merrier & C. Être l'âme de la fête \\
4. To be the life of the party & D. Fêter quelque chose en beauté \\
5. To paint the town red & E. Plus on est de fous, plus on rit \\
6. To let one's hair down & F. Se détendre / Se lâcher \\
7. A social butterfly & G. Quelqu'un de très sociable \\
8. To crash a party & H. S'inviter sans être invité \\
\hline
\end{tabularx}
\end{center}
\end{exercice}

\courssub{Je me prépare au Bac}

\begin{exercice}[Compréhension et langue : Texte original]
Lis attentivement le texte ci-dessous, puis réponds aux questions qui suivent.
\begin{textebox}[Buea Cultural Festival: A Bridge Between Generations]
Last weekend, the town of Buea hosted its annual Cultural Festival at the Molyko Stadium. The event, organised by the Students' Union of the University of Buea in collaboration with the Municipal Council, attracted hundreds of residents and visitors. The Guest of Honour, the Mayor of Buea, arrived at 10 a.m. accompanied by traditional chiefs and school authorities.

In his opening speech, the Mayor emphasised the importance of preserving the Bakweri heritage. "Our festivals are not just entertainment," he stated. "They are a platform for interacting with our elders and passing down values to the youth." He requested the students to respect the traditions displayed during the festivities.

The programme featured traditional wrestling, choral performances, and a food fair exhibiting local dishes like \eng{eru} and \eng{fufu}. Students from various faculties set up booths showcasing handicrafts. The atmosphere was electric; everyone had a blast. For many freshers, it was a perfect opportunity to break the ice and make new friends.

The Vice-Chancellor, Professor Ngomo, closed the ceremony by thanking the organising committee. He invited the authorities to attend next year's edition, promising an even bigger celebration.
\end{textebox}

\textbf{Questions de compréhension (réponds en anglais, phrases complètes) :}
\begin{enumerate}
\item Who organised the Buea Cultural Festival mentioned in the text?
\item What was the Mayor's main message in his opening speech?
\item Name two traditional activities that took place during the festival.
\item Why was the festival a good opportunity for "freshers" (new students)?
\item Who closed the ceremony and what did he promise?
\end{enumerate}

\textbf{Questions de langue :}
\begin{enumerate}
\setcounter{enumi}{5}
\item Find in the text a synonym for "party" or "celebration" (paragraph 1).
\item Find in the text the phrase that means "transmitting values to the younger generation" (paragraph 2).
\item \eng{Grammar:} Rewrite the sentence \eng{« He requested the students to respect the traditions »} using the structure \eng{« He asked... »}.
\item \eng{Word formation:} Complete the sentence with the right form of the word in brackets: The \eng{\_\_\_\_\_\_\_\_\_\_} (organise) committee did a great job.
\item \eng{Vocabulary:} Explain the meaning of \eng{« to have a blast »} and \eng{« to break the ice »} in the context of the text.
\end{enumerate}
\end{exercice}

\begin{exercice}[Traduction : Thème]
Traduis le paragraphe suivant en anglais. Soigne le registre de langue (formel) et le vocabulaire spécifique aux cérémonies officielles.
\begin{quote}
« Monsieur le Maire, au nom du comité d'organisation, je vous adresse une invitation formelle pour la fête culturelle de notre établissement qui aura lieu le 20 mai prochain. Votre présence honorerait notre cérémonie et encouragerait les élèves. Nous avons demandé l'autorisation aux autorités scolaires et tout est prêt pour accueillir les invités d'honneur. Nous espérons que vous pourrez assister à l'événement. »
\end{quote}
\end{exercice}

\begin{exercice}[Expression écrite : Sujet de composition]
Traite l'un des deux sujets suivants. Ton texte comportera entre \textbf{120 et 150 mots}. Il sera évalué sur la correction de la langue, la richesse du vocabulaire (thème des festivals, registres de politesse) et la cohérence.

\textbf{Sujet A :} You are the President of the Students' Union in your school. Write a formal letter to the Principal (or the Mayor of your city) inviting him/her to the upcoming School Cultural Festival. Mention the date, the programme highlights, the guests of honour expected, and why his/her presence is important.

\textbf{Sujet B :} "Community festivals are essential for social cohesion." Write an article for your school magazine discussing this statement. Give examples from your own experience (Ngondo, Medumba, Nyem-Nyem, school festivals, etc.) and explain how these events help people interact across generations and social backgrounds.
\end{exercice}

\newpage

\coursec{Corrigés des exercices}

\begin{corrige}[Exercice 1 — QCM : Vocabulaire des célébrations]
\begin{enumerate}
\item \textbf{a) organises.} \eng{“The school organises a cultural festival every year to promote local traditions.”} « L'école organise un festival culturel chaque année. » \eng{assists} est un faux ami : il signifie « aider », pas « assister à » ; \eng{attends} signifie « assister à » mais ne convient pas ici (l'école est l'organisatrice, pas la spectatrice).
\item \textbf{a) guest.} \eng{“The guest of honour at the ceremony was the Minister of Youth Affairs.”} « L'invité d'honneur de la cérémonie était le ministre de la Jeunesse. » La collocation figée est \eng{guest of honour}. \eng{host} = l'hôte (celui qui reçoit) ; \eng{visitor} = un simple visiteur.
\item \textbf{b) request.} \eng{“We need to request permission from the principal before putting up posters.”} « Nous devons demander l'autorisation au proviseur. » \eng{request} est le verbe formel approprié ; \eng{demand} est trop agressif (exiger) ; \eng{order} signifie « commander ».
\item \textbf{c) festivities.} \eng{“The festivities lasted three days and included traditional dances and food fairs.”} « Les festivités ont duré trois jours. » Le sujet est au pluriel (verbe \eng{lasted} sans -s), donc \eng{celebration} et \eng{party} (singuliers) sont impossibles.
\item \textbf{a) invited.} \eng{“My neighbours invited us to their house-warming party next Saturday.”} « Mes voisins nous ont invités à leur pendaison de crémaillère samedi prochain. » Il faut un verbe conjugué au prétérit ; \eng{inviting} est un participe et \eng{invitation} un nom.
\end{enumerate}
\end{corrige}

\begin{corrige}[Exercice 2 — Vrai ou Faux justifié]
\begin{enumerate}
\item \textbf{Faux.} \eng{to assist} est un faux ami : il signifie « aider ». Pour dire « assister à », il faut employer \eng{to attend} : \eng{“I attended the wedding ceremony yesterday.”}
\item \textbf{Vrai.} L'expression idiomatique \eng{The more the merrier} correspond exactement au proverbe français « Plus on est de fous, plus on rit ».
\item \textbf{Faux.} S'adresser au proviseur exige un registre formel ; on écrira plutôt \eng{“Dear Principal, we would be honoured to invite you to our party.”} et non une formule familière comme \eng{Hey}.
\item \textbf{Faux.} \eng{To break the ice} est une expression idiomatique qui signifie « briser la glace » au sens figuré, c'est-à-dire mettre les gens à l'aise, et non faire un trou dans un lac gelé.
\item \textbf{Vrai.} \eng{festivity} est bien un synonyme de \eng{celebration} et s'emploie très souvent au pluriel \eng{festivities} pour désigner l'ensemble des réjouissances.
\end{enumerate}
\end{corrige}

\begin{corrige}[Exercice 3 — Vocabulaire : Registres de langue]
\begin{center}
\begin{tabularx}{\linewidth}{|X|X|X|X|}
\hline
\rowcolor{popblueL} \textbf{Mate (Camarade)} & \textbf{Neighbour (Voisin)} & \textbf{Adult (Adulte)} & \textbf{School Authority (Autorité scolaire)} \\
\hline
\eng{Hey mate!} & \eng{Hi neighbour!} & \eng{Good morning Sir/Madam} & \eng{Good morning Sir/Madam} \\
\hline
\eng{What's up?} & \eng{Hi there!} & \eng{Excuse me, Mr. ...} & \eng{Dear Principal,} \\
\hline
\eng{Yo!} & \eng{Hello Mrs. ...} & \eng{Good afternoon everyone} & \eng{Respectfully yours} \\
\hline
\end{tabularx}
\end{center}
\textbf{Justifications :} \eng{Hey mate!}, \eng{What's up?} et \eng{Yo!} relèvent du registre familier, réservé aux camarades. \eng{Hi neighbour!} et \eng{Hi there!} sont amicaux, adaptés aux voisins. \eng{Good morning Sir/Madam}, \eng{Excuse me, Mr. ...} et \eng{Good afternoon everyone} sont polis et conviennent aux adultes. \eng{Dear Principal,} et la formule de clôture \eng{Respectfully yours} appartiennent au registre formel des courriers adressés aux autorités scolaires. \eng{Dear Mr. Mayor,} (autorité administrative) se rapproche de la colonne « School Authority » par son degré de formalité.
\end{corrige}

\begin{corrige}[Exercice 4 — Conjugaison et formation de mots]
\begin{enumerate}
\item \eng{“The committee has been organising the end-of-year party since January.”} Present perfect continu : action commencée en janvier (\eng{since}) et qui continue. On accepte aussi \eng{has organised}.
\item \eng{“We received a formal invitation from the village chief.”} Nom formé sur le verbe \eng{to invite} ; attention à l'orthographe : \eng{invitation}, pas \eng{invitacion}.
\item \eng{“The celebration of the 20th May is a national holiday in Cameroon.”} Nom dérivé de \eng{to celebrate}. « La célébration du 20 mai est un jour férié national au Cameroun. »
\item \eng{“Look! The dancers are performing on the stage right now.”} Present continu (\eng{Look!} + \eng{right now}) : action en cours au moment où l'on parle.
\item \eng{“If we had enough funds, we would invite a famous artist.”} Conditionnel de type 2 : \eng{If + prétérit, ... would + base verbale} pour une hypothèse irréelle au présent.
\end{enumerate}
\end{corrige}

\begin{corrige}[Exercice 5 — Traduction : Du français vers l'anglais]
\begin{enumerate}
\item \eng{“We attended a wonderful traditional ceremony in Bafoussam.”} Attention au faux ami : « assister à » se traduit par \eng{to attend}, jamais \eng{to assist}.
\item \eng{“The mayor gave his approval for the neighbourhood party.”} Collocation : \eng{to give one's approval} ; « quartier » = \eng{neighbourhood}.
\item \eng{“They organised a fund-raising event for the school.”} « Collecte de fonds » = \eng{fund-raising (event)} ; on ne dit pas \eng{a collection of funds}.
\item \eng{“ ‘The more the merrier!’ said my friend.”} Expression idiomatique figée à apprendre par cœur.
\item \eng{“We had a lot of fun at the Ngondo festival.”} « S'amuser » = \eng{to have fun} ; ne pas traduire par \eng{we amused ourselves well}. On accepte aussi \eng{“We had a blast at the Ngondo festival.”}
\end{enumerate}
\end{corrige}

\begin{corrige}[Exercice 6 — Transformation de phrases : Éviter les calques]
\begin{enumerate}
\item \eng{“We held / organised a party for the national day.”} En anglais, on ne « fait » pas une fête : collocation \eng{to hold / organise / throw a party}.
\item \eng{“I attended the meeting with the principal.”} Faux ami : \eng{to assist} = aider ; « assister à » = \eng{to attend}.
\item \eng{“We requested permission from the director.”} \eng{to demand} est trop fort, et \eng{permission} est indénombrable (pas de \eng{a}) ; la préposition est \eng{from}, pas \eng{to}.
\item \eng{“The entertainment was very funny.”} Faux ami : \eng{animation} en anglais désigne les dessins animés ; « l'animation » d'une fête = \eng{entertainment}.
\item \eng{“He invited me to attend / come to his birthday party.”} Deux corrections : \eng{attend} à la place de \eng{assist}, et la collocation \eng{birthday party} (on n'assiste pas à un « anniversaire » mais à une fête d'anniversaire).
\end{enumerate}
\end{corrige}

\begin{corrige}[Exercice 7 — Texte à trous : Festival camerounais]
\textbf{Texte complété :}
\begin{quote}
\eng{Every year, our school organises a big cultural celebration to mark the end of the academic year. The festivities usually last three days. This year, the guests of honour were the Regional Delegate of Education and a famous local musician. We requested them formally two months ago. Many parents and neighbours came to attend the event. At first, the students were shy, but the traditional dances helped with breaking the ice. Everyone had a blast. The local authorities praised the students for their discipline.}
\end{quote}

\textbf{Mots en trop (banque de 10--11 items pour 9 trous) :} Les formes utilisées dans le texte sont : \eng{organises} (verbe conjugué), \eng{celebration}, \eng{festivities}, \eng{guests of honour}, \eng{requested}, \eng{attend}, \eng{breaking the ice}, \eng{blast}, \eng{authorities}. Le mot \eng{invitation} (nom) n'est pas utilisé. L'autre intrus est la forme de base d'un mot employé sous forme fléchie (ex. \eng{organise} au lieu de \eng{organises}, \eng{request} au lieu de \eng{requested}, \eng{authority} au singulier, \eng{guest of honour} au singulier, ou \eng{break the ice} à l'infinitif).

\textbf{Remarques grammaticales :} \eng{organises} prend un -s (sujet singulier \eng{our school}, présent simple d'habitude) ; \eng{festivities} est au pluriel car le verbe \eng{last} est sans -s ; \eng{requested} est au prétérit (\eng{two months ago}).
\end{corrige}

\begin{corrige}[Exercice 8 — Appariement : Expressions idiomatiques]
\begin{center}
\begin{tabularx}{\linewidth}{|X|X|}
\hline
\rowcolor{popblueL} \textbf{English} & \textbf{Français} \\
\hline
1. To have a blast & B. S'amuser comme un fou \\
2. To break the ice & A. Mettre l'ambiance / Briser la gêne \\
3. The more the merrier & E. Plus on est de fous, plus on rit \\
4. To be the life of the party & C. Être l'âme de la fête \\
5. To paint the town red & D. Fêter quelque chose en beauté \\
6. To let one's hair down & F. Se détendre / Se lâcher \\
7. A social butterfly & G. Quelqu'un de très sociable \\
8. To crash a party & H. S'inviter sans être invité \\
\hline
\end{tabularx}
\end{center}
\textbf{Astuce méthodologique :} ces expressions sont figées ; il faut les mémoriser en bloc, sans chercher à les traduire mot à mot (\eng{to paint the town red} ne veut pas dire « peindre la ville en rouge » !).
\end{corrige}

\begin{corrige}[Exercice 9 — Compréhension et langue : Texte original]
\textbf{Barème indicatif : 20 points} — Compréhension : 10 pts (2 pts par question) ; Langue : 10 pts (2 pts par question).

\textbf{Questions de compréhension (réponses modèles) :}
\begin{enumerate}
\item \eng{“The festival was organised by the Students' Union of the University of Buea in collaboration with the Municipal Council.”}
\item \eng{“The Mayor's main message was that festivals are not just entertainment but a platform for interacting with elders and passing down values to the youth.”}
\item \eng{“Two traditional activities were traditional wrestling and choral performances.”} (On accepte aussi la foire gastronomique avec \eng{eru} et \eng{fufu}.)
\item \eng{“It was a good opportunity for freshers to break the ice and make new friends.”}
\item \eng{“The Vice-Chancellor, Professor Ngomo, closed the ceremony and promised an even bigger celebration next year.”}
\end{enumerate}

\textbf{Questions de langue :}
\begin{enumerate}
\setcounter{enumi}{5}
\item \eng{Festival} (ou \eng{event}) au paragraphe 1.
\item \eng{“passing down values to the youth”} (paragraphe 2).
\item \eng{“He asked the students to respect the traditions.”} La structure \eng{ask somebody to do something} remplace directement \eng{request somebody to do something}, avec un registre légèrement moins formel.
\item \eng{“The organising committee did a great job.”} Adjectif verbal en -ing placé devant le nom ; on accepte aussi \eng{organisation committee}, mais \eng{organising} est la forme attendue.
\item \eng{“To have a blast”} means to have a lot of fun, to enjoy oneself enormously — in the text, the atmosphere was electric and everyone enjoyed the festival. \eng{“To break the ice”} means to overcome initial shyness and start socialising — in the text, the festival helped new students make friends.
\end{enumerate}

\textbf{Points de vigilance :} répondre par des phrases complètes avec sujet et verbe ; ne pas recopier le texte sans le reformuler ; accorder les verbes au passé (\eng{was organised}, \eng{closed}, \eng{promised}).
\end{corrige}

\begin{corrige}[Exercice 10 — Traduction : Thème]
\textbf{Barème indicatif : 10 points} — 2 pts par segment correctement traduit (registre formel, vocabulaire des cérémonies, grammaire).

\textbf{Traduction modèle :}
\begin{quote}
\eng{“Dear Mr. Mayor, on behalf of the organising committee, I would like to extend a formal invitation to you for our school's cultural festival, which will take place on the 20th of May. Your presence would honour our ceremony and encourage the students. We have requested permission from the school authorities and everything is ready to welcome the guests of honour. We hope that you will be able to attend the event.”}
\end{quote}

\textbf{Points de vigilance :}
\begin{itemize}
\item « Au nom de » = \eng{on behalf of}, jamais \eng{in the name of} dans ce contexte de courrier.
\item « Adresser une invitation » = \eng{to extend / send an invitation} (collocation formelle).
\item « Demander l'autorisation » = \eng{to request permission} ; \eng{permission} est indénombrable.
\item « Invités d'honneur » = \eng{guests of honour} (collocation figée).
\item « Assister à » = \eng{to attend} ; bannir le faux ami \eng{to assist}.
\item « Établissement » (scolaire) = \eng{school}, pas \eng{establishment}.
\end{itemize}
\end{corrige}

\begin{corrige}[Exercice 11 — Expression écrite : Sujet de composition]
\textbf{Barème indicatif : 20 points} — Respect de la tâche et du format : 4 pts ; richesse du vocabulaire thématique (festivals, registres) : 6 pts ; correction grammaticale et orthographique : 6 pts ; cohérence et organisation : 4 pts. Longueur exigée : 120 à 150 mots.

\textbf{Proposition de rédaction modèle — Sujet A (lettre formelle) :}
\begin{quote}
\eng{Dear Principal,\\
On behalf of the Students' Union, I am honoured to invite you to our annual School Cultural Festival, which will take place on Friday, 15th March, in the school courtyard.\\
The programme will feature traditional dances, a choral competition, a drama performance and a food fair organised by the parents' association. The Regional Delegate of Secondary Education and the Mayor of our town have confirmed their attendance as guests of honour.\\
Your presence would greatly encourage the students and honour the organising committee. It would also be a wonderful opportunity for you to interact with parents and neighbours of the school community.\\
We sincerely hope you will be able to attend this celebration.\\
Respectfully yours,\\
Amina Mbarga, President of the Students' Union}
\end{quote}
(Environ 125 mots.)

\textbf{Proposition de rédaction modèle — Sujet B (article) :}
\begin{quote}
\eng{Community festivals are essential for social cohesion because they bring together people of all ages and backgrounds. During the Ngondo festival in Douala, for example, neighbours, elders and young people gather on the banks of the Wouri River to celebrate their shared heritage. Traditional dances, canoe races and food fairs create a joyful atmosphere where everyone can break the ice and interact freely.\\
In our schools too, cultural festivals allow students to meet school authorities in a relaxed setting and to learn values from their elders. These events are a platform for passing down traditions to the youth.\\
In conclusion, festivals are not just entertainment: they strengthen the bonds between generations and remind us that, as the saying goes, the more the merrier!}
\end{quote}
(Environ 125 mots.)

\textbf{Points de vigilance :}
\begin{itemize}
\item Sujet A : respecter la mise en page de la lettre formelle (formule d'appel \eng{Dear Principal,}, clôture \eng{Respectfully yours}), employer \eng{on behalf of}, \eng{guests of honour}, \eng{to attend}.
\item Sujet B : structurer l'article (introduction, exemples, conclusion) et mobiliser au moins trois expressions idiomatiques de la leçon.
\item Éviter les calques : \eng{to hold a party}, \eng{to have fun}, \eng{to attend} (et non \eng{to assist}).
\item Compter les mots : un texte hors limites (moins de 120 ou plus de 150) est pénalisé.
\end{itemize}
\end{corrige}

\newpage

\coursec{Fiche bilan}

\begin{popfigure}
\begin{tikzpicture}[
    mindmap,
    concept color=popblue,
    level 1 concept/.append style={level distance=3.5cm, sibling angle=60},
    level 2 concept/.append style={level distance=2.5cm},
    every node/.style={concept, execute at begin node=\small\bfseries, align=center},
    mate/.style={concept color=poporange},
    neighbour/.style={concept color=popgreen},
    adult/.style={concept color=poppurple},
    authority/.style={concept color=poppink},
    vocab/.style={concept color=popgold},
    grammar/.style={concept color=popblue!70!poppurple},
    trap/.style={concept color=popred},
]
\node[concept, align=center]{Community Celebrations \\ \& Festivals}
    child[mate] { node[align=center]{Interacting with\\ Mates} 
        child { node {Slang \& Informal} }
        child { node {Invitations} }
    }
    child[neighbour] { node[align=center]{Interacting with\\ Neighbours} 
        child { node {Polite \& Friendly} }
        child { node {Community Prep} }
    }
    child[adult] { node[align=center]{Interacting with\\ Adults} 
        child { node {Respectful Tone} }
        child { node {Advice \& Blessings} }
    }
    child[authority] { node[align=center]{Interacting with\\ Authorities} 
        child { node {Formal Register} }
        child { node {Permissions \& Reports} }
    }
    child[vocab] { node {Core Vocabulary} 
        child { node {Events \& Items} }
        child { node {People \& Roles} }
        child { node {Actions \& Verbs} }
    }
    child[grammar] { node {Grammar Tools} 
        child { node {Modals for Politeness} }
        child { node {Indirect Questions} }
    }
    child[trap] { node {Traps \& Tips} 
        child { node {False Friends} }
        child { node {Pronunciation} }
    };
\end{tikzpicture}
\legende{Carte mentale : Vocabulaire des célébrations communautaires et registres de langue}
\end{popfigure}

\begin{bilan}

\end{bilan}

\courssub{Lexique clé — English / Français}

\begin{center}
\begin{tabularx}{\linewidth}{|X|X|X|}
\hline
\rowcolor{popblueL} \textbf{Category / Catégorie} & \textbf{English} & \textbf{Français} \\
\hline
\cle{Events} & \eng{festival, celebration, ceremony, anniversary, feast, holiday, national day, cultural week} & festival, célébration, cérémonie, anniversaire, fête, jour férié, fête nationale, semaine culturelle \\
\hline
\cle{People} & \eng{mate, peer, neighbour, elder, adult, authority, principal, head teacher, prefect, organiser, participant, guest, host} & camarade, pair, voisin, aîné, adulte, autorité, proviseur, chef d'établissement, préfet, organisateur, participant, invité, hôte \\
\hline
\cle{Actions} & \eng{celebrate, organise, attend, participate, invite, welcome, greet, congratulate, thank, apologise, request permission, deliver a speech, perform} & célébrer, organiser, assister à, participer, inviter, accueillir, saluer, féliciter, remercier, s'excuser, demander la permission, prononcer un discours, performer \\
\hline
\cle{Items/Prep} & \eng{invitation card, programme, venue, date, time, dress code, traditional attire, decoration, refreshments, budget, committee, announcement} & carte d'invitation, programme, lieu, date, heure, code vestimentaire, tenue traditionnelle, décoration, rafraîchissements, budget, comité, annonce \\
\hline
\cle{Atmosphere} & \eng{joyful, lively, colourful, memorable, unity, solidarity, heritage, tradition, harmony, excitement} & joyeux, animé, coloré, mémorable, unité, solidarité, héritage, tradition, harmonie, excitation \\
\hline
\end{tabularx}
\end{center}

\courssub{Registres de langue et politesse (Register \& Politeness)}

\begin{center}
\begin{tabularx}{\linewidth}{|X|X|X|}
\hline
\rowcolor{poporangeL} \textbf{Interlocuteur} & \textbf{Register / Structures} & \textbf{Exemples} \\
\hline
\cle{Mates / Peers} & Informel, \eng{phrasal verbs}, contractions, slang léger & \eng{“Hey! Wanna join the fest?”} (« Hé ! Tu viens à la fête ? ») \\
\hline
\cle{Neighbours} & Semi-formel, poli, \eng{would like to, could} & \eng{“Good morning. We would like to invite you…”} (« Bonjour. Nous voudrions vous inviter… ») \\
\hline
\cle{Adults / Elders} & Respectueux, titres (\eng{Mr, Mrs, Auntie, Uncle}), \eng{may, might} & \eng{“Good afternoon, Uncle. May we have your blessing?”} (« Bonjour Tonton. Pourrions-nous avoir votre bénédiction ? ») \\
\hline
\cle{School Authorities} & Formel, \eng{passive voice}, modaux \eng{would, could, might}, indirect questions & \eng{“The Principal would be grateful if you could attend.”} (« Le Proviseur serait reconnaissant si vous pouviez assister. ») \\
\hline
\end{tabularx}
\end{center}

\courssub{Collocations et expressions idiomatiques}

\begin{itemize}
    \item \cle{to kick off} (commencer, lancer) — \eng{The festival kicks off at 10 a.m.}
    \item \cle{to be in full swing} (battre son plein) — \eng{By noon, the celebration was in full swing.}
    \item \cle{to lend a hand} (donner un coup de main) — \eng{Everyone lent a hand to decorate the hall.}
    \item \cle{to make a speech / deliver an address} (prononcer un discours)
    \item \cle{to grace the occasion} (honorer de sa présence — formel, autorités)
    \item \cle{a token of appreciation} (un geste de reconnaissance)
    \item \cle{to RSVP} (répondre à une invitation — \eng{Répondez S'il Vous Plaît})
\end{itemize}

\courssub{Faux amis et pièges fréquents (Francophones)}

\begin{center}
\begin{tabularx}{\linewidth}{|X|X|X|}
\hline
\rowcolor{popred!20} \textbf{Mot anglais} & \textbf{Piège (Français)} & \textbf{Sens réel / Correction} \\
\hline
\eng{actually} & \eng{actuellement} & \cle{in fact / en fait} — \eng{currently = at present} \\
\eng{assist} & \eng{assister à} & \cle{help / aider} — \eng{attend = assister à} \\
\eng{ceremony} & \eng{cérémonie (seulement officiel)} & \cle{ceremony} ok, mais \eng{party} pour fête entre amis \\
\eng{disguise} & \eng{déguisement (carnaval)} & \cle{costume / fancy dress} pour fête costumée ; \eng{disguise} = cacher son identité \\
\eng{entertain} & \eng{divertir (seulement)} & \cle{receive / host} pour « recevoir des invités » \\
\eng{festival} & \eng{festival (musique seulement)} & \cle{festival} = événement culturel/série de spectacles ; \eng{celebration} = terme générique \\
\eng{invite} (nom) & \eng{une invite} & \cle{invitation} (nom) ; \eng{invite} = verbe \\
\eng{programme} & \eng{programme (orthographe FR)} & \cle{program} (US) / \eng{programme} (UK) — attention \eng{programmer} n'existe pas (verbe : \eng{to schedule}) \\
\hline
\end{tabularx}
\end{center}

\courssub{Familles de mots, prononciation, orthographe}

\begin{center}
\begin{tabular}{|l|l|l|l|}
\hline
\rowcolor{popgreenL} \textbf{Verbe} & \textbf{Nom} & \textbf{Adjectif} & \textbf{Prononciation (approx.)} \\
\hline
\eng{celebrate} & \eng{celebration} & \eng{celebratory} & « sè-lè-breï-chn » / « sè-lè-breï-to-ri » \\
\eng{organise} & \eng{organisation, organiser} & \eng{organised} & « o-ga-naïz » / « o-ga-naï-zeï-chn » \\
\eng{invite} & \eng{invitation} & \eng{inviting} & « in-vait » / « in-vi-tei-chn » \\
\eng{participate} & \eng{participation, participant} & \eng{participatory} & « par-ti-si-peït » / « par-ti-si-peï-chn » \\
\eng{congratulate} & \eng{congratulations} & \eng{congratulatory} & « kon-grà-tchu-leït » / « kon-grà-tchu-leï-chn » \\
\eng{apologise} & \eng{apology} & \eng{apologetic} & « a-po-lo-djaïz » / « a-po-lo-dji » \\
\hline
\end{tabular}
\end{center}

\emph{Règles d'orthographe :} \eng{organise} (UK) vs \eng{organize} (US) — au Cameroun, norme britannique. Doublement du \eng{l} : \eng{travel $\rightarrow$ travelled}, \eng{cancel $\rightarrow$ cancelled}. \eng{programme} (nom UK), \eng{program} (verbe/informatique).

\courssub{Méthodes express}

\begin{methode}[Rédiger une invitation formelle (Authorities/Adults)]
\begin{enumerate}
    \item \textbf{En-tête :} Sender's address (top right), Date, Recipient's address (left).
    \item \textbf{Objet :} \eng{Re: Invitation to [Event Name]}.
    \item \textbf{Salutation :} \eng{Dear Sir/Madam,} ou \eng{Dear Mr/Mrs [Name],}.
    \item \textbf{Corps :} \eng{We have the honour to invite you...} / \eng{The Principal requests the pleasure of your company...}. Préciser \eng{date, time, venue, dress code, RSVP date}.
    \item \textbf{Formule de politesse :} \eng{Yours faithfully,} (si \eng{Dear Sir/Madam}) ou \eng{Yours sincerely,} (si nom connu).
    \item \textbf{Signature :} Name, Title.
\end{enumerate}
\end{methode}

\begin{methode}[S'excuser / Refuser poliment]
\begin{enumerate}
    \item Remercier : \eng{Thank you for the invitation.}
    \item Exprimer le regret : \eng{I regret to inform you that I cannot attend...} / \eng{I am afraid I will not be able to...}
    \item Donner une raison vague (optionnel) : \eng{due to prior commitments}.
    \item Souhaiter du succès : \eng{I wish you a successful celebration.}
\end{enumerate}
\end{methode}



\begin{aretenir}[Checklist avant le Bac]
\begin{itemize}
    \item $\square$ Je distingue 4 registres : \cle{Mates (informel)}, \cle{Neighbours (semi-formel)}, \cle{Adults (respectueux)}, \cle{Authorities (formel)}.
    \item $\square$ Je connais 15 mots de base du champ lexical \eng{celebrations} (noms, verbes, adjectifs).
    \item $\square$ J'utilise les modaux de politesse : \eng{can/could, will/would, may/might} selon l'interlocuteur.
    \item $\square$ Je forme des questions indirectes pour les autorités : \eng{“Could you tell me if...?”} / \eng{“I would like to know whether...”}.
    \item $\square$ J'évite les 7 faux amis listés (\eng{actually, assist, disguise, entertain, invite (n), programme, ceremony}).
    \item $\square$ J'écris correctement les familles de mots : \eng{organise/organisation, invite/invitation, participate/participation}.
    \item $\square$ Je prononce correctement : \eng{celebration} (accent sur \eng{-bra-}), \eng{invitation} (accent sur \eng{-ta-}), \eng{RSVP} (\eng{ar-es-vi-pi}).
    \item $\square$ Je sais rédiger une invitation formelle (structure + formules \eng{Yours faithfully/sincerely}).
    \item $\square$ Je sais refuser une invitation poliment par écrit et à l'oral.
    \item $\square$ J'utilise 3 collocations idiomatiques : \eng{kick off, in full swing, lend a hand, grace the occasion}.
    \item $\square$ Je contextualise : je cite un exemple camerounais (Fête de la Jeunesse, Ngondo, fête de l'école).
    \item $\square$ Je relis mes productions pour vérifier : registre adapté, pas de calque français, orthographe UK.
\end{itemize}
\end{aretenir}

\findecours

\end{document}
